\chapter{Réalisations spectrales, scalaires et nulles de la loi générale des masses : familles matérielles, transports CKM–PMNS et secteurs photonique et gluonique}
\label{chap:realizaciones-familias-mezclas-masas}

\section{Codomaines des réalisations matérielles : opérateur spectral, scalaire, nullité de jauge, composé et pôle résonant}
Une fibre peut produire un opérateur spectral, un scalaire après compression physique, une nullité exacte de jauge, un opérateur composé ou un pôle résonant doté d’une largeur. Ces sorties possèdent des codomaines distincts et ne sont pas comptées comme classes du domaine \(81/56/13/324\).

\section{Extensions nulles, propagation entre frontières et fermeture électronique}

Les secteurs photonique et gluonique se réalisent par nullité exacte, tandis que la fermeture électronique occupe le centre matériel. Le développement complet distingue propagation nulle, entrelacement aller--retour, interface physique et critères de fermeture sans transformer une nullité en absence de structure.

\begingroup
\let\chapter\section
\let\section\subsection
\let\subsection\subsubsection
\let\subsubsection\paragraph
\let\clearpage\relax
\let\cleardoublepage\relax
\chapter{Extensions nulles et clôture matérielle}
\label{chap:realizaciones-nula-material}
\index{photon!interface physique}\index{électron!clôture matérielle}
\index{extension nulle}\index{cône causal scindé}

\section{\(4\) réalisations de l'état matériel commun : nulle, massive, composée et topologique}
\label{sec:cuatro-modos-realizacion}

Un chiffre isolé détermine une coordonnée, mais n'identifie ni la structure
qui le produit, ni la mémoire qu'il conserve, ni l'opération qui le transforme en
grandeur. HMT se situe avant cette projection : son objet primaire est une
extension compatible qui conserve les lois arithmétiques, l'orientation
TRIT, les quotients, la phase et la mémoire de frontière. La mécanique dimensionnelle
spécifie ensuite quelle composante se réalise comme grandeur, laquelle demeure
comme mémoire et quelle transformation produit l'observable.

Les développements de cette partie distinguent quatre modes par leur opération de
fermeture :
\begin{enumerate}
 \item une extension ouverte et nulle transporte une perturbation entre
 frontières ;
 \item une extension fermée et persistante acquiert une norme intérieure et définit
 le candidat matériel ;
 \item une famille d'extensions pondérée par l'action définit un ensemble
 thermique ;
 \item une connexion qui compare des sections en des points distincts fournit le
 codomaine de la réalisation gravitationnelle.
\end{enumerate}
La source commune n'identifie pas les quatre codomaines. Les applications physiques
correspondantes restent discontinues :
\[
 \mathfrak R_\gamma:
 \Omega_{\HMT}^{\rm surv}\dashrightarrow\mathcal H_{\rm rad},
 \qquad
 \mathfrak R_m:
 \Omega_{\HMT}^{\rm surv}\dashrightarrow\mathcal H_{\rm mat},
\]
\[
 \mathfrak R_\Theta:
 \mathcal P(\Omega_{\HMT}^{\rm surv})\dashrightarrow\mathcal S_{\rm th},
 \qquad
 \mathfrak R_{\rm grav}:
 \Omega_{\HMT}^{\rm surv}\dashrightarrow\mathsf{Cartan}_{1,3}.
\]
Leur unité réside dans la généalogie ; leur différence, dans la lecture et dans les
structures supplémentaires exigées par chaque réalisation.

\section{Propagation nulle entre frontières}
\label{sec:propagacion-nula-fotonica}

Soit \(B\) une base de transport. Une \emph{section nulle} de signe
\(\pm\) est une application
\[
 s_\pm:B\longrightarrow\mathcal A_{\rm sp}P_\pm
\]
telle que \(s_\pm(x)\in\mathcal A_{\rm sp}P_\pm\) pour tout \(x\in B\).
Par la \cref{prop:descomposicion-espectral-partida},
\[
 N_{\rm sp}(s_\pm(x))=0
 \qquad(x\in B).
\]
Si \(\gamma\) est une voie orientée entre deux frontières,
\[
 \partial\gamma=x_{\rm abs}-x_{\rm em},
\]
l'interface photonique proposée attribue à \(\gamma\) une section nulle
étendue. L'adjectif \emph{étendue} indique que son identité appartient au
transport complet et non à un point isolé.

La voie conjuguée s'obtient par inversion,
\[
 \gamma^\vee=C_{\rm rev}\gamma,
 \qquad
 \partial\gamma^\vee=x_{\rm em}-x_{\rm abs}.
\]
Relativement au lecteur d'orientation déclaré, la paire
\((\gamma,\gamma^\vee)\) représente les lectures avancée et retardée. Les
deux branches sont ordonnées séquentiellement ou logées dans des fibres distinctes. Elles ne
s'ajoutent pas automatiquement dans une même fibre : si les deux composantes conjuguées
sont activées simultanément, l'identité exacte
\[
 \boxed{N_{\rm sp}(uP_++vP_-)=uv}
\]
montre que, pour \(u,v>0\), la section quitte la frontière nulle et entre dans
l'intérieur positif.

\subsection{Longueur, aller et retour}

La carte dimensionnelle de vitesse,
\(\widehat c=e^{-\mathcal E}\), et une base positive \(u_L\) de la droite de
longueur étant fixées, on définit
\[
 L_{\rm int}(\gamma)=\widehat c\,|\gamma|\,u_L.
\]
L'inversion conserve le nombre de pas :
\[
 L_{\rm int}(\gamma^\vee)=L_{\rm int}(\gamma),
 \qquad
 L_{\leftrightarrow}=2L_{\rm int}(\gamma).
\]
Ces identités sont cinématiques. Les transformer en propagation
électromagnétique exige qu'une même application reproduise la loi de Gauss, la
dynamique de jauge, la dispersion nulle et la réduction à deux degrés
transversaux.

\subsection{Entropie d'intrication entre les \(2\) lectures}
\label{sec:entrelazamiento-ida-retorno}

L'inversion d'une voie ne suffit pas pour affirmer une intrication quantique. Une
formulation minimale exige d'abord une complétion de Hilbert. Soient
\(\mathcal H_{\rm adv}\) et \(\mathcal H_{\rm ret}\) deux espaces complexes
avec des familles orthonormales
\(\{|i\rangle_{\rm adv}\}_{i\in I}\) et
\(\{|i\rangle_{\rm ret}\}_{i\in I}\), où \(I\) est fini. Pour une
probabilité \(p=(p_i)_{i\in I}\), nous définissons
\[
 |\Psi_p\rangle
 :=\sum_{i\in I}\sqrt{p_i}\,
 |i\rangle_{\rm adv}\otimes|i\rangle_{\rm ret}.
 \tag{E1}\label{eq:estado-entrelazado-ida-retorno}
\]

\begin{proposition}[Entropie bicouche de l'aller et du retour]
\label{prop:entropia-entrelazamiento-ida-retorno}
Les matrices réduites de l'état
\(\rho_p=|\Psi_p\rangle\langle\Psi_p|\) sont
\[
 \rho_{\rm adv}
 =\sum_{i\in I}p_i|i\rangle_{\rm adv}\langle i|,
 \qquad
 \rho_{\rm ret}
 =\sum_{i\in I}p_i|i\rangle_{\rm ret}\langle i|,
 \tag{E2}\label{eq:reducidas-ida-retorno}
\]
et possèdent la même entropie de von Neumann :
\[
 \boxed{
 S_{\rm ent}(p)
 =-\operatorname{Tr}(\rho_{\rm adv}\log\rho_{\rm adv})
 =-\sum_{i\in I}p_i\log p_i.}
 \tag{E3}\label{eq:entropia-ida-retorno}
\]
On adopte ici la convention \(0\log0:=0\). L'état est séparable
exactement lorsqu'une seule probabilité vaut un.
\end{proposition}

\begin{proof}
En effectuant la trace partielle, l'orthogonalité annule les termes avec
\(i\ne j\). Les valeurs propres non nulles de chaque matrice réduite sont donc
les \(p_i\), ce qui donne \eqref{eq:entropia-ida-retorno}. Le rang de
Schmidt est un exactement dans le cas déterministe.
\end{proof}

La proposition formalise, une fois la complétion choisie, une entropie entre
les lectures avancée et retardée. Elle ne déduit pas les probabilités \(p_i\),
n'identifie pas ces lectures aux propagateurs avancé et retardé d'une
théorie des champs et ne transforme pas une corrélation classique de voies en
intrication. La construction d'un isomorphisme qui envoie la mémoire HMT
sur des états physiques bipartites, préserve l'évolution et fixe \(p\) appartient à
l'\textsc{interface physique}. L'entropie nulle d'une distribution
déterministe constitue le contrôle négatif immédiat.

\begin{statusbox}
La proposition est
\textsc{démontrée avec une structure de départ explicite} : elle présuppose la
complétion de Hilbert, les bases orthonormales et la distribution \(p\). La construction
de l'opérateur d'entrelacement qui réalise physiquement les deux couches est une
\textsc{interface physique} ; elle ne fait pas partie de la proposition.
\end{statusbox}

\begin{definition}[Interface photonique de HMT--MD]
\label{def:interfaz-fotonica-hmt-md}
Une réalisation photonique de HMT--MD est une section ouverte portée par
une direction nulle de l'algèbre scindée, transportée entre des frontières
conjuguées et protégée par une symétrie qui conserve exactement la norme
nulle.
\end{definition}

Les idempotents \(P_\pm\) ne sont pas, à eux seuls, des photons physiques ni les deux
polarisations de Maxwell. La réalisation complète exige un
opérateur d'entrelacement
\[
 \mathfrak R_\gamma:
 \partial_{\rm null}\mathcal C_{\rm sp}^{+}
 \dashrightarrow\mathcal H_{\rm Maxwell}^{\rm phys}
\]
qui préserve le produit scalaire, l'évolution et l'action de jauge, et qui
sélectionne exactement deux polarisations transversales en \(3+1\)
dimensions.

\begin{criterion}[Secteur photonique]
\label{crit:sector-fotonico}
L'identification physique est réfutée si l'évolution quitte
\(\partial_{\rm null}\mathcal C_{\rm sp}^{+}\), engendre une masse sans le mécanisme
de brisure déclaré, ne reproduit pas la dispersion nulle ou ne laisse pas exactement
deux polarisations physiques après la réduction de jauge.
\end{criterion}

\section{Clôture matérielle et centre électronique}
\label{sec:clausura-material-electron}

Le couplage simultané de deux directions nulles conjuguées produit, pour
\(u,v>0\),
\[
 N_{\rm sp}(uP_++vP_-)=uv>0.
\]
Aucune composante isolée n'a de norme non nulle ; la norme intérieure
naît de leur couplage. Cet énoncé algébrique n'additionne pas des masses
préexistantes et n'identifie pas encore les directions nulles à des photons
\emph{sur couche de masse}.

   La classification structurelle de HMT contient un unique point fixe de l'inversion cellulaire,
\[
 (5,5).
\]
L'origine relative d'action de l'électron est \(v_e=0\). Dans la réalisation
scindée normalisée, on fixe
\[
 Z_e=P_++P_-=I,
 \qquad
 N_{\rm sp}(Z_e)=1.
\]
Le point \((5,5)\), l'origine \(v_e=0\), la signature brute
\((24,0,12,0,-12)\) et le mot transversal \(555555\) sont quatre objets
distincts.

\begin{definition}[Candidat de clôture électronique]
\label{def:candidato-clausura-electronica}
La clôture résonante de modes nuls électromagnétiques conjugués se
formalise comme le programme qui cherche une section matérielle centrale normalisée
obtenue par le couplage de deux modes nuls conjugués, stabilisée par
une clôture de voie et par son holonomie.
\end{definition}

L'existence de l'extension, du retour et du centre de l'atlas ne prouve pas
la compacité minimale ni la stabilité de la fermeture. Le théorème dynamique
restant à établir doit construire
\[
 \gamma_e=\xi_1\circ\cdots\circ\xi_N,
 \qquad
 \partial\gamma_e=0,
 \qquad
 \varepsilon(\gamma_e)=-1,
\]
où chaque \(\xi_j\) est localement nul, le bord total s'annule et
\(\varepsilon\) est l'holonomie de signe. La masse serait associée au cocycle
global de clôture, non à une somme de masses nulles locales. Ce programme ne
rouvre ni l'électron central \((5,5)\), ni l'atlas, ni la chaîne
extension--voie--registre--caractère déjà construits.

\begin{proposition}[Critère relativiste d'intérieur causal]
\label{prop:cierre-nulos-interior-causal}
Soient \(p_1,\ldots,p_N\) des impulsions nulles futures dans un espace de Minkowski de
signature temporelle positive, \(p_i^2=0\), et soit
\[
 P=\sum_{i=1}^{N}p_i.
\]
Si tous les \(p_i\) sont colinéaires, alors \(P^2=0\). Si au moins deux sont
futurs et non colinéaires, alors \(P^2>0\), et l'on peut définir
\[
 m^2c^2=P^2.
\]
\end{proposition}

\begin{proof}
Pour des vecteurs nuls futurs, on a
\(p_i\cdot p_j\geq0\), avec égalité exactement lorsqu'ils sont colinéaires.
Comme
\[
 P^2=\sum_i p_i^2+2\sum_{i<j}p_i\cdot p_j,
\]
la première somme est nulle. Tous les produits croisés s'annulent dans le cas
colinéaire ; s'il existe une paire non colinéaire, au moins l'un est strictement
positif.
\end{proof}

Ainsi, une voie fermée n'est pas massive par la seule topologie, et une voie
ouverte n'est pas nulle du seul fait qu'elle possède un bord. La réalisation exige que
la fermeture géométrique et la somme des impulsions entrent conjointement dans
l'intérieur causal.

\subsection{Möbius, conjugaison et occupation ultérieure}

Sous les hypothèses de relèvement et d'holonomie du
\cref{thm:dos-cuatro-pi-estadistica}, un tour inverse l'orientation et
deux la restaurent. Le ruban de Möbius organise la particule et l'antiparticule comme des
orientations conjuguées d'un même type de fermeture. Leur identification physique
exige une représentation qui inverse la charge, préserve la norme et
entrelace l'évolution. La géométrie de Möbius n'implique pas à elle seule
l'algèbre d'occupation fermionique ; ses relations opératoires sont
introduites comme hypothèses opératoires dans le chapitre suivant.

\begin{criterion}[Fermeture électronique]
\label{crit:cierre-electronico}
La réalisation est réfutée si le centre \((5,5)\) ne se relève pas en une
extension fermée, si l'holonomie ne sélectionne pas le secteur impair, si
l'hamiltonien ne possède pas de solution liée stable et unique à symétrie près, ou
si la charge, la dispersion et le facteur de forme ne concordent pas avec la réalisation
électronique.
\end{criterion}

\endgroup

\begingroup
\let\chapter\section
\let\section\subsection
\let\subsection\subsubsection
\let\subsubsection\paragraph
\let\clearpage\relax
\section{Réalisations physiques par objet, représentation et codomaine}
Chaque réalisation déclare son objet primaire, sa représentation ou sa fibre et son résultat natif. Une coordonnée scalaire n'est présentée que lorsque l'objet lui-même et son morphisme de réalisation l'imposent.


\endgroup
\begingroup
\let\chapter\section
\let\section\subsection
\let\subsection\subsubsection
\let\subsubsection\paragraph
\let\clearpage\relax

\clearpage
\section{Réalisations élémentaires du Modèle Standard}
\begingroup
\footnotesize
\setlength{\tabcolsep}{1.9pt}
\begin{longtable}{L{2.03cm}L{2.03cm}L{1.87cm}L{2.34cm}L{6.55cm}}
\toprule
\rowcolor{HMTNavy}\HMTtablehead{Particule} & \HMTtablehead{Charge} & \HMTtablehead{Spin} & \HMTtablehead{Statistique} & \HMTtablehead{Réalisation HMT--MD} \\
\midrule
\endfirsthead
\multicolumn{5}{l}{\sffamily\scriptsize\color{HMTGray}Suite du tableau}\\[-1pt]
\toprule
\rowcolor{HMTNavy}\HMTtablehead{Partícula} & \HMTtablehead{Carga} & \HMTtablehead{Espín} & \HMTtablehead{Estadística} & \HMTtablehead{Realización HMT--MD} \\
\midrule
\endhead
\midrule
\endfoot
\bottomrule
\endlastfoot
\(e^-\) & \(-1\) & \(1/\allowbreak{}2\) & fermiónica & punto fijo \((5,5)\), vacancia orientada \\
\rowcolor{HMTLight}\(\mu^-\) & \(-1\) & \(1/\allowbreak{}2\) & fermiónica & bloque leptónico de rango 8 \\
\(\tau^-\) & \(-1\) & \(1/\allowbreak{}2\) & fermiónica & bloque leptónico de rango 16 \\
\rowcolor{HMTLight}\(\nu_e\) & \(0\) & \(1/\allowbreak{}2\) & fermiónica & vector de interacción en soporte neutro 20 \\
\(\nu_\mu\) & \(0\) & \(1/\allowbreak{}2\) & fermiónica & vector de interacción en soporte neutro 20 \\
\rowcolor{HMTLight}\(\nu_\tau\) & \(0\) & \(1/\allowbreak{}2\) & fermiónica & vector de interacción en soporte neutro 20 \\
\(u\) & \(+2/\allowbreak{}3\) & \(1/\allowbreak{}2\) & fermiónica & rama up, generación 1, triplete de color \\
\rowcolor{HMTLight}\(d\) & \(-1/\allowbreak{}3\) & \(1/\allowbreak{}2\) & fermiónica & rama down, generación 1, triplete de color \\
\(c\) & \(+2/\allowbreak{}3\) & \(1/\allowbreak{}2\) & fermiónica & rama up, generación 2, triplete de color \\
\rowcolor{HMTLight}\(s\) & \(-1/\allowbreak{}3\) & \(1/\allowbreak{}2\) & fermiónica & rama down, generación 2, triplete de color \\
\(t\) & \(+2/\allowbreak{}3\) & \(1/\allowbreak{}2\) & fermiónica & rama up, generación 3, triplete de color \\
\rowcolor{HMTLight}\(b\) & \(-1/\allowbreak{}3\) & \(1/\allowbreak{}2\) & fermiónica & rama down, generación 3, triplete de color \\
\(\gamma\) & \(0\) & \(1\) & bosónica & sección nula, dos polarizaciones \\
\rowcolor{HMTLight}\(g^a\) & \(0\) & \(1\) & bosónica & adjunta de color, ocho generadores \\
\(W^\pm\) & \(\pm1\) & \(1\) & bosónica & bloque vectorial cargado conjugado \\
\rowcolor{HMTLight}\(Z^0\) & \(0\) & \(1\) & bosónica & bloque vectorial neutral \\
\(H\) & \(0\) & \(0\) & bosónica & sección escalar persistente \\
\end{longtable}
\endgroup
Las antipartículas se obtienen por \(C_M\). Los antiquarks pertenecen al triplete conjugado; \(W^-\) es el conjugado de \(W^+\); fotón, \(Z\), Higgs y gluones poseen las conjugaciones propias de sus representaciones.

El soporte B1 exacto de las clases elementales tiene rangos

\[
 1,\quad8,\quad16,\quad
 3\times20,\quad
 3\times16,\quad
 3\times20.
\]
Le nombre d'incidences classe--parcours est donc

\[
 1+8+16+60+48+60=193.
\]
Les incidences ne sont pas 193 parcours distincts : les trois saveurs neutriniques partagent le support neutre de rang vingt ; les saveurs up réutilisent un support de rang seize et les saveurs down, un support de rang vingt. Le projecteur de saveur doit raffiner ces supports. Dans l'union des quarks à 36 cellules, la couleur est exactement

\[
 12_R+12_G+12_B.
\]
Cette égalité certifie le support chromatique global, mais n'autorise pas encore à choisir une cellule par saveur.

\Needspace{7\baselineskip}
\section{Décomposition du secteur leptonique par génération, charge, spin et caractère de masse}
\Needspace{7\baselineskip}
\subsection{Électron : centre, lacunes et orientation}
L'électron est exceptionnel parce que sa fibre contient l'unique point fixe \((5,5)\). Sa réduction à une masse scalaire ne nécessite pas de briser une orbite. Le parcours actif est

\[
 e_{\rm act}=(2,2,5,-4,-3,-2)^2,
\]
tandis que l'enveloppe historique

\[
 e_{\rm env}=(4,3,5,-4,-3,-5)^2
\]
conserve le signe et le maximum absolu des deux feuillets. L'enveloppe retient le mot

\[
 \tau_e=+++---+++---,
\]
mais ne remplace pas l'histoire active.

Quatre données doivent demeurer distinctes :

\[
 \text{signature brute}=(24,0,12,0,-12),
\]
\[
 \text{signature paire CT--108}=(-12,-4,12,-6,12),
\]
\[
 \text{vecteur d'événements}=e_{\rm act},
\qquad
 \text{origine affine}=v_e=0.
\]
Le sélecteur de base \((-22,+15)\) possède deux lectures internes complémentaires. La décomposition régionale donne

\[
 -22=-(27-5),\qquad +15=5\cdot3=\binom62,
\]
et le cadre des lacunes

\[
 V_e=
 \begin{pmatrix}
 21&22\\
 6&7
 \end{pmatrix}
\]
satisfait

\[
 \det V_e=21\cdot7-22\cdot6=15.
\]
Le \(22\) mesure le déficit orienté de l'extension ; le \(15\), la rétention d'aire qui survit à la clôture. La particule électronique apparaît ainsi comme persistance centrale d'une lacune orientée, et non comme une masse insérée dans la cellule centrale.

L'équation complète est

\[
 \mathfrak e_e^\beta=
 \frac{\sqrt3}{4}
 \left[
 \exp\!\left(\frac{\varphi}{\pi^2}\right)
 -22\alpha_{\rm HMT}^{3}
 \right]
 (1+15\Delta_4)
 R_{\rm act}^{80-54\alpha_{\rm HMT}+6\Delta_4}.
\]
Dans la carte électronique,

\[
 \mathfrak e_e^\beta
 =0.5109989506900008086082571648\ldots\ {\rm MeV}.
\]
\Needspace{7\baselineskip}
\subsection{Réalisation du muon comme section leptonique de 2e génération : fibre, signature et caractère}
Le muon réside dans une multisection leptonique de rang huit. Son empreinte structurelle retrouvée comprend les cellules

\[
 \{21,23,25,39,43,57,59,61\},
\]
les dénombrements de rotation/changement \(22/25\), l'axe nord--sud,

\[
 (N,E,S,O)=(62,25,7,14),
\qquad
 (h_{369},r_5)=(57,9),
\]
et la coordonnée centrale

\[
 (n_A,s_C)=(42,-3).
\]
L'objet propre est l'opérateur comprimé

\[
 \widehat M_\mu=P_\mu\widehat M P_\mu.
\]
Sa réduction à un nombre exige de prouver que le projecteur de saveur et de génération commute avec l'opérateur, ou que l'état préparé a une variance nulle. Étant donnée la signature enrichie

\[
 (42,-3,8,0,2),
\]
l'évaluation du caractère est

\[
 105.658360294435829\ldots\ {\rm MeV}.
\]
Cette valeur est un contrôle exact de l'évaluation de cette signature ; la dérivation physique complète doit reproduire la signature à partir de l'enregistrement généalogique sélectionné par le couplage. La largeur muonique s'obtient à partir de l'opérateur temporel de désintégration, et non de \(\widehat M_\mu\).

\Needspace{7\baselineskip}
\subsection{Réalisation du tau comme section leptonique de 3e génération : fibre, signature et caractère}
Le tau occupe une multisection leptonique de rang seize. Son empreinte comprend

\[
 17/24,\qquad \text{axe est--ouest},\qquad
 (N,E,S,O)=(29,55,7,17),
\]
\[
 (h_{369},r_5)=(59,6),\qquad
 (n_A,s_C)=(64,0).
\]
La signature enrichie

\[
 (64,0,25,-1,6)
\]
produit

\[
 1776.8609176314323\ldots\ {\rm MeV}.
\]
Comme pour le muon, la structure établie est le bloc opératoriel ; la masse scalaire s'obtient lorsque le couplage tau sélectionne un bloc scalaire ou un pôle isolé.

\Needspace{7\baselineskip}
\subsection{Fibre neutrinique commune et séparation des saveurs par le lecteur PMNS}
Les trois saveurs neutriniques partagent l'espace

\[
 \mathcal H_\nu=
 \bigoplus_{j=1}^{4}\mathcal H_{\nu,G_j},
\qquad
 \dim\mathcal H_\nu=2+4+6+8=20.
\]
Les vecteurs

\[
 f_e,\quad f_\mu,\quad f_\tau
\]
constituent une base d'interaction ; ce ne sont pas trois vecteurs propres de masse. La profondeur \(G_j\) n'est pas non plus une quatrième saveur stérile : une interprétation stérile exige un projecteur physique supplémentaire.

L'extracteur neutre se définit par

\[
 T_\nu(m)=\sum_{t\in E_m}
 \omega(t)\varepsilon(t)\theta(t),
\qquad
 e_m^\nu=T_\nu(m)-T_\nu(m+6),
\]
\[
 \tau_m^\nu=\operatorname{sgn}(e_m^\nu).
\]
On en obtient

\[
 b_{120}^\nu
 =\sum_{a=0}^{3}
 \operatorname{sgn}\!\left(\sum_{m\in I_a}e_m^\nu\right),
\]
\[
 b_\zeta^\nu=\sum_m\tau_m^\nu\tau_{m+3}^\nu.
\]
La signature neutrinique prend la forme

\[
 v_\nu^{(i)}
 =(0,0,k_{\Delta_4}^{(i)},b_{120}^{\nu_i},b_\zeta^{\nu_i}).
\]
L'inversion bidirectionnelle change \(b_{120}\) et conserve \(b_\zeta\) ; cette asymétrie produit intérieurement les deux orientations de hiérarchie.

\Needspace{7\baselineskip}
\subsection{Nature fermionique du neutrino}
La neutralité ne détermine pas la statistique. Le secteur neutrinique se situe dans la complétion extérieure et porte l'holonomie spinorielle :

\[
 U_\nu(2\pi)=-I,\qquad U_\nu(4\pi)=I,
\]
\[
 \{a_{\nu,i},a_{\nu,j}^\dagger\}=\delta_{ij}.
\]
Chaque mode neutrinique physique est donc fermionique. La question Dirac/Majorana est ultérieure : elle détermine si la conjugaison échange deux sections distinctes ou identifie une section à sa conjuguée au sein d'une représentation réelle. Aucune des deux possibilités ne transforme le neutrino en boson.

\Needspace{7\baselineskip}
\subsection{Opérateur de masse et PMNS}
Soit \(P_\nu\) le projecteur de rang trois sur le sous-espace des modes massiques. On diagonalise d'abord

\[
 P_\nu\widehat M_\nu P_\nu
 =\sum_{i=1}^{3}m_i\,|n_i\rangle\langle n_i|.
\]
On construit ensuite le transport entre bases :

\[
 (f_e,f_\mu,f_\tau)
 =(n_1,n_2,n_3)U_{\rm PMNS}^\dagger.
\]
Un noyau de rang plein lisant feuillet, orientation, retenue, lacune et holonomie peut s'écrire

\[
 G_{ab}=
 \frac1{\sqrt{|L_a||N_b|}}
 \sum_{c\in L_a}\sum_{d\in N_b}K_{\rm reg}(c,d),
\]
\[
 U_{\rm PMNS}=\operatorname{polar}(G),
\qquad \det G\ne0.
\]
La phase interne de rang un produit les angles

\[
 (19.6807^\circ,56.4970^\circ,29.6224^\circ)
\]
et ne peut donc constituer le noyau PMNS physique complet. Le rang plein est une exigence mathématique, et non une retouche numérique.

Les observables

\[
 m_1,m_2,m_3,\qquad
 \Delta m_{ij}^2,\qquad
 m_\beta,\qquad
 m_{\beta\beta},\qquad
 \sum_i m_i
\]
sont des applications différentes du même spectre. Elles ne doivent pas être confondues dans une unique « masse du neutrino ».

\Needspace{7\baselineskip}
\section{Quarks, couleur, générations et mélange CKM}
\Needspace{7\baselineskip}
\subsection{Position structurelle des \(6\) saveurs de quarks}
Chaque quark se définit par

\[
 \mathfrak q_X=
 (\text{branche de charge},\text{génération},
 \mathbb Z_3^{\rm color},\text{multisection},
 \text{registre généalogique},\text{feuillet}).
\]
La position structurelle des six saveurs est résumée par leurs empreintes antérieures à toute carte de masse :

\Needspace{7\baselineskip}
\begingroup
\fontsize{7.2}{8.2}\selectfont
\setlength{\tabcolsep}{1.9pt}
\begin{longtable}{L{0.94cm}L{2.81cm}L{2.03cm}L{1.87cm}L{3.28cm}L{2.18cm}L{1.87cm}}
\toprule
\rowcolor{HMTNavy}\HMTtablehead{Saveur} & \HMTtablehead{Branche/\allowbreak{}génération} & \HMTtablehead{Giro/\allowbreak{}cambio} & \HMTtablehead{Eje} & \HMTtablehead{\((N,E,S,O)\)} & \HMTtablehead{\((h_{369},r_5)\)} & \HMTtablehead{\((n_A,s_C)\)} \\
\midrule
\endfirsthead
\multicolumn{7}{l}{\sffamily\scriptsize\color{HMTGray}Suite du tableau}\\[-1pt]
\toprule
\rowcolor{HMTNavy}\HMTtablehead{Sabor} & \HMTtablehead{Rama/\allowbreak{}generación} & \HMTtablehead{Giro/\allowbreak{}cambio} & \HMTtablehead{Eje} & \HMTtablehead{\((N,E,S,O)\)} & \HMTtablehead{\((h_{369},r_5)\)} & \HMTtablehead{\((n_A,s_C)\)} \\
\midrule
\endhead
\midrule
\endfoot
\bottomrule
\endlastfoot
\(u\) & positiva/\allowbreak{}1 & \(21/\allowbreak{}30\) & N--S & \((34,31,31,12)\) & \((51,13)\) & \((11,7)\) \\
\rowcolor{HMTLight}\(d\) & negativa/\allowbreak{}1 & \(20/\allowbreak{}32\) & E--O & \((27,35,18,28)\) & \((47,17)\) & \((17,8)\) \\
\(c\) & positiva/\allowbreak{}2 & \(21/\allowbreak{}29\) & E--O & \((20,38,24,26)\) & \((49,16)\) & \((61,8)\) \\
\rowcolor{HMTLight}\(s\) & negativa/\allowbreak{}2 & \(17/\allowbreak{}27\) & N--S & \((50,32,12,14)\) & \((57,11)\) & \((41,-3)\) \\
\(t\) & positiva/\allowbreak{}3 & \(21/\allowbreak{}26\) & N--S & \((46,27,16,19)\) & \((57,9)\) & \((100,-1)\) \\
\rowcolor{HMTLight}\(b\) & negativa/\allowbreak{}3 & \(21/\allowbreak{}30\) & E--O & \((12,63,9,24)\) & \((47,7)\) & \((71,-5)\) \\
\end{longtable}
\endgroup
Estas posiciones no son seis números elegidos: son seis clases de representación que actúan sobre tripletes de color. La ruta física puede ser una órbita o un bloque de rutas; no tiene que ser un punto.

\Needspace{7\baselineskip}
\subsection{Invarianza escalar del carácter de masa bajo la acción de color}
Sea \(P_q\) el proyector de sabor y generación y \(P_{\rm RGB}\) el proyector a la órbita de color. La masa quark es

\[
 \widehat M_q=
 P_qP_{\rm RGB}\widehat M P_{\rm RGB}P_q.
\]
Si

\[
 [\widehat M_q,U(g)]=0,\qquad g\in SU(3),
\]
alors

\[
 \widehat M_q=m_qI_3
\]
sur le triplet irréductible. Cette égalité explique pourquoi la saveur porte une masse commune bien que ses trois couleurs occupent des positions distinctes dans l'atlas.

\Needspace{7\baselineskip}
\subsection{Carte du schéma et de l'échelle}
La coordonnée HMT interne doit être transportée vers une définition renormalisée :

\[
 m_q^{\mathcal S}(\mu)
 =Z_m^{\mathcal S}(\mu,\mu_0)\,
 \mathcal R_q(\widehat M_q),
\]
avec la loi de composition

\[
 Z_m(\mu_2,\mu_0)
 =Z_m(\mu_2,\mu_1)Z_m(\mu_1,\mu_0).
\]
Pour \(u,d,s\), la comparaison naturelle utilise \(\overline{\rm MS}\) à \(2\,{\rm GeV}\) ; pour \(c,b\), l'échelle de la masse renormalisée elle-même ; pour \(t\), il faut déclarer si la carte est de pôle, MSR ou cinématique. Sans ce transport, une différence numérique n'est pas encore un résidu de la loi.

\Needspace{7\baselineskip}
\subsection{Évaluations des signatures déclarées}
Les coordonnées centrales historiques produisent :

\Needspace{7\baselineskip}
\begingroup
\footnotesize
\setlength{\tabcolsep}{2.1pt}
\begin{longtable}{L{4.68cm}L{10.61cm}}
\toprule
\rowcolor{HMTNavy}\HMTtablehead{Saveur} & \HMTtablehead{Évaluation} \\
\midrule
\endfirsthead
\multicolumn{2}{l}{\sffamily\scriptsize\color{HMTGray}Suite du tableau}\\[-1pt]
\toprule
\rowcolor{HMTNavy}\HMTtablehead{Sabor} & \HMTtablehead{Evaluación} \\
\midrule
\endhead
\midrule
\endfoot
\bottomrule
\endlastfoot
\(u\) & \(2.165924\allowbreak{}536173\allowbreak{}907\ldots\ {\rm MeV}\) \\
\rowcolor{HMTLight}\(d\) & \(4.679550\allowbreak{}162948\allowbreak{}874\ldots\ {\rm MeV}\) \\
\(c\) & \(1.270500\allowbreak{}838641\allowbreak{}911\ldots\ {\rm GeV}\) \\
\rowcolor{HMTLight}\(s\) & \(92.940093\allowbreak{}907845\allowbreak{}64\ldots\ {\rm MeV}\) \\
\(t\) & \(172.597479\allowbreak{}544019\allowbreak{}1\ldots\ {\rm GeV}\) \\
\rowcolor{HMTLight}\(b\) & \(4.190119\allowbreak{}763299\allowbreak{}790\ldots\ {\rm GeV}\) \\
\end{longtable}
\endgroup
Estas evaluaciones comprueban el carácter una vez dada la firma. La deducción completa exige obtener cada firma, su refinamiento de torre y su carta de esquema a partir del proyector de sabor y del registro genealógico correspondiente.

\Needspace{7\baselineskip}
\subsection{Determinación de los ángulos CKM desde la reflexión racional de las rutas quark}
Sean

\[
 A=1000\alpha_{\rm HMT},\qquad
 h=\frac{C^*}{6},\qquad
 \Delta=\frac{180}{\pi}\Delta_4.
\]
La grammaire angulaire donne

\[
 \begin{pmatrix}
 \theta_{12}\\
 \theta_{23}\\
 \theta_{13}
 \end{pmatrix}
 =
 \begin{pmatrix}
 2&-5&28\\
 0&7&-25/2\\
 1&-20&-2/3
 \end{pmatrix}
 \begin{pmatrix}
 A\\h\\\Delta
 \end{pmatrix},
\qquad
 \delta_{\rm CP}=9A.
\]
Le critère linéaire exact de réfutation est

\[
 3857\,\delta_{\rm CP}
 =9\bigl(
 1528\theta_{12}
 +3380\theta_{23}
 +801\theta_{13}
 \bigr).
\]
À partir des trois angles et de \(\delta_{\rm CP}\), on construit la matrice unitaire \(V_{\rm CKM}\). Sa fonction dans la loi des masses n'est pas de produire des valeurs propres, mais de transporter les opérateurs down vers la base up :

\[
 B_d^{(u)}=V_{\rm CKM}B_dV_{\rm CKM}^*.
\]
\Needspace{7\baselineskip}
\subsection{Invariant de Jarlskog}
L'invariant de violation CP est

\[
 J_{\rm CKM}
 =c_{12}c_{23}c_{13}^2
 s_{12}s_{23}s_{13}\sin\delta_{\rm CP},
\]
et l'évaluation interne est

\[
 J_{\rm CKM}
 =3.1189723326632947\times10^{-5}.
\]
De manière équivalente, la compatibilité entre spectres et mélange peut se lire dans le déterminant du commutateur des opérateurs hermitiens up/down. Une valeur non nulle de \(J_{\rm CKM}\) enregistre une orientation physique non éliminable ; ce n'est pas un correcteur de masse ajouté à chaque quark.

\Needspace{7\baselineskip}
\section{Secteurs de jauge et scalaire}
\Needspace{7\baselineskip}
\subsection{Carte du secteur nul : parcours de masse \(0\), représentations de jauge et persistance}
Soit

\[
 \mathcal A_{\rm sp}=\mathbb R[R]/(R^2-1),
 \qquad
 P_\pm=\frac{1\pm R}{2}.
\]
Pour

\[
 x=r_+P_++r_-P_-,
\]
la norme déployée est

\[
 N_{\rm sp}(x)=r_+r_-.
\]
La section nulle est une direction de norme nulle. Sa masse est exactement

\[
 m=0,
\]
sans passer par une limite de caractères positifs et sans attribuer \(\kappa=0\). La nullité appartient au type de la section.

\Needspace{7\baselineskip}
\subsection{Réalisation du photon dans le secteur nul électromagnétique : hélicité, jauge et caractère}
Le photon est la réalisation vectorielle transverse d'une section de jauge nulle. Le projecteur physique doit laisser exactement deux polarisations :

\[
 \operatorname{rg}P_\gamma^{\rm trans}=2,
\qquad
 P_\gamma^{\rm trans}k=0.
\]
La masse au repos est

\[
 m_\gamma=0.
\]
La propagation bidirectionnelle peut porter une anisotropie orientée

\[
 \epsilon=q_-^{90}-q_+^{90},
\qquad
 c_\pm=\frac{c}{1\pm\epsilon},
\]
tandis que la moyenne harmonique conserve

\[
 \frac{2}{1/c_++1/c_-}=c.
\]
Ainsi, l'aller et le retour peuvent différer sans modifier l'invariant two--way. Ce biais n'est pas une masse de photon.

Une réalisation Proca/Stückelberg distincte apparaît si un gap spectral de l'enregistrement généalogique satisfait

\[
 \omega^2=k^2+m_T^2
\]
et s'il existe un entrelaceur dimensionnel

\[
 m_T\longmapsto m_\gamma.
\]
Le retard induit obéirait à

\[
 \frac{\Delta T}{T}
 \simeq\frac12
 \left(\frac{m_\gamma c^2}{\hbar\omega}\right)^2.
\]
Le secteur Maxwell nul et la réalisation Proca massive sont deux branches distinctes ; l'une n'invalide pas l'autre.

\Needspace{7\baselineskip}
\subsection{Réalisation chromodynamique du secteur nul : octet de gluons et représentation adjointe de \(SU(3)\)}
Le trit résiduel engendre l'espace de couleur, dont la partie sans trace réalise

\[
 \mathfrak{sl}_3\longrightarrow\mathfrak{su}_3.
\]
Le gluon occupe la représentation adjointe de dimension huit. Tant que la symétrie de couleur demeure exacte,

\[
 m_g=0.
\]
Confinement, écrantage et gap spectral collectif sont des propriétés du spectre interactif ; ils ne constituent pas une masse de pôle élémentaire du gluon. Une masse effective doit préciser l'observable, la jauge, le régime et l'opérateur qui la définissent.

\Needspace{7\baselineskip}
\subsection{Réalisation des bosons \(W^\pm\) : charge, polarisation et caractère électrofaible}
Le secteur \(W\) est un bloc vectoriel chargé avec conjugaison

\[
 C_M:W^+\longleftrightarrow W^-.
\]
Son empreinte historique est

\[
 15/37,\qquad \text{axe est--ouest},\qquad
 (n_A,s_C)=(94,-1),
\]
et la signature enrichie déclarée est

\[
 (94,-1,0,-2,4).
\]
L'évaluation correspondante est

\[
 80.37896490970455\ldots\ {\rm GeV}.
\]
La réalisation complète est le pôle de la résolvante du bloc \(P_W\widehat M P_W\) ; sa partie imaginaire détermine la largeur. La dégénérescence de \(W^+\) et de \(W^-\) procède de la conjugaison de l'opérateur.

\Needspace{7\baselineskip}
\subsection{Réalisation du boson \(Z^0\) : mélange neutre, polarisation et caractère électrofaible}
Le \(Z^0\) est le bloc vectoriel neutre du couplage électrofaible. Son empreinte historique est

\[
 20/23,\qquad \text{axe est--ouest},\qquad
 (n_A,s_C)=(95,-1),
\]
et la signature enrichie

\[
 (95,-1,-11,0,-4).
\]
L'évaluation du caractère donne

\[
 91.18753344431341\ldots\ {\rm GeV}.
\]
Le projecteur neutre doit être construit avant la carte de pôle ; masse et largeur sont deux coordonnées du même pôle, et non une seule valeur numérique.

\Needspace{7\baselineskip}
\subsection{Réalisation du secteur de Higgs : mode scalaire, couplage et caractère de masse}
Le Higgs est une section scalaire persistante du bloc de jauge, de parité CP déclarée. Son empreinte est

\[
 24/29,\qquad \text{axe nord--sud},\qquad
 (n_A,s_C)=(97,9).
\]
L'évaluation conditionnelle du caractère produit

\[
 125.2924660425361\ldots\ {\rm GeV}.
\]
La définition HMT n'identifie pas « scalaire » à « masse » : le spin zéro fixe la représentation des rotations ; la masse procède du spectre du bloc scalaire couplé. La largeur et les canaux de désintégration exigent l'opérateur temporel.

\Needspace{7\baselineskip}
\subsection{Le moment magnétique anomal du muon}
La fonctionnelle interne

\[
 a_\mu=
 \frac{\alpha_{\rm HMT}}{2\pi}
 +\frac{\alpha_{\rm HMT}(\varphi-1)}{1000}
 +\frac{\alpha_{\rm HMT}^2}{1000(54+1)}
\]
donne

\[
 a_\mu=0.001165920713008014\ldots,
\]
et

\[
 g_\mu=2(1+a_\mu)
 =2.002331841426016028\ldots.
\]
Les trois termes correspondent respectivement à la rotation de base, au retour orienté dans la complétion \(1000\) et à la correction quadratique médiée par le défaut \(54\). Cet observable vérifie la structure interne du secteur muonique, mais ne sélectionne pas, à lui seul, son parcours de masse.

\Needspace{7\baselineskip}
\section{Construction des particules composées par composition non additive des parcours et des liaisons}
\Needspace{7\baselineskip}
\subsection{Opérateur de composition}
La composition des parcours n'est pas strictement additive, car frontière, ordre temporel et retenue interagissent. Pour deux enregistrements,

\[
 \mathcal L_1\star\mathcal L_2
 =\mathcal L_1+\mathcal L_2
 +\mathfrak b(\mathcal L_1,\mathcal L_2),
\]
où \(\mathfrak b\) est le cocycle de liaison. L'associativité exige

\[
 \mathfrak b(L_1,L_2)
 +\mathfrak b(L_1\star L_2,L_3)
 =
 \mathfrak b(L_2,L_3)
 +\mathfrak b(L_1,L_2\star L_3).
\]
Le lecteur de masse agit après cette composition :

\[
 (\Gamma_i,\mathcal T,\partial)
 \xrightarrow{\mathfrak C_{12}}
 \mathcal L_{\rm comp}
 \xrightarrow{L_{\rm tick}}
 \sigma_{\rm comp}
 \xrightarrow{\chi_{\rm full}}
 \widehat M_{\rm comp}.
\]
\Needspace{7\baselineskip}
\subsection{Réalisation mésonique par composition quark–antiquark et loi cocyclique de liaison}
Le domaine mésonique contient 432 expressions orientées

\[
 M(q,q')=q\otimes(q')^\vee
\]
à couleur compatible. Le projecteur physique impose le singulet de couleur et fixe

\[
 (J^{PC},I,I_3,S,C,B,\ldots),
\]
avec l'arbre temporel de liaison. La masse d'un méson n'est pas \(m_q+m_{\bar q}\) ; c'est une valeur propre de l'enregistrement composite.

Les six évaluations historiques de contrôle comprennent

\[
 m_{\pi^0}=134.9767243278721\ldots\ {\rm MeV},
\]
\[
 m_{\pi^\pm}=139.5704851715722\ldots\ {\rm MeV},
\]
\[
 m_{K^\pm}=493.6770856495306\ldots\ {\rm MeV},
\]
\[
 m_{K^0}=497.6111864977505\ldots\ {\rm MeV}.
\]
Ces valeurs sont des évaluations exactes d'enregistrements déclarés. La clôture généalogique état par état exige de régénérer chaque enregistrement au moyen de \(\mathfrak C_{12}\), y compris sa frontière de liaison.

\Needspace{7\baselineskip}
\subsection{Réalisation baryonique par composition trilinéaire des parcours de quarks et neutralisation de couleur}
Le domaine baryonique contient 1728 expressions RGB. Le projecteur physique doit imposer simultanément :

\begin{enumerate}
\item le singulet de couleur ;
\item la statistique fermionique ;
\item la saveur et l'isospin ;
\item le spin et la symétrisation ;
\item l'arbre triple avec retenue et mémoire.
\end{enumerate}
Pour le proton et le neutron, les signatures historiques

\[
 \sigma_p=(59,0,9,1,0),\qquad
 \sigma_n=(59,1,-34,0,1)
\]
donnent

\[
 m_p=938.2717515469849\ldots\ {\rm MeV},
\]
\[
 m_n=939.5658104725070\ldots\ {\rm MeV}.
\]
La précision apparente ne remplace pas la dérivation de l'enregistrement généalogique de liaison. Le signe et la grandeur de

\[
 m_n-m_p
\]
doivent découler de l'orientation, de l'isospin, des lacunes et de la frontière de liaison avant l'ouverture de la comparaison externe.

\Needspace{7\baselineskip}
\subsection{États excités et résonances hadroniques}
Une identité composée peut posséder plusieurs pôles. L'état physique est individualisé par

\[
 (\text{constituants},\mathcal T,\partial,
 J^{PC},I,\text{radial},\text{orbital},\text{canal}).
\]
La même composition de saveurs ne suffit pas à distinguer les excitations. Chaque pôle

\[
 z_k=M_k-\frac{i}{2}\Gamma_k
\]
doit procéder d'un vecteur propre de l'opérateur composé et d'un canal de désintégration. Cela étend la loi de quelques masses stables à tout le spectre des mésons et des baryons, sans transformer le catalogue externe en générateur.

\Needspace{7\baselineskip}
\subsection{Inventaire exhaustif des largeurs et des taux de transition des réalisations de l'atlas}
Soit \(P\) le sous-espace préparé et \(Q=I-P\). Le Hamiltonien effectif est

\[
 H_{\rm eff}(z)
 =PHP+PHQ(z-QHQ)^{-1}QHP.
\]
Les pôles de

\[
 \det(z-H_{\rm eff}(z))=0
\]
déterminent masses et largeurs. Le caractère de masse fournit la partie réelle de base de l'opérateur ; le terme d'autoénergie de \(Q\) fournit déplacement et désintégration. Les 384 largeurs de l'inventaire exigent donc un lecteur de canaux supplémentaire au lecteur des 471 masses.

\Needspace{7\baselineskip}
\section{Secteurs topologiques et particules hypothétiques}
\Needspace{7\baselineskip}
\subsection{Secteur topologique de Fibonacci : anneau basé, fusion et représentation de tresses}
L'incidence surfacique--volumétrique

\[
 F_{\rm av}=
 \begin{pmatrix}
 0&1\\1&1
 \end{pmatrix}
\]
réalise l'anneau

\[
 \tau^2=1+\tau,\qquad
 \operatorname{FPdim}\tau=\varphi.
\]
Le résultat primaire est une règle de fusion et une dimension quantique. Une particule anyonique de Fibonacci exige en outre des matrices \(F\) et \(R\), un pentagone, un hexagone, une chiralité, un Hamiltonien à gap spectral et une application de l'état TPK vers l'espace physique. Ce n'est qu'alors que l'on peut définir

\[
 m_{\rm brecha espectral}=\frac{\Delta E}{c_{\rm eff}^2}.
\]
\Needspace{7\baselineskip}
\subsection{Secteur d'Ising–Majorana : objets simples, règles de fusion et représentation de tresses}
La catégorie d'Ising satisfait

\[
 \psi^2=1,\qquad
 \psi\sigma=\sigma,\qquad
 \sigma^2=1+\psi,\qquad
 d_\sigma=\sqrt2.
\]
Quatre niveaux doivent être distingués :

\begin{enumerate}
\item l'involution des parcours \(C_M=-\operatorname{rev}\) ;
\item la parité de la complétion CAR ;
\item un spineur relativiste autoconjugué ;
\item un mode zéro topologique d'Ising.
\end{enumerate}
Un fermion relativiste de Majorana exige une représentation spinorielle réelle, une dynamique et un terme de masse. Un mode zéro de Majorana se caractérise par un gap spectral, une séparation de niveaux et une protection ; il ne possède pas de masse universelle déductible du seul anneau de fusion. L'identité d'un parcours avec son conjugué est une condition nécessaire, mais non suffisante, de la réalisation de Majorana.

\Needspace{7\baselineskip}
\subsection{Représentation de tresses sur \(\mathbb Z/6\mathbb Z\) et caractères topologiques}
Soit

\[
 w_6:B_n\longrightarrow\mathbb Z/6,
\qquad
 \chi_k(b)=e^{2\pi ikw(b)/6}.
\]
Les représentations

\[
 \rho_q(\sigma_i)=qP_i
\]
descendent au groupe symétrique seulement lorsque \(q^2=1\). Les quatre autres valeurs sont véritablement tressées. Leur existence mathématique n'implique pas quatre particules : une réalisation physique exige un Hamiltonien, un gap spectral, des opérateurs de tresse, une stabilité et un entrelaceur avec une multisection.

\Needspace{7\baselineskip}
\subsection{Impossibilité d'un tachyon persistant}
La droite d'action est positive,

\[
 \hbar_{\rm int}>0,
\]
le caractère satisfait

\[
 \chi_{\rm full}>0,
\]
et l'opérateur bêta s'obtient par conjugaison positive :

\[
 \widehat M^\beta
 =R_{\rm act}^{\widehat K/2}
 \widehat M^{(0)}
 R_{\rm act}^{\widehat K/2}\ge0.
\]
Par conséquent,

\[
 \operatorname{Spec}(\widehat M^\beta)\subset[0,\infty).
\]
Une section persistante stable ne peut engendrer

\[
 m^2<0.
\]
Si la linéarisation d'un état produit une courbure négative ou une fréquence imaginaire, cet état est une instabilité qui quitte la section, et non une particule supraluminique. Si le prolongement se termine, c'est une section virtuelle à obstruction terminale. Ainsi, « tachyon » se décompose en deux objets HMT correctement typés --- instabilité ou virtualité --- et ne constitue pas une espèce persistante.

\Needspace{7\baselineskip}
\subsection{Gravitation sans particule élémentaire obligatoire de spin \(2\)}
La réalisation gravitationnelle suit la chaîne

\[
 \text{holonomie de route}
 \longrightarrow(e,\omega)
 \longrightarrow(T,R)
 \longrightarrow
 S_{\rm PCH},
\]
où \(e\) est la costructure, \(\omega\) la connexion,

\[
 T=de+\omega\wedge e,\qquad
 R=d\omega+\omega\wedge\omega,
\]
et \(S_{\rm PCH}\) est l'action de Palatini--Cartan--Holst. La gravitation apparaît comme réponse globale de torsion et de courbure du réseau persistant. La vibration collective de l'univers modèle peut posséder une composante tensorielle de spin deux, mais HMT n'a pas besoin de postuler une particule élémentaire supplémentaire pour transporter la gravitation.

La constante gravitationnelle interne est typée par

\[
 G_{\rm HMT}
 =\frac{c_{\rm int}^3}{\hbar_{\rm int}}
 \bigl(\delta_\theta\Lambda_G\bigr)^2,
\]
où \(\delta_\theta\) est la torsion angulaire élémentaire et \(\Lambda_G\) la longueur spectrale sélectionnée par l'holonomie gravitationnelle. Cette équation convertit torsion et échelle de parcours en couplage gravitationnel sans utiliser de masse comme normalisation. La vibration totale est représentée par le Hessien

\[
 \mathcal H_{\rm univ}
 =\operatorname{Hess}_{(e,\omega)}S_{\rm PCH}
\]
sur l'état global. Ses modes sont des excitations collectives de l'univers modèle ; une composante tensorielle n'est pas automatiquement une particule séparée.

Une excitation appelée « graviton » doit démontrer trois propriétés :

\begin{enumerate}
\item un sous-espace tensoriel persistant ;
\item un pôle isolé de l'opérateur linéarisé ;
\item un couplage physique qui ne double pas la torsion déjà présente.
\end{enumerate}
Sans ces trois conditions, l'objet correct est le mode collectif de la géométrie, et non une nouvelle ligne de masse.

\Needspace{7\baselineskip}
\subsection{Barbero–Immirzi et la liaison hexade–octade}
Avec

\[
 A=1000\alpha_{\rm HMT},
\qquad
 q_\pm=\exp\!\left[-\frac{\pi}{180}(A\pm C^*)\right],
\]
la fonctionnelle hexade--octade utilise deux incidences distinctes des tours temporelles. La combinaison

\[
 \varepsilon_{\rm BI}
 =12\Delta S_{90}-\Delta S_{120}
\]
et le terme angulaire produisent

\[
 \gamma_{\rm BI}
 =\frac{\pi}{180}AC^*
 +\varepsilon_{\rm BI}
 =0.2740076726944593\ldots.
\]
Le poids \(12:-1\) procède de la lecture des douze secteurs dodécaphasiques
et de l'unique couture orientée de retour. L'équation
\(4a-3b=45\) vérifie ensuite que la paire \((12,1)\) est la solution
entière positive minimale. Cet objet ne doit pas être confondu avec l'extracteur dodécaphasique de masse ni avec le résidu angulaire \(90\)--\(120\) ; il partage l'opérateur harmonique, mais possède un autre domaine et une autre fonction.

\Needspace{7\baselineskip}
\subsection{Extensions de l'atlas vers les secteurs sombres : domaines, lecteurs et statut physique}
Un secteur sombre HMT est une multisection persistante avec

\[
 P_{\rm vis}P_{\rm dark}=0
\]
pour les couplages électromagnétique et chromatique, mais avec une incidence gravitationnelle ou de frontière possible :

\[
 P_{\rm grav}P_{\rm dark}\ne0.
\]
Son « obscurité » est une propriété du noyau de couplage, et non une étiquette de masse. La loi prédit d'abord un opérateur et son spectre ; seul un canal physique détermine si le mode est stable, métastable ou résonant. Cette définition permet de rechercher des multisections invisibles sans inventer une masse à partir de la valeur observée d'une densité cosmologique.

\Needspace{7\baselineskip}
\subsection{Section neutrinique stérile : compatibilité de fibre, mélange et absence de charge de jauge}
Un neutrino stérile n'est pas la profondeur \(G_4\). C'est un vecteur propre \(n_s\in\mathcal H_\nu\) qui satisfait

\[
 P_{\rm weak}n_s=0,\qquad
 P_{\rm mass}n_s=n_s.
\]
Il peut se mélanger aux trois modes actifs au moyen d'un noyau neutre étendu, mais son existence exige

\[
 \dim\ker(P_{\rm weak}|_{\mathcal H_\nu})>0
\]
et une valeur propre de masse dans ce noyau. Cette définition produit un critère interne de recherche : diagonaliser d'abord le bloc neutre, puis vérifier son couplage faible.

\Needspace{7\baselineskip}
\subsection{Sections axionique et pseudoscalaires : caractère topologique et couplage de bord}
Un candidat axionique est une section pseudoscalaire dont l'orientation change sous CP,

\[
 CP\,a\,CP^{-1}=-a,
\]
et dont le déplacement compense une phase topologique du secteur de couleur. Dans HMT, il doit correspondre à un mode impair de torsion/holonomie et à un projecteur qui entrelace ce mode avec la densité topologique de jauge. Sa masse s'obtient à partir de la dérivée seconde du potentiel effectif,

\[
 m_a^2=
 \left.\frac{\partial^2V_{\rm eff}}{\partial a^2}\right|_{a=a_0},
\]
une fois \(V_{\rm eff}\) construit à partir de l'enregistrement généalogique des lacunes et de la frontière. La seule existence d'un canal \(C^*\mapsto-C^*\) ne suffit pas à affirmer l'existence d'une particule axionique.

\Needspace{7\baselineskip}
\subsection{Photon sombre et vecteur massif caché}
Un photon sombre exige une seconde section de jauge \(A'\) et un couplage de mélange

\[
 \mathcal L_{\rm mix}
 =-\frac{\varepsilon_{\rm kin}}2F_{\mu\nu}F'^{\mu\nu}.
\]
En termes de fibres, \(\varepsilon_{\rm kin}\) est un entrelaceur entre deux sous-espaces nuls. Si le second secteur acquiert un gap spectral de Stückelberg ou de frontière,

\[
 \widehat M_{A'}^2\ne0,
\]
un vecteur massif sombre apparaît. HMT prédit que le paramètre de mélange doit procéder d'une incidence de mémoire commune ; il ne peut être fixé avec la masse que l'on cherche à expliquer.

\Needspace{7\baselineskip}
\subsection{WIMP et état sombre stable}
Un WIMP est une section sombre qui satisfait en outre :

\[
 [P_{\rm weak},P_X]\ne0,\qquad
 [P_{\rm em},P_X]=[P_{\rm color},P_X]=0,
\]
et possède un mode stable

\[
 |\lambda_X|=1.
\]
La stabilité peut être protégée par une parité de l'enregistrement généalogique. Sa masse est une valeur propre du bloc sombre faiblement couplé. « WIMP » ne constitue pas une nouvelle famille mathématique : c'est une condition conjointe de couplage et de persistance sur une multisection.

\Needspace{7\baselineskip}
\subsection{Section monopolaire : charge magnétique, condition de Dirac et compatibilité topologique}
Un monopôle se définit par une classe topologique non triviale du fibré de jauge :

\[
 \frac1{2\pi}\int_{S^2}F=n\ne0.
\]
HMT doit réaliser cet entier comme holonomie nette d'une famille fermée de parcours. Sa masse, si le mode est dynamique et stable, appartient au spectre de l'opérateur de défaut topologique. Une lacune ponctuelle ou une retenue non nulle ne constituent pas, à elles seules, un monopôle ; la descente à la classe entière doit être démontrée.

\Needspace{7\baselineskip}
\subsection{Extensions supersymétriques de l'atlas : appariement des représentations et des caractères}
Une supersymétrie HMT exigerait un opérateur impair \(Q\) échangeant les complétions extérieure et symétrique :

\[
 Q:\mathcal F_-\leftrightarrow\mathcal F_+,
\qquad
 Q^2=H,
\]
et entrelaçant la loi des masses,

\[
 Q\widehat M_-=\widehat M_+Q.
\]
Ce n'est qu'alors que les spectres bosonique et fermionique seraient appariés. L'existence séparée de CAR et de CCR ne démontre pas la supersymétrie et ne prédit pas, à elle seule, des sélectrons, neutralinos ou gluinos. Chaque partenaire exige un entrelaceur non nul et une règle de brisure déterminant sa séparation de masse.

\Needspace{7\baselineskip}
\subsection{Tétraquarks, pentaquarks, boules de glu et hybrides}
La grammaire de composition se prolonge sans changer de loi :

\[
 \mathfrak C_{12}^{(n)}:
 \Gamma_1\times\cdots\times\Gamma_n
 \times\mathcal T_n\times\partial_n
 \longrightarrow\mathcal L_{\rm comp}.
\]
Un tétraquark utilise \(qq\bar q\bar q\) ; un pentaquark, \(qqqq\bar q\) ; une boule de glu, des enregistrements de l'adjointe de couleur projetés sur le singulet ; un hybride, des parcours de quarks et de jauge dans un même arbre. La classification physique dépend du projecteur de couleur, du spin, de la parité et de l'arbre de liaison, et non de la seule liste des constituants.

\Needspace{7\baselineskip}
\subsection{Sections de dilaton et de radion : modes scalaires de volume et de compactification}
Le dilaton est un mode de la section d'échelle ; le radion, un mode qui modifie une longueur ou un volume interne. Dans la théorie des masses, ils apparaissent comme fluctuations de la carte et de la mémoire géométrique :

\[
 t_0\mapsto e^{\phi_d}t_0,\qquad
 \mathcal V_{\rm int}\mapsto e^{\phi_r}\mathcal V_{\rm int}.
\]
Pour devenir des particules, ils doivent posséder un terme cinétique, un potentiel, un projecteur persistant et un pôle. Tant qu'ils ne les possèdent pas, ce sont des degrés collectifs de la carte dimensionnelle, et non des espèces scalaires supplémentaires.

\Needspace{7\baselineskip}
\subsection{Exclusion des fantômes physiques}
La mesure spectrale d'un état physique est positive :

\[
 \mu_{P,a}^{\rm mass}(B)\ge0.
\]
Une direction de norme négative n'appartient pas à l'espace des états préparé. Si une paramétrisation introduit un « fantôme », la réalisation physique doit quotienter la redondance de jauge ou rejeter le couplage. Cette exclusion est indépendante de l'interdiction tachyonique : un fantôme viole la positivité de la norme ; un tachyon persistant violerait la positivité de l'opérateur de masse.

\Needspace{7\baselineskip}
\section{Inventaire universel des masses et des largeurs}
\Needspace{7\baselineskip}
\subsection{Correspondance ultérieure entre les réalisations HMT–MD et le dénombrement physique externe}
L'inventaire de comparaison contient

\[
 1170\ \text{identités},\qquad
 438\ \text{groupes d'identité},
\]
\[
 471\ \text{observables de masse},\qquad
 384\ \text{observables de largeur}.
\]
Sa partition est

\Needspace{7\baselineskip}
\begingroup
\footnotesize
\setlength{\tabcolsep}{2.1pt}
\begin{longtable}{L{5.04cm}L{5.04cm}L{5.04cm}}
\toprule
\rowcolor{HMTNavy}\HMTtablehead{Domaine} & \HMTtablehead{Masses} & \HMTtablehead{Largeurs} \\
\midrule
\endfirsthead
\multicolumn{3}{l}{\sffamily\scriptsize\color{HMTGray}Suite du tableau}\\[-1pt]
\toprule
\rowcolor{HMTNavy}\HMTtablehead{Dominio} & \HMTtablehead{Masas} & \HMTtablehead{Anchuras} \\
\midrule
\endhead
\midrule
\endfoot
\bottomrule
\endlastfoot
partículas generales y leptónicas & 53 & 9 \\
\rowcolor{HMTLight}quarks & 8 & 1 \\
mesones & 243 & 225 \\
\rowcolor{HMTLight}bariones & 166 & 149 \\
registro gravitatorio de contraste & 1 & 0 \\
\rowcolor{HMTLight}\textbf{Total} & \textbf{471} & \textbf{384} \\
\end{longtable}
\endgroup
Las 471 masas no son 471 especies elementales: incluyen estados excitados, resonancias, distintas cargas y múltiples realizaciones compuestas. Las 384 anchuras son observables temporales y no deben ser leídas por el carácter de masa.

\Needspace{7\baselineskip}
\subsection{Identidad interna \(\mathcal I_X\) de una realización material y criterio de individualización frente a observaciones múltiples}
Cada estado comparado debe poseer una identidad interna

\[
 \mathcal I_X=
 (\mathcal C_X,\mathcal R_X,\Gamma_X,\mathcal L_X,
 \sigma_X,\eta_X,P_X,\mathfrak M_X,\mathcal T_X),
\]
où :

\begin{itemize}
\item \(\mathcal C_X\) est le type physique de réalisation déterminé par l'objet, la représentation et le codomaine ;
\item \(\mathcal R_X\) est la représentation physique ;
\item \(\Gamma_X\) est un parcours, une orbite ou une composition orientée ;
\item \(\mathcal L_X\) est l'enregistrement généalogique ;
\item \(\sigma_X\) est la signature complète ;
\item \(\eta_X\) est la tour ;
\item \(P_X\) est le projecteur du couplage ;
\item \(\mathfrak M_X\) est le canal de mesure ;
\item \(\mathcal T_X\) est l'arbre de composition, le cas échéant.
\end{itemize}
Deux observations externes appartiennent au même état HMT lorsqu'elles partagent cette identité et ne diffèrent que par l'expérience, la carte ou l'estimateur. L'existence de deux mesures ne crée pas deux particules.

\Needspace{7\baselineskip}
\subsection{Morphisme de comparaison entre invariants internes et observables métrologiques}
La comparaison exige le diagramme

\[
 \mathcal L_M^{\rm HMT}
 \xrightarrow{\;\varsigma_{\mathcal S,\mu}\;}
 \mathbb R_{>0}u_M^{\mathcal S,\mu}
 \xleftarrow{\;\text{mesure}\;}
 \mathcal O_X^{\rm ext}.
\]
Une ligne scientifiquement complète contient :

\[
 (m_X^{\rm HMT},m_X^{\rm ext},
 u,\mathcal S,\mu,\Sigma_{\rm exp},\Sigma_{\rm HMT},
 \text{convention de pôle}).
\]
Pour les quarks, schéma et échelle sont obligatoires ; pour les résonances, la convention de pôle ; pour les limites, le niveau de confiance ; pour les observables corrélés, la matrice de covariance.

\Needspace{7\baselineskip}
\subsection{Évaluation explicite de \(17\) réalisations représentatives de l'atlas}
Les évaluations suivantes conservent la séparation entre la loi et la donnée de signature :

\Needspace{7\baselineskip}
\begingroup
\footnotesize
\setlength{\tabcolsep}{2.1pt}
\begin{longtable}{L{4.68cm}L{10.61cm}}
\toprule
\rowcolor{HMTNavy}\HMTtablehead{État} & \HMTtablehead{Valeur du caractère dans sa carte} \\
\midrule
\endfirsthead
\multicolumn{2}{l}{\sffamily\scriptsize\color{HMTGray}Suite du tableau}\\[-1pt]
\toprule
\rowcolor{HMTNavy}\HMTtablehead{État} & \HMTtablehead{Valeur du caractère dans sa carte} \\
\midrule
\endhead
\midrule
\endfoot
\bottomrule
\endlastfoot
\(\mu\) & \(105.658360\allowbreak{}294435\allowbreak{}829\ldots\ {\rm MeV}\) \\
\rowcolor{HMTLight}\(\tau\) & \(1776.860917\allowbreak{}631432\allowbreak{}3\ldots\ {\rm MeV}\) \\
\(u\) & \(2.165924\allowbreak{}536173\allowbreak{}907\ldots\ {\rm MeV}\) \\
\rowcolor{HMTLight}\(d\) & \(4.679550\allowbreak{}162948\allowbreak{}874\ldots\ {\rm MeV}\) \\
\(c\) & \(1.270500\allowbreak{}838641\allowbreak{}911\ldots\ {\rm GeV}\) \\
\rowcolor{HMTLight}\(s\) & \(92.940093\allowbreak{}907845\allowbreak{}64\ldots\ {\rm MeV}\) \\
\(t\) & \(172.597479\allowbreak{}544019\allowbreak{}1\ldots\ {\rm GeV}\) \\
\rowcolor{HMTLight}\(b\) & \(4.190119\allowbreak{}763299\allowbreak{}790\ldots\ {\rm GeV}\) \\
\(W\) & \(80.378964\allowbreak{}909704\allowbreak{}55\ldots\ {\rm GeV}\) \\
\rowcolor{HMTLight}\(Z\) & \(91.187533\allowbreak{}444313\allowbreak{}41\ldots\ {\rm GeV}\) \\
\(H\) & \(125.292466\allowbreak{}042536\allowbreak{}1\ldots\ {\rm GeV}\) \\
\rowcolor{HMTLight}\(\pi^0\) & \(134.976724\allowbreak{}327872\allowbreak{}1\ldots\ {\rm MeV}\) \\
\(\pi^\pm\) & \(139.570485\allowbreak{}171572\allowbreak{}2\ldots\ {\rm MeV}\) \\
\rowcolor{HMTLight}\(K^\pm\) & \(493.677085\allowbreak{}649530\allowbreak{}6\ldots\ {\rm MeV}\) \\
\(K^0\) & \(497.611186\allowbreak{}497750\allowbreak{}5\ldots\ {\rm MeV}\) \\
\rowcolor{HMTLight}\(p\) & \(938.271751\allowbreak{}546984\allowbreak{}9\ldots\ {\rm MeV}\) \\
\(n\) & \(939.565810\allowbreak{}472507\allowbreak{}0\ldots\ {\rm MeV}\) \\
\end{longtable}
\endgroup
L'électron bêta n'appartient pas à ce tableau, car son parcours central, ses coefficients de base et son lecteur complet sont déterminés par la construction elle-même.

\Needspace{7\baselineskip}

\endgroup
\begingroup
\let\chapter\section
\let\section\subsection
\let\subsection\subsubsection
\let\subsubsection\paragraph
\let\clearpage\relax
\subsection{Du catalogue des réalisations à l'inventaire de 471 masses}
L'extension universelle suit cinq opérations :

\[
\begin{aligned}
 \text{identité physique}
 &\longrightarrow
 \text{projecteur de représentation}
 \longrightarrow
 \text{registre généalogique élémentaire ou composé}\\
 &\longrightarrow
 \text{opérateur/pôle}
 \longrightarrow
 \text{carte de comparaison}.
\end{aligned}
\]
Pour les états élémentaires, la deuxième application agit sur une fibre de l'atlas. Pour les mésons et les baryons, elle agit sur les 432 et 1728 expressions composées. Pour les résonances, la quatrième application donne un pôle complexe. Pour les secteurs topologiques, elle donne d'abord un gap spectral et une représentation de tresse. Ce typage permet de couvrir tout l'inventaire sans déclarer que chaque nom externe correspond à une cellule élémentaire.

\Needspace{7\baselineskip}
\subsection{Des 384 largeurs à l'opérateur de survie}
Chaque largeur exige

\[
 X\longmapsto Q_XUQ_X
 \longmapsto \rho(Q_XUQ_X)
 \longmapsto\Gamma_X
\]
ou, de manière équivalente, le pôle de \(H_{\rm eff}\). Le lecteur doit incorporer canaux de désintégration, multiplicité, seuil et espace des phases. L'égalité

\[
 \Gamma_{\bar X}=\Gamma_X
\]
ne se déduit que si l'opérateur de désintégration est covariant par conjugaison.

\Needspace{7\baselineskip}
\section{Masse, température et statistique}
\Needspace{7\baselineskip}
\subsection{Transducteur thermique entre caractère de masse, température et occupation statistique}
La constante de Boltzmann est le transducteur

\[
 k_B:\mathcal L_\Theta\longrightarrow\mathcal L_E.
\]
L'horloge CT--108 fixe l'identité

\[
 k_B\Theta_{\rm clk}\log3
 =\hbar_{\rm int}\frac{2\pi}{108t_0}.
\]
Cette équation relie une unité d'information tritique, une fréquence et une énergie. Elle n'introduit pas la coordonnée SI de \(k_B\) dans la loi ; elle détermine le produit typé une fois la carte thermique choisie.

\Needspace{7\baselineskip}
\subsection{Projection fermionique de l'atlas et distribution de Fermi–Dirac}
Pour les niveaux

\[
 \epsilon=(0,3,6,9),\qquad
 g=(1,6,12,8),
\]
et contraintes

\[
 (N,E)=(12,66),
\]
la complétion extérieure produit

\[
 2\,104\,494
\]
microétats fermioniques compatibles. Les multiplicateurs effectifs sont

\[
 (\beta,\mu)
 =(0.1530701135,\;4.5028651976).
\]
Cette réalisation illustre comment masse, énergie et exclusion émergent des parcours fermioniques sans que la statistique soit un ajustement supplémentaire.

\Needspace{7\baselineskip}
\subsection{Projection bosonique de l'atlas et distribution de Bose–Einstein}
Sur les mêmes niveaux et sous les mêmes contraintes, la complétion symétrique produit

\[
 259\,436\,430
\]
microétats, avec

\[
 (\beta,\mu)
 =(0.0545672796,\;-15.9256978191).
\]
La différence entre Fermi--Dirac et Bose--Einstein est donc une différence dans la règle de composition des parcours, et non dans la valeur d'une masse individuelle.

\Needspace{7\baselineskip}
\subsection{Critère de condensation des sections bosoniques et occupation macroscopique}
Un condensat bosonique apparaît lorsqu'une composante symétrique acquiert une occupation macroscopique dans la valeur propre minimale de l'opérateur. Une mer fermionique, en revanche, remplit un ensemble de valeurs propres distinctes jusqu'à une frontière. HMT décrit les deux au moyen du même spectre de masse/énergie et de deux foncteurs de Fock distincts. Cette séparation empêche d'identifier « beaucoup de particules sur un parcours » à une particule plus massive.

\Needspace{7\baselineskip}
\section{Propagation différentielle des incertitudes et stabilité structurelle du caractère de masse}
\Needspace{7\baselineskip}
\subsection{Invariance du graphe génératif sous la confrontation métrologique ultérieure}
Le résultat interne est figé dans l'ordre

\[
 \mathcal R_{324}
 \to P_X
 \to\mathcal L_X
 \to\sigma_X
 \to\eta_X
 \to\widehat M_X
 \to\mu_X.
\]
Ce n'est qu'ensuite que la carte externe est ouverte. Une permutation arbitraire des 471 masses ou des 384 largeurs doit laisser invariants tous les objets précédents.

\Needspace{7\baselineskip}
\subsection{Décomposition de l'incertitude structurelle en contributions expérimentale, de carte, de parcours, de tour et de couplage}
L'incertitude totale se décompose comme

\[
 \Sigma_{\rm tot}
 =\Sigma_{\rm ext}
 +\Sigma_{\rm carta}
 +\Sigma_{\rm route}
 +\Sigma_{\rm tower}
 +\Sigma_{\rm coupling}.
\]
Ici :

\begin{itemize}
\item \(\Sigma_{\rm ext}\) est la covariance expérimentale ;
\item \(\Sigma_{\rm carta}\) quantifie schéma, échelle et conversion des unités ;
\item \(\Sigma_{\rm route}\) compare tous les parcours de l'orbite admissible ;
\item \(\Sigma_{\rm tower}\) utilise la borne de reste de la série harmonique ;
\item \(\Sigma_{\rm coupling}\) mesure la variation entre les entrelaceurs autorisés.
\end{itemize}
Elles ne s'additionnent pas nécessairement comme des nombres indépendants ; l'équation exprime une décomposition des sources qui doit être convertie en covariance conjointe.

\Needspace{7\baselineskip}
\subsection{Dérivée de l'observable de masse par rapport à la signature et aux lecteurs}
Pour

\[
 \log m_X=
 \log m_\star+
 \log\chi_0(\sigma_X)
 +\kappa_X\log R_{\rm act}
 +\sum_h\eta_{X,h}s_h,
\]
la variation est

\[
 d\log m_X
 =d\log m_\star
 +\langle\Theta_0,d\sigma_X\rangle
 +\langle\sigma_X,d\Theta_0\rangle
 +\log R_{\rm act}\,d\kappa_X
 +\kappa_X\,d\log R_{\rm act}
 +\sum_h(s_h\,d\eta_{X,h}+\eta_{X,h}\,ds_h).
\]
Cette formule permet d'attribuer chaque décimale au parcours, à une constante, au lecteur, à la tour ou à la carte. Un résultat est robuste lorsque les perturbations préservant les prémisses maintiennent sa classe spectrale et son incertitude calculée.

\Needspace{7\baselineskip}
\subsection{Essais de dépendance structurelle par suppression contrôlée du feuillet, de l'orientation, de la retenue, de la mémoire, de l'enroulement et de la frontière}
Un test de dépendance élimine séparément feuillet, orientation, retenue, mémoire, enroulement ou frontière, puis recalcule l'objet. Le résultat attendu n'est pas que tout change, mais que changent exactement les coordonnées dépendant de la composante éliminée. Si la suppression de la mémoire laisse invariant un coefficient défini par le retour, la construction de ce coefficient est réfutée.

\Needspace{7\baselineskip}
\subsection{Obstructions à l'identification des parcours, des caractères et des grandeurs métrologiques}
La loi est réfutée dans une application concrète si l'un des faits suivants se produit :

\begin{enumerate}
\item une masse externe modifie le parcours ou le projecteur choisi ;
\item la signature présentée n'est pas régénérée à partir de l'enregistrement généalogique ;
\item le raffinement ne converge pas ou dépend d'une troncature choisie selon le résidu ;
\item deux parcours équivalents reçoivent des masses distinctes sans brisure de symétrie ;
\item une antiparticule acquiert une masse distincte sans violation déclarée de la conjugaison ;
\item un triplet de couleur se scinde sans brisure de couleur ;
\item un scalaire est présenté alors que la variance spectrale est positive ;
\item une largeur est obtenue à partir du caractère de masse sans opérateur temporel ;
\item une carte de quark omet le schéma ou l'échelle ;
\item une instabilité est interprétée comme un tachyon persistant.
\end{enumerate}
\Needspace{7\baselineskip}

\endgroup

\section{Leptons, neutrinos, quarks et secteurs de jauge}
Les leptons chargés, les six saveurs de quarks et les fibres neutriniques sont réalisés comme multisections typées. La couleur est l’orbite résiduelle ternaire \(12+12+12\) ; saveur, génération, schéma et échelle sont des données distinctes. Les neutrinos appartiennent à la complétion fermionique ; la condition de Majorana exige en outre l’autoconjugaison et un terme de masse compatible. Photon et gluon appartiennent à la carte nulle exacte.

\begingroup
\let\chapter\section
\let\section\subsection
\let\subsection\subsubsection
\let\subsubsection\paragraph
\let\clearpage\relax
% Vista científica exacta materializada desde una rama condicional.
% Fuente: source/demostraciones_primarias/09_06h_atlas_familias_rev6.tex; rama: HMTIncludeAtlasPhysical.
% No modifica el contenido matemático de la rama de origen.
\chapter[Correspondance avec les secteurs physiques]{Correspondance entre classes internes et secteurs physiques}
\label{sec:censo-fisico-atlas-rev7}

La réalisation physique conserve les quotients précédents et ajoute les données
dont une représentation a besoin sans en faire de nouvelles cellules. Pour une
section persistante \(X\), nous écrivons de manière typée
\[
 \mathcal P_X=
 \bigl[q(X),\,[X]_{\mathcal R},\,\operatorname{Rep}(X),\,
       \gamma_X,\,\tau_{12}(X),\,\sigma_X\bigr].
\]
Ici, \(q(X)\) fixe la multisection, \([X]_{\mathcal R}\) la famille de chemins,
\(\operatorname{Rep}(X)\) la représentation, \(\gamma_X\) l'histoire
persistante, \(\tau_{12}(X)\) son mot dodécaphasique et \(\sigma_X\) la signature
entière extraite après le registre. Les niveaux \(G0,\ldots,G4\) restent
des profondeurs de l'atlas ; ils ne sont pas renommés comme les trois générations
physiques.

L'information de saveur est conservée au moyen de trois coordonnées distinctes :
\[
 \operatorname{flav}(X)=
 \bigl(\sigma_{ud}(X),g(X),\chi_{\rm fina}(X)\bigr).
\]
La première sépare les branches leptonique, neutrino, de type haut et de type bas ;
la deuxième distingue la génération physique ; la troisième conserve
l'orientation et la mémoire qui séparent les sections au sein d'une même branche. Pour
les quarks, l'orbite résiduelle \(R/G/B\) est ajoutée après la fixation de la saveur et
avant l'évaluation du caractère. Couleur, saveur et masse sont donc des lectures de
types différents.

% El inventario humano sectorial se conserva en su totalidad,
% pero se monta después del cierre electrónico, en la residencia sectorial.

Ce recensement n'épuise pas les réalisations de la loi. Les mésons et les baryons se
construisent en composant des registres enrichis avant la projection de la signature ;
les résonances et les particules virtuelles se distinguent par leur horizon de
persistance. Les secteurs Fibonacci, \(Z_6\) et Majorana--Ising conservent
des codomaines topologiques propres et ne sont pas forcés dans une table ordinaire
de masses.

\section{Signatures de trajectoire et couples dodécaphasiques coindexés par saveur}

Pour l'ensemble des saveurs
\(\mathcal Q=\{u,d,c,s,t,b\}\), le tableau définit une correspondance finie
entre une empreinte géométrique et un couple dodécaphasique. Dans chaque ligne, on a
\[
 W_\pm=6n_A\pm s_C.
\]
Ce tableau n'est pas une comparaison a posteriori, mais le graphe d'une application
finie sur le domaine nominal certifié. Si
\(\mathscr R_{\mathcal Q}^{\rm nom}=\{\rho_u,\rho_d,\rho_c,\rho_s,\rho_t,\rho_b\}\)
est l'ensemble des empreintes CT--108 qui suivent, nous définissons
\[
 H_{\rm flav}:\mathscr R_{\mathcal Q}^{\rm nom}\longrightarrow\mathbb Z^2,
 \qquad H_{\rm flav}(\rho_q)=(n_A(q),s_C(q)).
\]
Les six images sont uniques, le tableau épuise le domaine et aucune colonne
ne contient de masse. Ainsi, \(H_{\rm flav}\) est un extracteur causal clos sur
le domaine nominal : il lit d'abord empreinte, orientation et mémoire, et ce n'est qu'ensuite que
le caractère évalue la masse. L'orbite ternaire de couleur est quotientée sans changer
\(H_{\rm flav}\) ; le renversement conserve la classe de masse et transporte la
représentation particule--antiparticule. Le tableau est donc la définition
exécutable du morphisme et non une obligation future :

\begin{longtable}{c c c c c c c}
\toprule
Saveur & Génération & Giros/cambios & Eje & \(N,E,S,O\)
& \((n_A,s_C)\) & \((W_+,W_-)\)\\
\midrule
\endfirsthead
\toprule
Saveur & Génération & Giros/cambios & Eje & \(N,E,S,O\)
& \((n_A,s_C)\) & \((W_+,W_-)\)\\
\midrule
\endhead
\(u\)&1&21/30&N--S&34,31,31,12&(11,7)&(73,59)\\
\(d\)&1&20/32&E--O&27,35,18,28&(17,8)&(110,94)\\
\(c\)&2&21/29&E--O&20,38,24,26&(61,8)&(374,358)\\
\(s\)&2&17/27&N--S&50,32,12,14&(41,-3)&(243,249)\\
\(t\)&3&21/26&N--S&46,27,16,19&(100,-1)&(599,601)\\
\(b\)&3&21/30&E--O&12,63,9,24&(71,-5)&(421,431)\\
\bottomrule
\end{longtable}

La coïncidence de deux rotations ou d'une hauteur spectrale n'identifie pas les
espèces : le chemin cardinal, l'orientation, le registre dodécaphasique et la
représentation restent dans l'objet. La chaîne
cellule--chemin--registre chronologique enrichi--signature--caractère est définie
indépendamment de \(H_{\rm flav}\). Les
quatre orientations initiales de chaque cellule produisent
\(81\cdot4=324\) généalogies exécutables ; ce cardinal compte des histoires, non des
particules, des familles ni des masses.

\endgroup
\begingroup
\let\chapter\section
\let\section\subsection
\let\subsection\subsubsection
\let\subsubsection\paragraph
\let\clearpage\relax
\section{Stratification générationnelle, supports de quarks et réalisation neutrino}
\label{sec:generaciones-quarks-neutrinos-rev2}
\index{générations fermioniques!stratification HMT--MD}
\index{quark!supports cellulaires et registre}
\index{PMNS!transport des bases}

La loi générale agit sur des multisections et non sur une liste de masses.
L'apparition de trois générations s'exprime donc comme compatibilité de
trois structures : la stratification de l'atlas, les registres de saveur et
l'espace tridimensionnel de mélange. Dans le secteur leptonique chargé, les
multisections
\begin{equation}
 G_0^{\ell},\qquad G_2^{\ell},\qquad G_4^{\ell}
 \label{eq:tres-multisecciones-leptonicas-rev2}
\end{equation}
ont pour rangs \(1,8,16\). Le centre \(G_0^\ell\) contient la section
électronique de rang un ; \(G_2^\ell\) et \(G_4^\ell\) conservent, en revanche,
des spectres opératoriels non triviaux. L'organisation générationnelle complète
est
\begin{center}
\small
\begin{tabularx}{0.96\textwidth}{@{}c c c c X@{}}
\toprule
\textbf{Génération} & \textbf{Lepton} & \textbf{Par quark} &
\textbf{Sabor neutro} & \textbf{Soporte estructural}\\
\midrule
1 & \(e\) & \((u,d)\) & \(\nu_e\) & centre \(G_0^\ell\), registre de
première génération et premier vecteur de la base d'interaction\\
2 & \(\mu\) & \((c,s)\) & \(\nu_\mu\) & multisection \(G_2^\ell\),
registre de deuxième génération et deuxième vecteur de la base d'interaction\\
3 & \(\tau\) & \((t,b)\) & \(\nu_\tau\) & multisection \(G_4^\ell\),
registre de troisième génération et troisième vecteur de la base d'interaction\\
\bottomrule
\end{tabularx}
\end{center}
La conjugaison particule--antiparticule inverse les coordonnées de charge et
d'orientation de la représentation, mais conserve cette stratification ; par
conséquent, elle ne double pas le nombre de générations.

La dimension trois dispose en outre d'un porteur d'incidence propre. La
représentation orientée dodécaphasique contient
\begin{equation}
 V_{12}\cong V_1\oplus V_3\oplus V_{3'}\oplus V_5,
 \qquad \operatorname{rank}P_3=3,
 \label{eq:descomposicion-rango-tres-generaciones-rev2}
\end{equation}
et l'entrelaceur dodécaphasique--Witt transporte le secteur positif selon
\begin{equation}
 P_G\xleftrightarrow{\;W_{GK}\;}P_3
 \xleftrightarrow{\;U_W\;}Q_3,
 \qquad
 \operatorname{rank}P_G=\operatorname{rank}P_3
 =\operatorname{rank}Q_3=3.
 \label{eq:transporte-rango-tres-generaciones-rev2}
\end{equation}
Ce porteur tridimensionnel permet de comparer les bases générationnelles. Les
chemins, les signatures et les opérateurs de masse fournissent ensuite le contenu
spectral que CKM et PMNS transportent.

\subsection{Position interne et registres des \(6\) saveurs de quark}

La position d'une saveur de quark est la donnée composée
\begin{equation}
 \mathfrak q_X=
 (\text{branche de charge},\text{génération},(\mathbb Z/3\mathbb Z)_{\rm color},
  \text{multisection},\text{registre de route}).
 \label{eq:posicion-quark-compuesta-rev2}
\end{equation}
Les branches de charge \(-1/3\) possèdent les supports
\begin{align}
 D_1&=\{(4,5),(6,5)\},\nonumber\\
 D_2&=\{(4,3),(4,7),(6,3),(6,7)\},\nonumber\\
 D_3&=\{(2,3),(2,5),(2,7),(8,3),(8,5),(8,7)\},\nonumber\\
 D_4&=\{(2,1),(2,9),(4,1),(4,9),(6,1),(6,9),(8,1),(8,9)\},
 \label{eq:soportes-quark-down-rev2}
\end{align}
tandis que les branches de charge \(+2/3\) occupent
\begin{align}
 U_1&=\{(4,4),(4,6),(6,4),(6,6)\},\nonumber\\
 U_3&=\{(2,2),(2,4),(2,6),(2,8),(4,2),(4,8),
       (6,2),(6,8),(8,2),(8,4),(8,6),(8,8)\}.
 \label{eq:soportes-quark-up-rev2}
\end{align}
Chaque support se relève à quatre orientations par cellule. Le quotient de couleur
organise \(36=12_{\rm R}+12_{\rm G}+12_{\rm B}\) biographies orientées et conserve une même
masse par saveur avant la réalisation de couleur.

Le transducteur nominal
\begin{equation}
 \mathcal N_{\rm ruta}(\Gamma_f)
 =((N,E,S,O),h_{369},r_5)
 \in\mathbb Z_{\ge0}^{4}\times\mathbb Z_{\ge0}^{2}
 \label{eq:transductor-nominal-quark-rev2}
\end{equation}
publie les coordonnées suivantes avant toute réalisation scalaire :
\begin{center}
\small
\begin{tabularx}{0.96\textwidth}{@{}c c c c c X@{}}
\toprule
\textbf{Saveur} & \textbf{Génération} & \textbf{Branche} &
\textbf{\((N,E,S,O)\)} & \textbf{\((h_{369},r_5)\)} &
\textbf{Eje dominante}\\
\midrule
\(u\)&1&\(+2/3\)&\((34,31,31,12)\)&\((51,13)\)&N--S\\
\(d\)&1&\(-1/3\)&\((27,35,18,28)\)&\((47,17)\)&E--O\\
\(c\)&2&\(+2/3\)&\((20,38,24,26)\)&\((49,16)\)&E--O\\
\(s\)&2&\(-1/3\)&\((50,32,12,14)\)&\((57,11)\)&N--S\\
\(t\)&3&\(+2/3\)&\((46,27,16,19)\)&\((57,9)\)&N--S\\
\(b\)&3&\(-1/3\)&\((12,63,9,24)\)&\((47,7)\)&E--O\\
\bottomrule
\end{tabularx}
\end{center}
Les signatures nominales associées sont
\begin{center}
\small
\begin{tabular}{@{}c c c@{}}
\toprule
\textbf{Saveur} & \textbf{\((n_A,s_C)\)} & \textbf{\((W_+,W_-)\)}\\
\midrule
\(u\)&\((11,7)\)&\((73,59)\)\\
\(d\)&\((17,8)\)&\((110,94)\)\\
\(c\)&\((61,8)\)&\((374,358)\)\\
\(s\)&\((41,-3)\)&\((243,249)\)\\
\(t\)&\((100,-1)\)&\((599,601)\)\\
\(b\)&\((71,-5)\)&\((421,431)\)\\
\bottomrule
\end{tabular}
\end{center}
Ces coordonnées distinguent des saveurs qui partagent une branche ou une génération et closent
le morphisme de saveur sans consulter aucune masse. Pour une section physique (P),
nous définissons
\begin{equation}
 \begin{aligned}
 \operatorname{sabor}(P)
 &=\bigl(\sigma_{\rm ud}(P),g(P),\chi_{\rm fina}(P)\bigr),\\
 \chi_{\rm fina}(P)
 &=(\text{rotations},\text{commutations},\text{axe},N,E,S,O,h_{369},r_5),
 \end{aligned}
 \label{eq:morfismo-sabor-ruta-cerrado-rev2}
\end{equation}
et, pour les quarks, on conserve en outre l'orbite ternaire
\([\gamma_P]_{\rm col}=\{\gamma_P^R,\gamma_P^G,\gamma_P^B\}\). La réalisation
physique est l'application déjà déterminée par les données précédentes,
\begin{equation}
 \mathcal S_{\rm phys}(P,\mathfrak M):
 \bigl(\sigma_{\rm ud},g,\chi_{\rm fina},[\gamma_P]_{\rm col}\bigr)
 \longmapsto \gamma^{\rm surv}_{P,\mathfrak M}
 \longmapsto
 \bigl(\sigma(\gamma^{\rm surv}_{P,\mathfrak M}),
       B_F^D,B_F^\Omega\bigr),
 \label{eq:realizacion-fisica-sabor-ruta-cerrada-rev2}
\end{equation}
où \(\mathfrak M\) est le couplage compatible et
\(\gamma^{\rm surv}_{P,\mathfrak M}\) le chemin qui persiste sous celui-ci. Pour
les leptons chargés, l'orbite de couleur se réduit à l'objet terminal. Le caractère
de masse s'évalue alors sur cette même généalogie :
\begin{equation}
 \boxed{
 \frac{m_{\mathfrak M}(P)}{m_e^{\rm HMT}}
 =\chi_{\rm mass}\!\left(
   \sigma\!\left(\gamma^{\rm surv}_{P,\mathfrak M}\right)\right).}
 \label{eq:masa-como-caracter-persistencia-rev2}
\end{equation}
Par conséquent, saveur, génération et couleur n'attendent pas un sélecteur futur : elles publient la
classe de chemin et son vecteur propre avant d'ouvrir la comparaison externe. La masse ne
choisit pas rétrospectivement le chemin ; elle ne fait que lire multiplicativement la persistance
que le couplage a stabilisée.

\subsection{Fermionicité, valeurs propres de masse et transport PMNS}

La clôture extérieure du foncteur d'échange construit les relations CAR
et fixe la fermionicité. Dans le secteur neutre, l'étiquette
\(\nu_e,\nu_\mu,\nu_\tau\) désigne une base d'interaction ; les valeurs de
masse sont des valeurs propres d'un opérateur neutre mélangé. Les référentiels internes
\(F_\ell,F_\nu:\mathbb C^3\to\operatorname{im}P_3\) comparent ces bases, et
la matrice
\begin{equation}
 U_{\rm HMT}=F_\ell^*\Phi_LF_\nu
 \label{eq:transporte-pmns-generaciones-rev2}
\end{equation}
est un transport unitaire. PMNS change de base entre les états d'interaction
et les vecteurs propres de masse ; l'échelle neutre appartient au spectre de l'opérateur,
non au changement de base.

\subsection{Extracteur neutre Z0 antérieur à la projection}
\label{sec:extractor-neutro-z0-torsion-rev2}

La construction récupérée du secteur neutre agit sur le registre
dodécaphasique enrichi, avant la contraction du feuillet, de la torsion et du bit
bidirectionnel. Soient \(E_m\), \(m\in\mathbb Z/12\mathbb Z\), les douze
événements de neuf ticks du registre CT--108. À chaque tick \(t\), on conserve
\[
 \theta(t)=\mathbf1_{\{d(t)=9\}},\qquad
 \varepsilon(t)\in\{-1,+1\},\qquad
 \kappa(t)=\mathbf1_{\{9\ {\rm en\ corona}\}},
 \qquad \omega(t)=(-1)^{\kappa(t)}.
\]
Le premier facteur détecte la torsion nonadique, le deuxième conserve le feuillet
two--ways et le troisième change le signe lors de la rotation de \(180^\circ\) de la
couronne. L'entier torsionnel de chaque événement et sa différence antipodale sont
\begin{equation}
 T_\nu(m)=\sum_{t\in E_m}
   \omega(t)\varepsilon(t)\theta(t),\qquad
 e_m^\nu=T_\nu(m)-T_\nu(m+6),
 \label{eq:extractor-z0-torsion-antipodal-rev2}
\end{equation}
et le trit neutre est fixé, sans paramètres continus, par
\begin{equation}
 \tau_m^\nu=\operatorname{sgn}(e_m^\nu)
 \in\{-1,0,+1\}.
 \label{eq:trit-neutro-torsion-antipodal-rev2}
\end{equation}

Pour les quatre orbites de saut quatre
\[
 I_a=\{a,a+4,a+8\}\pmod{12},\qquad a=0,1,2,3,
\]
le coefficient d'événement qui entre dans la signature et l'observable de majorité
des signes sont, respectivement,
\begin{equation}
 b_{120}^\nu
 =\sum_{a=0}^{3}\operatorname{sgn}
   \left(\sum_{m\in I_a}e_m^\nu\right),
 \qquad
 \nu_{120}^{\nu,\mathrm{sgn}}
 =\sum_{a=0}^{3}\operatorname{sgn}
   \left(\sum_{m\in I_a}\tau_m^\nu\right).
 \label{eq:canales-120-neutros-distintos-rev2}
\end{equation}
Ces deux entiers ne sont pas identifiés : le premier est la coordonnée
\(b_{120}=\nu_{120}^{\rm ev}\) du caractère, tandis que le second est un
diagnostic de majorité. La mémoire quadratique est
\begin{equation}
 b_\zeta^\nu=\nu_{270}^\nu
 =\sum_{m=0}^{11}\tau_m^\nu\tau_{m+3}^\nu.
 \label{eq:canal-270-neutro-rev2}
\end{equation}
Par conséquent, pour chaque histoire neutre survivante
\(\gamma_{\nu_i}\), le registre publie avant toute comparaison la signature
\[
 v_\nu(\gamma_{\nu_i})
 =\bigl(0,0,k_{\Delta_4}^{(i)},b_{120}^{\nu_i},
        b_\zeta^{\nu_i}\bigr)
\]
et l'exposant relatif
\begin{equation}
 \mathcal A_\nu^{(i)}
 =k_{\Delta_4}^{(i)}\Delta_4
  +b_{120}^{\nu_i}s_{120}
  +b_\zeta^{\nu_i}\zeta_{270},
 \qquad
 \zeta_{270}=135\Delta_4^2.
 \label{eq:caracter-relativo-neutro-z0-rev2}
\end{equation}
Ici, \(\zeta_{270}\) n'est pas remplacé par le moment de tour
\(s_{270}=q_-^{270}-q_+^{270}\).

La conjugaison globale two--ways
\(\tau^\nu\mapsto-\tau^\nu\) inverse
\(b_{120}^\nu\) et conserve \(b_\zeta^\nu\). Elle engendre donc deux
branches de spectre relatif,
\begin{equation}
 \mathcal A_{\nu,\pm}^{(i)}
 =k_{\Delta_4}^{(i)}\Delta_4
  \pm b_{120}^{\nu_i}s_{120}
  +b_\zeta^{\nu_i}\zeta_{270},
 \qquad
 \text{NON}\longleftrightarrow\text{IO},
 \label{eq:bifurcacion-neutra-no-io-rev2}
\end{equation}
qui se distinguent par le signe ultérieur de
\(\Delta m_{31}^2\), non par une sélection avec des données PDG. Cette construction
clôt l'extracteur discret Z0 et la bifurcation relative. L'échelle absolue
commune \(m_0\) et le noyau physique PMNS de rang complet restent des objets
séparés : ils ne sont pas introduits dans le sélecteur et ne sont pas déclarés
déduits de la seule parité two--ways.

La réalisation avec une unique phase de rang un produit
\[
 (\theta_{12},\theta_{23},\theta_{13})
 =(19.6807^\circ,56.4970^\circ,29.6224^\circ)
\]
et ne reproduit pas la phénoménologie du mélange. Ce contrôle sélectionne comme domaine approprié un
noyau de registre de rang complet,
\begin{equation}
 G_{ab}=\frac1{\sqrt{|L_a||N_b|}}
 \sum_{c\in L_a}\sum_{d\in N_b}
 K_{\rm registro}(c,d;\xi_c,\xi_d),
 \qquad U_{\rm int}=\operatorname{polar}(G),
 \label{eq:kernel-registro-rango-completo-rev2}
\end{equation}
dont la source technique est développée dans le chapitre du transport CKM et
PMNS. Ainsi sont séparés l'opérateur de masse, ses valeurs propres et l'application de
lecture par saveur.

\endgroup
\begingroup
\let\chapter\section
\let\section\subsection
\let\subsection\subsubsection
\let\subsubsection\paragraph
\let\clearpage\relax
\section{Transport des opérateurs de masse par CKM}
\label{sec:transporte-ckm-operadores-rev11}

L'incidence N42 fixe le sextuplet \((4,5,7,6,8,3)\) et la matrice rationnelle
\[
 M_{\rm CKM}=
 \begin{pmatrix}
 2&-5&28\\
 0&7&-25/2\\
 1&-20&-2/3
 \end{pmatrix},
 \qquad \det M_{\rm CKM}=-\frac{3857}{6}.
\]
Les coordonnées angulaires sont construites par
\[
 \begin{pmatrix}\theta_{12}\\\theta_{23}\\\theta_{13}\end{pmatrix}
 =M_{\rm CKM}
 \begin{pmatrix}A\\ C^*/6\\ (180/\pi)\Delta_4\end{pmatrix},
 \qquad \delta_{\rm CP}=9A.
\]
Cette transformation appartient à la clôture mathématique de HMT. Dans sa réalisation physique, soient \(B_u\) et \(B_d\) les opérateurs bidirectionnels déjà construits dans les fibres de quarks de types up et down. Le changement de base agit selon
\begin{equation}
 B_d^{(u)}=V_{\rm CKM}B_dV_{\rm CKM}^*,
 \qquad
 R_{\rm act}^{B_d^{(u)}}
 =V_{\rm CKM}R_{\rm act}^{B_d}V_{\rm CKM}^*.
 \label{eq:transporte-ckm-kappa-rev11}
\end{equation}
CKM transporte les valeurs propres entre bases et conserve le spectre et le déterminant. Le contrôle de la violation CP prend la forme de l'identité
\[
 \left|\det[B_u,V_{\rm CKM}B_dV_{\rm CKM}^*]\right|
 =2|J|\prod_{i<j}|\kappa_i^u-\kappa_j^u|
       \prod_{a<b}|\kappa_a^d-\kappa_b^d|.
\]
La matrice de mélange et le lecteur de route remplissent donc des fonctions complémentaires : le lecteur détermine l'opérateur spectral et CKM détermine sa représentation dans la base de saveur.

\section{Repères géométriques de rang \(3\) pour le secteur leptonique}
\label{sec:frames-pmns-rev11}

La construction leptonique utilise l'involution diagonale de signe
\[
 S_{\rm sign}=\operatorname{diag}(1,1,1,-1,-1,-1,1,1,1,-1,-1,-1),
 \qquad P_-=(I-S_{\rm sign})/2,
\]
et le projecteur icosaédrique canonique \(P_3\) de rang trois. La compression
\[
 C_-=P_-P_3P_-\big|_{\operatorname{im}P_-}
\]
possède le spectre
\[
 \left\{\frac{5+\sqrt5}{10},\frac12,
 \frac{5-\sqrt5}{10},0,0,0\right\}.
\]
Les trois valeurs propres positives sont simples. Leurs projecteurs spectraux, appliqués à l'indicatrice de l'hexade négative orientée et normalisés après projection par \(P_3\), construisent un repère orthonormal chargé \(F_\ell\) qui satisfait
\[
 F_\ell^*F_\ell=I_3,
 \qquad F_\ell F_\ell^*=P_3.
\]

Pour le secteur des neutrinos, soient \(H_1,H_2,H_3\) les trois hexades positives orientées et
\[
 X=(P_3\mathbf1_{H_1},P_3\mathbf1_{H_2},P_3\mathbf1_{H_3}).
\]
Sa matrice de Gram est définie positive et satisfait
\[
 \det(X^*X)=\frac{5-\sqrt5}{40}>0.
\]
Par conséquent,
\[
 F_\nu=X(X^*X)^{-1/2},
 \qquad F_\nu^*F_\nu=I_3,
 \qquad F_\nu F_\nu^*=P_3.
\]
Le transport réel
\begin{equation}
 O=F_\ell^*F_\nu\in SO(3)
 \label{eq:transporte-frames-pmns-rev11}
\end{equation}
s'obtient sans masses ni angles PMNS.

\HMTClaim{MD-R-PMNS-RANKONE-UNITARY}
\begin{theorem}[Unitarité de la phase dodécaphasique minimale]
\label{thm:unitaria-interna-pmns-rev11} Soit \(b_K\in\operatorname{im}P_3\) le vecteur normalisé du témoin de Witt, \(Q_K=|b_K\rangle\langle b_K|\), et
\[
 \Phi_L=I+(\zeta_{12}-1)Q_K,
 \qquad \zeta_{12}=e^{i\pi/6}.
\]
Alors
\begin{equation}
 U_{\rm HMT}=F_\ell^*\Phi_LF_\nu
 \label{eq:unitaria-interna-pmns-rev11}
\end{equation}
est une matrice unitaire indépendante de valeurs métrologiques cibles. L'inversion d'orientation conjugue \(U_{\rm HMT}\).
\end{theorem}

\begin{proof}
\(Q_K\) est un projecteur orthogonal, donc \(\Phi_L\) est unitaire. Les deux repères sont des isométries de même image \(\operatorname{im}P_3\). Par conséquent
\[
 U_{\rm HMT}^*U_{\rm HMT}
 =F_\nu^*\Phi_L^*P_3\Phi_LF_\nu=I_3.
\]
Le changement \(\pi/6\mapsto-\pi/6\) produit \(\Phi_L^*\), et toutes les autres primitives sont réelles.
\end{proof}

\subsection{Contrôle négatif de rang un et transport PMNS de rang complet}
\label{sec:no-go-pmns-rango-uno-rev11} \HMTClaim{MD-R-PMNS-RANKONE-PHYSICAL-NOGO}

Une fois la construction figée, la comparaison dans l'ordre déclaré produit
\[
 (\theta_{12},\theta_{23},\theta_{13})
 =(19.6807^\circ,56.4970^\circ,29.6224^\circ),
 \qquad J=-0.0426621.
\]
La meilleure des 36 permutations discrètes laisse \(\theta_{13}\simeq17.02^\circ\). La phase de rang un engendre une matrice unitaire interne, mais conserve un stabilisateur \(U(2)\) et est exclue comme réalisation physique de PMNS.

L'exclusion du porteur de rang un ne laisse pas la réalisation PMNS indéterminée. Soient \(M_\ell\) et \(M_\nu\) les opérateurs chargé et neutre produits par la Loi générale des masses, et soient \(F_\ell,F_\nu\) leurs repères spectraux orthonormaux complets sur l'image commune \(\operatorname{im}P_3\). Le transport leptonique de rang complet est alors construit par
\[
 U_{\rm PMNS}:=F_\ell^*F_\nu,
 \qquad
 U_{\rm PMNS}^*U_{\rm PMNS}=I_3.
\]
Pour des spectres simples, cette sortie est déterminée à des réajustements diagonaux de phase et à l'ordre spectral déclaré près ; en présence de dégénérescences, la sortie exacte est la classe double de repères correspondante. Le noyau du registre constitue un certificat fini supplémentaire de cette réalisation, et non son générateur ou une condition encore requise pour son existence. Pour les fibres chargées \(L_a\) et neutriniques \(N_b\), sa forme agrégée est
\[
 G_{ab}=\frac{1}{\sqrt{|L_a||N_b|}}
 \sum_{c\in L_a}\sum_{d\in N_b}K_{\rm reg}(c,d),
 \qquad U_{\rm int}=\operatorname{polar}(G),
\]
avec naturalité sous l'orientation. Lorsque les entrées explicites du registre vérifient \(\det G\ne0\), son facteur polaire fournit un certificat alternatif du transport déjà construit. Le noyau lit la retenue, le feuillet, l'orientation, la vacance et l'holonomie ; ses paramètres sont figés avant l'examen des angles externes.

\paragraph{Domaine et codomaine des réalisations CKM et PMNS}
La matrice rationnelle CKM, les repères \(F_\ell,F_\nu\), le transport \(O\), la phase \(\Phi_L\), l'unitarité de \(U_{\rm HMT}\) et le transport complet \(U_{\rm PMNS}=F_\ell^*F_\nu\) sont des résultats internes exacts sous les primitives déclarées. CKM constitue un transport physique une fois fixées les bases de saveur et de masse. La comparaison réfute uniquement la phase PMNS de rang un ; elle ne diminue pas la portée de la construction spectrale complète. La naturalité et le critère \(\det G\ne0\) appartiennent au certificat alternatif du registre et restent distincts de l'existence et de l'unitarité du transport PMNS.

\endgroup
\begingroup
\let\chapter\section
\let\section\subsection
\let\subsection\subsubsection
\let\subsubsection\paragraph
\let\clearpage\relax
\section{Involution des feuillets dans \(4\) représentations}

L'involution \(C^*\mapsto-C^*\) organise quatre réponses distinctes :
\begin{equation}
\label{rm:eq:modular-sheet-parities}
\begin{aligned}
S_{\rm bi}(A,-C^*)&=S_{\rm bi}(A,C^*),\\
\Gamma_{6\to8}(A,-C^*)&=-\Gamma_{6\to8}(A,C^*),\\
(\widehat\mu,\widehat\varepsilon)&\longleftrightarrow
(\widehat\varepsilon,\widehat\mu),\\
\widehat c&\longmapsto\widehat c,
\qquad D_A\longmapsto D_A.
\end{aligned}
\end{equation}
L'information bicouche est paire ; l'orientation aréale est impaire ; les deux
constitutives forment un doublet, et les déterminants restent fixes. Ce tableau
des parités précède toute interprétation comme symétrie physique.

\section{Transformation rationnelle involutive dans l'espace des paramètres CKM}

Soit
\begin{equation}
\label{rm:eq:modular-Mckm}
M_{\rm CKM}=
\begin{pmatrix}
2&-5&28\\
0&7&-25/2\\
1&-20&-2/3
\end{pmatrix},
\qquad
\det M_{\rm CKM}=-\frac{3857}{6}.
\end{equation}
L'inverse existe sur \(\mathbb Q\) et est
\begin{equation}
\label{rm:eq:modular-Mckm-inverse}
M_{\rm CKM}^{-1}=
\begin{pmatrix}
1528/3857&3380/3857&801/3857\\
75/3857&176/3857&-150/3857\\
6/551&-30/551&-12/551
\end{pmatrix}.
\end{equation}
Avec
\begin{equation}
\label{rm:eq:modular-ckm-coordinates}
h=\frac{C^*}{6},
\qquad
\Delta=\frac{180}{\pi}\Delta_4,
\qquad
\boldsymbol\theta=M_{\rm CKM}(A,h,\Delta)^T,
\end{equation}
la coordonnée impaire entre dans la direction
\begin{equation}
\label{rm:eq:modular-ch}
c_h=\begin{pmatrix}-5\\7\\-20\end{pmatrix}.
\end{equation}

Si \(c_A=(2,0,1)^T\) et
\(c_\Delta=(28,-25/2,-2/3)^T\), les trois colonnes
\((c_A,c_h,c_\Delta)\) sont précisément la base transportée par
\(M_{\rm CKM}\). La deuxième ligne de
\cref{rm:eq:modular-Mckm-inverse} est le covecteur \(\ell_h\) qui extrait la
coordonnée impaire :
\begin{equation}
\label{rm:eq:modular-dual-coordinate-check}
\ell_hc_A=0,
\qquad
\ell_hc_h=1,
\qquad
\ell_hc_\Delta=0.
\end{equation}

\begin{theorem}[Réflexion CKM induite par le feuillet]
\label{rm:thm:modular-ckm-reflection}
Soit
\begin{equation}
\label{rm:eq:modular-Dh-lh}
D_h=\operatorname{diag}(1,-1,1),
\qquad
\ell_h=\frac1{3857}\begin{pmatrix}75&176&-150\end{pmatrix}.
\end{equation}
L'involution induite dans l'espace de \(\boldsymbol\theta\) est
\begin{equation}
\label{rm:eq:modular-Rckm}
\boxed{
\mathcal R_{\rm CKM}
=M_{\rm CKM}D_hM_{\rm CKM}^{-1}
=I-2c_h\ell_h.}
\end{equation}
Satisfait
\begin{equation}
\label{rm:eq:modular-Rckm-properties}
\mathcal R_{\rm CKM}^2=I,
\qquad
\det\mathcal R_{\rm CKM}=-1,
\qquad
\mathcal R_{\rm CKM}M_{\rm CKM}=M_{\rm CKM}D_h.
\end{equation}
\end{theorem}

\begin{proof}
La multiplication de
\cref{rm:eq:modular-Mckm,rm:eq:modular-Mckm-inverse} donne l'identité ; en
particulier, \cref{rm:eq:modular-dual-coordinate-check} prouve que
\(c_h\ell_h\) est le projecteur de rang un sur \(\mathbb Qc_h\) le
long de \(\operatorname{span}\{c_A,c_\Delta\}\). Le produit exact
\(\ell_hc_h=1\) implique
\[
(I-2c_h\ell_h)^2
=I-4c_h\ell_h+4c_h(\ell_hc_h)\ell_h=I.
\]
L'opérateur fixe point par point l'hyperplan
\(\ker\ell_h=\operatorname{span}\{c_A,c_\Delta\}\) et change le signe de
\(c_h\). Ses valeurs propres sont donc \(1,1,-1\) ; d'où
\(\det\mathcal R_{\rm CKM}=-1\) et
\(\Tr\mathcal R_{\rm CKM}=1\). Enfin,
\[
\mathcal R_{\rm CKM}M_{\rm CKM}
=M_{\rm CKM}D_hM_{\rm CKM}^{-1}M_{\rm CKM}
=M_{\rm CKM}D_h,
\]
ce qui prouve l'identité d'entrelacement.
\end{proof}

La matrice explicite est
\begin{equation}
\label{rm:eq:modular-Rckm-matrix}
\mathcal R_{\rm CKM}=
\begin{pmatrix}
4607/3857&1760/3857&-1500/3857\\
-150/551&199/551&300/551\\
3000/3857&7040/3857&-2143/3857
\end{pmatrix}.
\end{equation}
Bien que \(\mathcal R_{\rm CKM}\) ne soit pas une réflexion orthogonale pour le
produit euclidien des trois coordonnées angulaires, elle l'est pour la
métrique rationnelle transportée
\begin{equation}
\label{rm:eq:modular-ckm-metric}
G_{\rm CKM}=(M_{\rm CKM}^{-1})^TM_{\rm CKM}^{-1},
\qquad
\mathcal R_{\rm CKM}^TG_{\rm CKM}\mathcal R_{\rm CKM}=G_{\rm CKM}.
\end{equation}
En effet, dans la base \((c_A,c_h,c_\Delta)\), la métrique est l'identité
et l'involution est \(D_h\). C'est pourquoi le terme ``réflexion'' a un
contenu métrique précis et ne décrit pas seulement le fait que
\(\mathcal R_{\rm CKM}^2=I\).
La phase \(\delta_{\rm CP}=9A\) reste fixe sous cette réflexion. C'est
pourquoi \(\mathcal R_{\rm CKM}\) n'est pas la conjugaison CP physique, mais le
transport de l'inversion de feuillet HMT à l'espace de mélange.

L'inverse rationnel de \(M_{\rm CKM}\) retrouve
\begin{align}
A&=\frac{1528\theta_{12}+3380\theta_{23}+801\theta_{13}}{3857},
\label{rm:eq:modular-A-from-ckm}\\
C^*&=\frac{6(75\theta_{12}+176\theta_{23}-150\theta_{13})}{3857}.
\label{rm:eq:modular-C-from-ckm}
\end{align}
La troisième coordonnée se retrouve également sans perte :
\begin{equation}
\label{rm:eq:modular-Delta-from-ckm}
\Delta
=\frac{6\theta_{12}-30\theta_{23}-12\theta_{13}}{551}.
\end{equation}
Ces équations sont des changements de coordonnées. Elles ne font pas de CKM la cause
rétrospective de \(A,C^*\) ou de \(\gammaH\).

Comme reconnaissance terminale, le constructeur préalable détermine
\begin{equation}
\label{rm:eq:modular-ckm-values}
(\theta_{12},\theta_{23},\theta_{13},\delta_{\rm CP})
=(13.003313589509^\circ,
2.398056157332^\circ,
0.213977382244^\circ,
65.676173123554^\circ).
\end{equation}
La réflexion rationnelle est démontrée avant l'utilisation de ces décimales.

\begin{keyresult}
La transition \(6\to8\) est d'abord formalisée par l'inclusion
\(\iota_{6,8}:\mathcal F_6\hookrightarrow\mathcal F_8\), dont on
obtient
\(90,120\) et le défaut \(-1/12\). Le mot de six trits fournit
ensuite la bijection \(729\leftrightarrow729\), et
\(\mathcal J_{6\to8}\) transporte les deux modes collectifs au bord. Sur
cet antécédent, l'aire et le hamiltonien modulaire reconstruisent exactement
les deux familles ; la purification thermochamp double détermine l'entropie
d'intrication, et la première loi
finie conserve le défaut d'entropie relative. Dans CKM, l'inversion de
feuillet devient une réflexion rationnelle de rang impair un.
\end{keyresult}

\begin{falsifier}
La généalogie initiale échoue si \(\iota_{6,8}\) n'est pas injective, si les
décomptes ne sont pas \(90/120\), si \(\delta_{6,8}^{\rm inc}\neq-1/12\), si
\(\beta_{729}\) n'est pas bijective ou si elle est présentée à tort comme un
isomorphisme additif. Le transport collectif échoue si
\(\mathcal J_{6\to8}D_{\rm inc}\neq
D_{\rm area}\mathcal J_{6\to8}\). La reconstruction échoue si le
déterminant des observables de bord n'est pas trente ou si
\cref{rm:eq:modular-inverse-channels} ne retrouve pas les occupations.
La TFD échoue si sa trace partielle n'est pas \(\rho_A\). La première loi ne
s'étend pas aux changements finis sans l'entropie relative de
\cref{rm:eq:modular-finite-first-law}. Le type
\(III_{\lambda_\partial}\) est invalidé si un canal ne persiste pas, s'il n'y a pas de
produit compatible ou si le groupe des rapports change. La réflexion CKM
est réfutée par n'importe laquelle des conditions
\(\mathcal R^2\neq I\), \(\det\mathcal R\neq-1\) ou
\(\mathcal RM\neq MD_h\). Une formule
\(\boldsymbol\theta=f(\gammaH)\) sans entrelaceur viole l'indépendance
typée des deux réalisations fonctionnelles de l'antécédent commun.
\end{falsifier}

\begin{statusbox}
Inclusion \(H\subset O\), dénombrements \(90/120\), deux normalisations du
défaut \(-1/12\), bijection de six trits \(729\leftrightarrow729\),
entrelaceur collectif, échelle \(a_0\), opérateurs finis, spectre d'aire,
TFD, défaut orienté, première loi locale et finie et réflexion CKM :
\status{Théorème interne exact avec ses structures initiales explicites}.
Classification
\(III_{\lambda_\partial}\) : \status{Théorème avec hypothèses explicites de
persistance}. Interprétation de la copie thermochamp double comme extérieur physique :
\status{Réalisation physique conditionnelle}. La grandeur \(\gammaH\) est reprise
du chapitre consacré à Barbero--Immirzi et à l'aire minimale, où elle a été dérivée ; elle intervient ici
comme donnée structurelle déjà établie.
\end{statusbox}

\endgroup
\begingroup
\let\chapter\section
\let\section\subsection
\let\subsection\subsubsection
\let\subsubsection\paragraph
\let\clearpage\relax
% La primera mitad de la fuente de mezcla coincide byte por byte con el
% desarrollo REV.2 ya publicado arriba. Se conserva su continuación exclusiva
% desde una extracción editor-ready y ambas procedencias permanecen intactas.
% Extracto editor-ready del contenido exclusivo de masas/12_mezcla.tex.
% Las líneas 1--104 están publicadas byte-exactamente desde
% 74_rev2_realizaciones_i.tex:254--296 y 406--468. El propietario íntegro
% permanece preservado; este extracto conserva su desarrollo exclusivo.

\subsection{Matrice PMNS comme lecteur de mélange de la fibre neutrinique commune}
Soient \(F_\ell\) et \(F_\nu\) des repères orthonormaux réels de même image tridimensionnelle, avec des orientations choisies de sorte que \(\det(F_\ell^*F_\nu)=+1\), dérivés des opérateurs chargé et neutre. Le transport réel est alors

\[
O=F_\ell^*F_\nu\in SO(3).
\]
Sans choix corrélatif de l'orientation, on conclut seulement \(O\in O(3)\). Soient en outre \(b_K\) le témoin de Witt normalisé, \(Q_K=|b_K\rangle\langle b_K|\) et

\[
\Phi_L=I+(e^{i\pi/6}-1)Q_K.
\]
La phase minimale produit

\[
U_{\rm HMT}=F_\ell^*\Phi_LF_\nu.
\]
Cette compression n'est unitaire que sous la condition de compatibilité

\[
\Phi_L(\operatorname{Im}F_\nu)=\operatorname{Im}F_\ell.
\]
Sans celle-ci, elle produit une contraction ; lorsqu'elle est inversible, son facteur polaire définit le candidat unitaire. Même sous la condition de compatibilité, la phase minimale de rang un est réfutée comme PMNS physique : sa meilleure permutation laisse \(\theta_{13}\simeq17.02^\circ\). Le noyau physique doit lire l'enregistrement complet :

\[
G_{ab}=
\frac1{\sqrt{|L_a||N_b|}}
\sum_{c\in L_a}\sum_{d\in N_b}K_{\rm reg}(c,d),
\qquad
U_{\rm int}=\operatorname{polar}(G),
\]
avec \(\det G\ne0\). Retenue, feuillet, orientation, lacune et holonomie font partie du noyau. Le critère de réfutation de la phase de rang un reste explicitement exposé afin d'éviter de confondre unitarité formelle et clôture physique.

La condition Majorana des neutrinos exige en outre une autoconjugaison fermionique compatible avec l'opérateur de masse. La neutralité, \(C_M\) et le secteur topologique d'Ising ne suffisent pas séparément.

\Needspace{7\baselineskip}
\subsection{Mélange CKM et violation CP}
La grammaire angulaire s'organise au moyen de

\[
M_{\rm CKM}=
\begin{pmatrix}
2&-5&28\\
0&7&-25/2\\
1&-20&-2/3
\end{pmatrix},
\qquad
\det M_{\rm CKM}=-\frac{3857}{6}.
\]
Avec

\[
\begin{pmatrix}\theta_{12}\\\theta_{23}\\\theta_{13}\end{pmatrix}
=M_{\rm CKM}
\begin{pmatrix}A\\C^*/6\\{\left(180/\pi\right)\Delta_4}\end{pmatrix},
\qquad
\delta_{\rm CP}=9A,
\]
on obtient

\[
(\theta_{12},\theta_{23},\theta_{13},\delta_{\rm CP})
=
(13.003313589509,
2.398056157332,
0.213977382244,
65.676173123554)^\circ.
\]
L'invariant de Jarlskog est

\[
J=c_{12}c_{23}c_{13}^2s_{12}s_{23}s_{13}\sin\delta_{\rm CP}
=3.1189723326632974\times10^{-5}.
\]
Le critère linéaire exact de réfutation est

\[
3857\delta_{\rm CP}
=9(1528\theta_{12}+3380\theta_{23}+801\theta_{13}).
\]
Dans la paramétrisation unitaire standard,

\[
V_{\rm CKM}=
\begin{pmatrix}
c_{12}c_{13}&s_{12}c_{13}&s_{13}e^{-i\delta}\cr
-s_{12}c_{23}-c_{12}s_{23}s_{13}e^{i\delta}&
c_{12}c_{23}-s_{12}s_{23}s_{13}e^{i\delta}&s_{23}c_{13}\cr
s_{12}s_{23}-c_{12}c_{23}s_{13}e^{i\delta}&
-c_{12}s_{23}-s_{12}c_{23}s_{13}e^{i\delta}&c_{23}c_{13}
\end{pmatrix}.
\]
L'évaluation interne précédente produit

\[
|V_{\rm CKM}|=
\begin{pmatrix}
0.974350259&0.225005836&0.003734601\cr
0.224873110&0.973489279&0.041841465\cr
0.008582400&0.041121745&0.999117283
\end{pmatrix}.
\]
Une comparaison marginale figée donne les écarts normalisés \((-0.24,-0.07,-0.34,-0.31,-0.37)\) pour les quatre paramètres angulaires et \(J\), respectivement. Ces valeurs ne remplacent pas la comparaison covariante : le texte doit identifier l'édition externe et sa matrice complète de covariance avant de former un \(\chi^2\) conjoint.

CKM assure le transport entre bases de quarks. Si \(B_u,B_d\) sont les opérateurs diagonaux de masse,

\[
B_d^{(u)}=V_{\rm CKM}B_dV_{\rm CKM}^*.
\]
L’identité

\[
|\det[B_u,V_{\rm CKM}B_dV_{\rm CKM}^*]|
=2|J|
\prod_{i<j}|\kappa_i^u-\kappa_j^u|
\prod_{a<b}|\kappa_a^d-\kappa_b^d|
\]
montre que mélange et spectre se couplent sans que la matrice de mélange crée les valeurs propres.

\Needspace{7\baselineskip}

\endgroup

\section{Composition, résonances et secteurs non ordinaires}
La masse composée s’obtient à partir des registres constituants, de la frontière et de l’opérateur de liaison ; ce n’est pas une simple somme de masses. La largeur provient du pôle de survie et demeure séparée de la signature massique. Les sections virtuelles, les secteurs de Fibonacci, d’Ising--Majorana et d’ordre six, ainsi que les exclusions des tachyons et des fantômes conservent leurs domaines propres.

\begingroup
\let\chapter\section
\let\section\subsection
\let\subsection\subsubsection
\let\subsubsection\paragraph
\let\clearpage\relax
\chapter{Composition des routes, des registres et des signatures}
\label{chap:composicion-especies-rev6}
\index{composition des routes}
\index{composition des registres}
\index{méson!expression compositionnelle}
\index{baryon!expression compositionnelle}

Une espèce composée ne peut entrer dans la loi des masses comme un nom auquel on attribue rétrospectivement une ligne de cinq entiers. Avant la signature doit exister un objet susceptible d'être composé. L'atlas du \cref{chap:atlas-genealogia-masas} fournit l'alphabet fini ; l'extension survivante et le registre du \cref{chap:extension-ruta-caracter} fournissent la mémoire susceptible de composition. Ce chapitre a pour objet de construire le niveau intermédiaire sans consulter de masses cibles :
\[
\begin{gathered}
 \text{configuration compositionnelle typée}
 \longrightarrow
 \text{registres complets de ses constituants}\\
 \longrightarrow
 \text{registre composé}
 \longrightarrow
 \ZZ^5.
\end{gathered}
\]

L’ordre des flèches fait partie du résultat. Les expressions se forment avant de recevoir une réalisation physique ; les enregistrements se composent avant l’extraction de la signature ; la signature n’est évaluée par le caractère qu’après la composition de la mémoire. Ainsi, le domaine composé ne se réduit pas aux espèces figurant dans une table métrologique.

\section{Classification interne et grammaire compositionnelle}
\label{sec:alfabeto-composicional-rev6}

Soit
\[
 \mathcal C_{81}=I_9\times I_9,
 \qquad I_9=\{1,\ldots,9\},
\]
l’atlas des cellules, et soit
\[
 \rho(r,c)=(10-r,10-c)
\]
son inversion centrale. Nous notons \(\mathcal L,\mathcal N,\mathcal D,\mathcal U\) les sous-ensembles des leptons chargés, des neutrinos, des quarks de type bas et des quarks de type haut, respectivement, et posons
\[
 \mathcal Q=\mathcal D\sqcup\mathcal U.
\]
Chaque cellule possède un niveau combinatoire \(\operatorname{gen}(x)\). Dans le secteur des quarks, le résidu de couleur est également conservé
\[
 \kappa(x)=r(x)-c(x)\pmod3.
\]
Écrivons
\[
 \mathcal Q_a=\{q\in\mathcal Q:\kappa(q)=a\},
 \qquad a\in\mathbb Z/3\mathbb Z.
\]
Par la \cref{prop:color-atlas},
\[
 |\mathcal Q_0|=|\mathcal Q_1|=|\mathcal Q_2|=12.
\]

\begin{definition}[Dual formel et identité de cellule]
\label{def:dual-identidad-composicional} Pour une section de cellule orientée \(x\), son dual formel \(x^\vee\) inverse l’orientation de charge et de couleur, et a pour support \(\rho(x)\). La notation conserve la distinction
\[
 x^\vee\ne \rho(x)
\]
lorsque \(\rho(x)\) est considérée comme cellule sans orientation. Pour chaque \(x\in\mathcal C_{81}\), nous notons \(\operatorname{id}_x:x\to x\) la flèche identité. L’ensemble des identités est
\[
 \mathfrak I
 :=\{\operatorname{id}_x:x\in\mathcal C_{81}\}.
\]
\end{definition}

L’identité \(\operatorname{id}_x\) est une flèche neutre du système compositionnel. Cette seule condition n’en fait pas une réalisation d’un boson neutre. De même, le dual formel intègre une orientation et ne peut être remplacé par une seconde cellule nue.

\begin{definition}[Expressions mésoniques et baryoniques]
\label{def:expresiones-mesonicas-barionicas} L’ensemble des expressions mésoniques neutres en couleur est
\[
 \mathfrak M
 :=
 \left\{
   M(q,q')=q\otimes{(q')}^\vee:
   q,q'\in\mathcal Q,\ \kappa(q)=\kappa(q')
 \right\}.
\]
Les paires sont ordonnées. Leur résidu total est \(\kappa(q)-\kappa(q')=0\).

L’ensemble des expressions baryoniques à positions de couleur fixées est
\[
 \mathfrak B
 :=
 \mathcal Q_0\times\mathcal Q_1\times\mathcal Q_2,
\]
et nous écrivons
\[
 B(q_0,q_1,q_2)=q_0\otimes q_1\otimes q_2.
\]
Leur neutralité résiduelle est l’identité
\[
 0+1+2=0\pmod3.
\]
\end{definition}

Ces expressions sont des objets syntaxiques typés. En particulier, aucun quotient n’a encore été pris selon les permutations, le spin, la saveur, la statistique, l’énergie de liaison ou l’équivalence dynamique.

\begin{definition}[Flèches chargées compatibles]
\label{def:flechas-cargadas-compatibles} Les flèches leptoniques orientées sont
\[
 \mathfrak A_\ell
 :=
 \left\{
 (x,y)\in
   (\mathcal L\times\mathcal N)
   \sqcup
   (\mathcal N\times\mathcal L):
 \operatorname{gen}(x)=\operatorname{gen}(y)
 \right\}.
\]
Les flèches de quarks orientées sont
\[
 \mathfrak A_q
 :=
 \left\{
 (x,y)\in
   (\mathcal U\times\mathcal D)
   \sqcup
   (\mathcal D\times\mathcal U):
 \begin{array}{l}
 \operatorname{gen}(x)=\operatorname{gen}(y),\\[-2pt]
 \kappa(x)=\kappa(y)
 \end{array}
 \right\}.
\]
Nous posons
\[
 \mathfrak A_{\rm ch}:=\mathfrak A_\ell\sqcup\mathfrak A_q.
\]
\end{definition}

L’égalité de niveau type la transition dans les deux secteurs ; l’égalité de couleur est en outre exigée dans le secteur des quarks. L’orientation distingue \(x\to y\) de \(y\to x\).

\begin{definition}[Système fini de compositions typées]
\label{def:gramatica-celular-compuesta} Le système compositionnel fini est la réunion disjointe typée
\[
 \mathfrak G_{\rm comp}
 :=
 \mathcal C_{81}
 \sqcup\mathfrak I
 \sqcup\mathfrak M
 \sqcup\mathfrak B
 \sqcup\mathfrak A_{\rm ch}.
\]
Les \(81\) cellules sont les générateurs primitifs. Les \(81\) identités sont des flèches sur ces mêmes générateurs et non un second atlas.
\end{definition}

\section{Dénombrement combinatoire des compositions de parcours indépendant des données métrologiques}
\label{sec:censo-composicional-rev6}

\begin{theorem}[Dénombrement du système fini de compositions]
\label{thm:censo-gramatica-composicional-rev6} Le système de la \cref{def:gramatica-celular-compuesta} contient
\[
 \boxed{\begin{gathered}
        81\ \text{identités},\qquad
        432\ \text{expressions mésoniques},\\
        1728\ \text{expressions baryoniques},\qquad
        372\ \text{flèches chargées}
 \end{gathered}}
\]
Les flèches chargées se décomposent exactement comme suit
\[
 372=320_{\text{leptoniques}}+52_{\text{quark}}.
\]
Le dénombrement ne dépend que de l’atlas, de son inversion, des niveaux et du résidu de couleur. Il n’utilise ni masses, ni unités, ni signatures publiées. Les quatre cardinaux comptent des expressions compositionnelles internes ou des flèches admissibles du système ; ils ne comptent ni espèces physiques distinctes, ni particules physiques, ni masses.
\end{theorem}

\begin{proof}
L’application \(x\mapsto\operatorname{id}_x\) est bijective, de sorte que
\[
 |\mathfrak I|=|\mathcal C_{81}|=81.
\]
Pour les mésons, la condition de neutralité permet de choisir indépendamment une paire ordonnée dans chaque classe de couleur. Comme chaque classe contient douze cellules,
\[
 |\mathfrak M|
 =\sum_{a\in\mathbb Z/3\mathbb Z}|\mathcal Q_a|^2
 =3\cdot12^2
 =432.
\]
Pour les baryons, les positions \(0,1,2\) sont fixées, et donc
\[
 |\mathfrak B|
 =|\mathcal Q_0||\mathcal Q_1||\mathcal Q_2|
 =12^3
 =1728.
\]

Les familles leptoniques chargées ont pour tailles
\[
 |\mathcal L_{G0}|=1,\qquad
 |\mathcal L_{G2}|=8,\qquad
 |\mathcal L_{G4}|=16,
\]
tandis que les familles neutriniques ont pour tailles
\[
 |\mathcal N_{G1}|=2,\quad
 |\mathcal N_{G2}|=4,\quad
 |\mathcal N_{G3}|=6,\quad
 |\mathcal N_{G4}|=8.
\]
Seuls \(G2\) et \(G4\) sont des niveaux communs. Comme les deux orientations sont comptées,
\[
 |\mathfrak A_\ell|
 =2(8\cdot4+16\cdot8)
 =320.
\]

Pour les flèches de quarks, dans l’ordre des couleurs \((0,1,2)\), les multiplicités de type haut et de type bas aux niveaux communs sont
\[
\begin{array}{ccc}
 &\mathcal U&\mathcal D\\
 G1&(2,1,1)&(0,1,1)\\
 G3&(4,4,4)&(2,2,2).
\end{array}
\]
Le nombre de paires d’une orientation et de même couleur est donc
\[
 (2\cdot0+1\cdot1+1\cdot1)
 +(4\cdot2+4\cdot2+4\cdot2)
 =2+24=26.
\]
En incluant l’orientation opposée, on obtient
\[
 |\mathfrak A_q|=2\cdot26=52.
\]
Enfin,
\[
 |\mathfrak A_{\rm ch}|=320+52=372.
\]
Aucune de ces règles ne contient de valeur cible ni de vecteur de masse.
\end{proof}

\begin{corollary}[Constructeur compositionnel sans cibles]
\label{cor:constructor-composicional-sin-blancos} Le constructeur qui énumère conjointement les identités, les expressions mésoniques, les expressions baryoniques et les flèches chargées se factorise par les champs typés
\[
 (\text{secteur},\text{famille},\text{niveau},\kappa,\rho)
\]
de l’atlas. Avec la notation précédente, ses quatre codomaines sont \(\mathfrak I,\mathfrak M,\mathfrak B,\mathfrak A_{\rm ch}\). En particulier, le dénombrement peut être fixé avant l’ouverture de toute table de signatures.
\end{corollary}

\begin{proof}
Les quatre définitions et le calcul du \cref{thm:censo-gramatica-composicional-rev6} n’emploient que ces champs. La table de signatures n’appartient au domaine d’aucune des règles.
\end{proof}

L’indépendance à l’égard des cibles ne signifie pas que les expressions ont déjà acquis un nom physique. Elle signifie quelque chose d’antérieur et de plus précis : il existe un domaine compositionnel construit sans utiliser comme critère la quantité qui sera évaluée ensuite.

\section{Transport de mémoire du parcours composé à la signature entière}
\label{sec:memoria-composicion-especies-rev6}

Soit \(\gamma\in\Gamma_{108}^{\rm enr}\) une présentation enrichie d’une extension. Son enregistrement lu contient
\[
 \mathcal E(\gamma)
 =
 (n_A,e_0,\ldots,e_{11},T_\zeta,C_\zeta),
\]
outre les données de phase, de feuillet, d’orientation et de bord qui typent une composition. Les douze événements ne peuvent être remplacés par leur total avant la composition. Pour \(0\le i<6\), nous définissons
\[
 p_i=e_i+e_{i+6},
 \qquad
 a_i=e_i-e_{i+6},
\]
\[
 s_C=\sum_{i=0}^{5}p_i,
 \qquad
 M_i=p_i-p_5\quad(0\le i<5).
\]
Avec \(M_5=(M_0,\ldots,M_4)\) et \(A_6=(a_0,\ldots,a_5)\), la reconstruction est
\[
 p_5=\frac{s_C-\sum_{i=0}^{4}M_i}{6},
 \qquad
 p_i=p_5+M_i,
\]
\[
 e_i=\frac{p_i+a_i}{2},
 \qquad
 e_{i+6}=\frac{p_i-a_i}{2}.
\]
D’après le \cref{thm:memoria-1-5-6}, on obtient sur l’image entière la décomposition réversible
\[
 \boxed{(s_C,M_5,A_6),\qquad 1+5+6=12.}
\]
Le premier bloc conserve le total visible ; le deuxième, la forme relative dans le feuillet symétrique ; le troisième, la mémoire entre les deux demi-couronnes.

La coordonnée torsionnelle forte requiert en outre la cochaîne au niveau du pas :
\[
 k_{\Delta_4}=T_\zeta-2C_\zeta.
\]
Elle n’est pas une fonction du seul motif d’événements nuls. Cette différence sera essentielle pour interpréter le losange de non-descente.

\begin{definition}[Donnée de composition orientée]
\label{def:dato-empalme-orientado-rev6} Soient \(\Gamma,\Lambda\) deux enregistrements enrichis. Une donnée de composition est une paire \((g,B)\), où \(g\) aligne la phase, le feuillet, l’orientation et la position à l’interface, et \(B\) est la sous-trace commune, avec sa contribution à l’action, ses douze événements, ses prédécesseurs torsionnels et sa couronne. Sur le domaine où ces données sont compatibles, on définit
\[
 \boxed{
 {(n,e)}_\Gamma\star_{(g,B)}{(n,e)}_\Lambda
 =
 \bigl(
 n_\Gamma+n_\Lambda-n_B,\,
 e_\Gamma+g e_\Lambda-e_B
 \bigr).}
\]
Le mot torsionnel et la couronne se composent au même niveau, avant le calcul de \(T_\zeta,C_\zeta\) et de \(k_{\Delta_4}\).
\end{definition}

La formule est une inclusion-exclusion d’histoires sur une sous-trace identifiée. Elle est démontrée avec une structure initiale explicite : les enregistrements, le transport \(g\), le bord \(B\), l’orientation et le mot des prédécesseurs font partie des données d’entrée. Le système compositionnel ne choisit de lui-même aucune de ces données.

\begin{definition}[Extracteur postérieur à la composition]
\label{def:extractor-posterior-empalme-rev6} Pour un enregistrement enrichi complet, on définit
\[
 L_{\rm tick}(\gamma)
 =
 \bigl(
 n_A,s_C,k_{\Delta_4},
 \nu_{120}^{\rm ev},\nu_{270}
 \bigr)\in\ZZ^5,
\]
où, si \(\tau_m=\operatorname{sgn}(e_m)\) et \(I_a=\{a,a+4,a+8\}\pmod {12}\),
\[
 \nu_{120}^{\rm ev}
 =
 \sum_{a=0}^{3}
 \operatorname{sgn}\left(\sum_{m\in I_a}e_m\right),
 \qquad
 \nu_{270}
 =
 \sum_{m=0}^{11}\tau_m\tau_{m+3}.
\]
La signature d’une composition est définie par
\[
 \operatorname{Sig}(\Gamma,\Lambda;g,B)
 :=
 L_{\rm tick}
 \bigl(\Gamma\star_{(g,B)}\Lambda\bigr).
\]
\end{definition}

Les signes et les autocorrélations s’appliquent à l’enregistrement déjà composé. La définition n’équivaut donc pas à additionner les signatures des constituants.

\section{Diagramme commutatif de composition et de préservation de la fibre}
\label{sec:rombo-no-descenso-composicional-rev6}

Le catalogue historique B1 emploie une projection d’événements qu’il convient de typer séparément. Soit
\[
 \Sigma_{12}=(+,+,+,-,-,-,+,+,+,-,-,-),
 \qquad \tau_m=\operatorname{sgn}(e_m),
\]
et définissons le proxy d’événements nuls
\[
 K_0(e)
 =
 \sum_{\tau_m=0}\Sigma_{12}(m).
\]
La projection exacte est
\[
 F_0(n_A,e)
 =
 \left(
 n_A,\,
 \sum_m e_m,\,
 K_0(e),\,
 \sum_{a=0}^{3}
   \operatorname{sgn}\sum_{m\in I_a}e_m,\,
 \sum_m\tau_m\tau_{m+3}
 \right)\in\ZZ^5.
\]
Pour deux enregistrements alignés, nous écrivons
\[
 (n,e)\boxplus(n',e')=(n+n',e+e').
\]
Cette opération est le cas sans sous-trace commune et avec alignement identité de la composition d’événements.

\begin{proposition}[Obstruction à la descente de la projection B1]
\label{prop:no-descenso-proyeccion-b1-rev6} Il n’existe aucune opération binaire
\[
 \mu:\ZZ^5\times\ZZ^5\longrightarrow\ZZ^5
\]
tel que
\[
 F_0(\Gamma\boxplus\Lambda)
 =
 \mu\bigl(F_0(\Gamma),F_0(\Lambda)\bigr)
\]
pour tous les enregistrements B1.
\end{proposition}

\begin{proof}
Considérons les parcours, écrits dans la carte \((r,c,h_0,v_0)\),
\[
 a=(1,3,-1,+1),
 \qquad
 b=(6,7,-1,+1),
 \qquad
 c=(1,2,-1,-1).
\]
Les deux premiers ont des enregistrements d’événements distincts :
\[
\begin{aligned}
 e(a)&=(5,4,4,-4,-4,-6,\;5,4,4,-4,-4,-6),\\
 e(b)&=(4,5,4,-4,-6,-4,\;4,5,4,-4,-6,-4),
\end{aligned}
\]
mais produisent le même quintuplet :
\[
 F_0(a)=F_0(b)=(8,-2,0,0,-12).
\]
Le parcours de test possède
\[
 e(c)=(1,4,1,-2,-4,-3,\;1,4,1,-2,-4,-3)
\]
et
\[
 F_0(c)=(28,-6,0,-4,-12).
\]
En additionnant d’abord les événements puis en réévaluant, on obtient
\[
 F_0(a\boxplus c)=(36,-8,0,0,-12),
\]
\[
 F_0(b\boxplus c)=(36,-8,0,-2,-12).
\]
Les deux paires d’entrées de l’hypothétique \(\mu\) sont égales, mais les sorties diffèrent en \(\nu_{120}^{\rm ev}\). Cela contredit l’existence de \(\mu\).
\end{proof}

Le dénombrement exhaustif du même domaine contient \(324\) parcours, \(53\) quintuplets proxy et \(198\) états d’événements distincts. Quarante-deux fibres de la projection contiennent plus d’un état enrichi ; quarante d’entre elles sont séparées par un parcours de test. Le losange précédent est le témoin canonique le plus bref de cette perte d’information.

\begin{remark}[Séparation typée entre le représentant \(K_0\) et la torsion forte]
\label{rem:proxy-no-torsion-fuerte-rev6} La troisième coordonnée de \(F_0\) est \(K_0(e)\), qui, dans B1, prend ses valeurs dans \(\{-4,-2,0,2,4\}\). Elle n’est pas identifiée à \(k_{\Delta_4}=T_\zeta-2C_\zeta\). Il existe en effet des parcours de même \(n_A\) et de même vecteur de douze événements auxquels le lecteur torsionnel attribue des valeurs distinctes parce qu’il conserve les prédécesseurs de chaque itération et la couronne.

La \cref{prop:no-descenso-proyeccion-b1-rev6} démontre précisément que la projection B1 à cinq coordonnées ne transporte pas la superposition des enregistrements. Elle n’affirme aucune impossibilité universelle pour toute structure future sur \(\ZZ^5\) et n’autorise pas à remplacer \(K_0\) par \(k_{\Delta_4}\). Dans le lecteur fort, la sortie est définie par la \cref{def:extractor-posterior-empalme-rev6} : on compose d’abord le mot enrichi, puis on en extrait à nouveau les cinq coordonnées.
\end{remark}

\section{Théorème de composition des enregistrements indépendant des données métrologiques}
\label{sec:teorema-composicional-sin-blancos-rev6}

\begin{theorem}[Clôture compositionnelle sans cibles]
\label{thm:cierre-composicional-sin-blancos-rev6} Une fois fixés l’atlas \(\mathcal C_{81}\), la dualité formelle, le résidu de couleur et l’enregistrement enrichi du \cref{chap:extension-ruta-caracter}, les affirmations suivantes sont vérifiées.
\begin{enumerate}
 \item Il existe le système fini de compositions typées \(\mathfrak G_{\rm comp}\) avec \(81\) identités, \(432\) expressions mésoniques, \(1728\) expressions baryoniques et \(372\) flèches chargées, décomposées comme \(320+52\).

 \item Le système est construit sans masses cibles. Sa sortie est une expression, une orbite ou une multisection typée, et non une masse ni une étiquette physique individualisée.

 \item Étant donnés une expression et des enregistrements constituants munis d’un arbre explicite de compositions compatibles \((g,B)\), l’opération produit un enregistrement enrichi. Sa signature est obtenue par l’application
 \[
 \boxed{
 \bigl(
 \text{expression},
 \text{registres},
 \text{compositions orientées}
 \bigr)
 \xrightarrow{\ \mathfrak C_{12}\ }
 \text{registre composé}
 \xrightarrow{\ L_{\rm tick}\ }
 \ZZ^5.}
 \]
 Ici, \(\mathfrak C_{12}\) désigne l’itération de \(\star_{(g,B)}\) avec le parenthésage et les bords inclus dans les données.

 \item La projection B1 vers cinq entiers ne constitue pas à elle seule un espace compositionnel : le diagramme avec \(\boxplus\) ne descend pas à \(\ZZ^5\). Par conséquent, remplacer chaque enregistrement par son quintuplet avant la composition fait perdre une information nécessaire à la réévaluation d’au moins \(\nu_{120}^{\rm ev}\).
\end{enumerate}
\end{theorem}

\begin{proof}
Le premier point est le \cref{thm:censo-gramatica-composicional-rev6}. Le deuxième est le \cref{cor:constructor-composicional-sin-blancos}. Pour le troisième, chaque arête de l’arbre de composition est évaluée au moyen de la \cref{def:dato-empalme-orientado-rev6} ; comme les bords et les transports font partie du domaine, l’itération produit un enregistrement sur lequel \(L_{\rm tick}\) est défini. Aucune signature intermédiaire ne remplace l’enregistrement. Le quatrième point est la \cref{prop:no-descenso-proyeccion-b1-rev6}.
\end{proof}

La formule correcte de la signature composée est donc
\[
 \operatorname{Sig}_{\rm comp}
 =
 L_{\rm tick}\circ\mathfrak C_{12}.
\]
Il ne s’agit pas d’une somme de lignes de masse. Ce n’est qu’ensuite que le caractère formel ou sa spécialisation globale peuvent être appliqués :
\[
 \text{registre composé}
 \xrightarrow{L_{\rm tick}}\ZZ^5
 \xrightarrow{\chi_{\rm form}}\operatorname{Fun}(\mathcal P,\RR_{>0})
 \longrightarrow
 \text{carte dimensionnelle}.
\]

\section{Reconstruction généalogique du mélange, de la génération et de la persistance de saveur}
\label{sec:significado-genealogico-composicion-rev6}

Le théorème précédent complète une strate qu’une liste de masses ne peut montrer. Les cellules primitives ne se limitent pas à recevoir des noms ; elles engendrent des identités, des paires duales, des triplets de couleur et des flèches orientées selon des règles internes. Les objets composés apparaissent ainsi dans le domaine avant toute comparaison métrologique.

La généalogie ne se réduit pas non plus à l’arbre syntaxique. Deux expressions comportant le même nombre de constituants peuvent se composer différemment parce que la phase, le feuillet, l’orientation et la sous-trace commune font partie de la donnée. L’enregistrement conserve ces différences jusqu’à ce que les signes, les corrélations et la torsion aient agi. La masse, lorsqu’une réalisation dimensionnelle existe, sera une lecture ultérieure de cette mémoire composée.

Cet ordre explique pourquoi \(\ZZ^5\) n’est pas le domaine de composition, bien qu’il soit celui du caractère. Le groupe \((\ZZ^5,+)\) existe et le caractère y est multiplicatif ; ce qui n’existe pas pour la projection B1, c’est une loi qui récupère à partir des quintuplets la superposition ayant déjà oublié la disposition des douze événements. La linéarité du caractère ne restitue pas l’information supprimée par l’extracteur.

\section{Du système compositionnel interne à la réalisation physique}
\label{sec:delimitacion-composicion-especies-rev6}

Le \cref{thm:cierre-composicional-sin-blancos-rev6} construit le domaine antérieur aux noms physiques : expressions ordonnées, duaux, triplets de couleur et flèches compatibles. Ses cardinaux comptent des possibilités internes de composition. Une espèce physique, en revanche, est une représentation persistante dotée d’une saveur, d’une couleur, d’un spin, d’une statistique et d’un canal déterminés. Confondre ces deux niveaux transformerait un dénombrement d’expressions en un faux catalogue de particules.

\begin{definition}[Donnée de réalisation physique]
\label{def:dato-realizacion-fisica-rev6} Une donnée de réalisation pour une espèce \(X\) est le tuple
\[
 \mathfrak D_X=
 \bigl(
 \mathfrak t_X,\mathfrak s_X,
 (\mathcal L_{X_i})_{i=1}^{r},
 \mathcal T_X,(g_X,B_X)
 \bigr),
\]
où :
\begin{enumerate}
 \item \(\mathfrak t_X\) fixe le type physique --- élémentaire, nul, mésonique, baryonique, multiquark ou topologique --- et sa représentation ;
 \item \(\mathfrak s_X\) est la multisection ou l’expression interne réalisée ;
 \item \(\mathcal L_{X_i}\) sont les enregistrements enrichis des constituants ;
 \item \(\mathcal T_X\) fixe la saveur, la couleur, le spin, la statistique et la symétrisation ;
 \item \((g_X,B_X)\) fixe l’arbre de composition, ses transports et ses sous-traces communes.
\end{enumerate}
Le tuple est déclaré sans consulter la masse ou la largeur qui seront évaluées ensuite.
\end{definition}

\begin{proposition}[La réalisation modifie l’enregistrement, non la loi]
\label{prop:realizacion-compuesta-ley-unica-rev6} Toute donnée \(\mathfrak D_X\) détermine un enregistrement composé
\[
 \mathcal L_X
 =\mathfrak C_{12}
   \bigl((\mathcal L_{X_i})_{i=1}^{r};
         \mathcal T_X,g_X,B_X\bigr),
\]
une signature \(v_X=L_{\rm tick}(\mathcal L_X)\) et le caractère central universel
\[
 \chi_0(v_X-v_e).
\]
Si la généalogie apporte en outre un raffinement spectral typé \(\eta_X\in\mathcal E_{\mathfrak S}\), l’évaluation complète est
\[
 m_X
 =m_e^{\rm HMT}\,
   \chi_{\rm full}(v_X-v_e,\eta_X).
\]
Le caractère central et le raffinement appartiennent à des niveaux distincts ; la complétude de la tour ne remplace pas l’application prospective produisant \(\eta_X\). Si \(X\) est instable, sa largeur relève du lecteur temporel \(\Gamma_X=\hbar_{\rm int}/\tau_X^{\rm life}\), et non d’une sixième coordonnée de la signature de masse.
\end{proposition}

\begin{proof}
La donnée de réalisation contient exactement les entrées de \(\mathfrak C_{12}\). Le \cref{thm:cierre-composicional-sin-blancos-rev6} produit l’enregistrement enrichi et la \cref{def:extractor-posterior-empalme-rev6} produit sa signature. Le caractère central et son relèvement spectral sont ceux de \cref{eq:caracter-masa-global-rev7,eq:ley-masa-todos-registros-rev6}. La largeur utilise une droite temporelle et possède donc un autre domaine.
\end{proof}

La séparation des niveaux est ainsi fixée :
\[
\boxed{
\begin{gathered}
\text{système compositionnel interne}
\longrightarrow\text{donnée de réalisation}
\longrightarrow\text{registre composé}\\
\longrightarrow\text{signature}
\longrightarrow\text{caractère central}
\longrightarrow\text{raffinement spectral}
\longrightarrow\text{comparaison physique}.
\end{gathered}}
\]
Les \(432\) expressions mésoniques, les \(1728\) expressions baryoniques et les \(372\) flèches chargées n’occupent que la première case et ne constituent pas un inventaire de particules. Le Modèle standard conserve son ensemble typé d’espèces, et chaque espèce occupe la chaîne au moyen de sa propre donnée de réalisation. La règle causale est unique dans tous les secteurs :
\[
 \boxed{\text{composer les extensions et les registres ;\quad lire ensuite la signature}.}
\]

\endgroup
\begingroup
\let\chapter\section
\let\section\subsection
\let\subsection\subsubsection
\let\subsubsection\paragraph
\let\clearpage\relax
\chapter{Bande de particule, anti-section et virtualité}
\label{chap:banda-particula-virtual}
\index{particule!bande de parcours}
\index{antiparticule!anti-section orientée}
\index{particule virtuelle}

La persistance d’un parcours produit une particule avant l’évaluation de sa masse. Cette priorité permet de distinguer trois objets :

\begin{enumerate}
\item la \emph{section orientée}, qui conserve la direction du parcours ;
\item la \emph{bande}, qui réunit une section et sa conjuguée ;
\item la \emph{persistance}, qui décide si la signature atteint ou non une branche globale.
\end{enumerate}

\begin{definition}[Mots dodécaphasiques et conjugaison matérielle]
\label{def:palabras-cm}
Soit
\[
 \mathscr T_{12}^{\rm word}:=\{-1,0,1\}^{12}.
\]
Pour \(\mathbf t=(t_1,\ldots,t_{12})\), nous écrivons
\[
 \operatorname{rev}(\mathbf t)
 =(t_{12},\ldots,t_1),
 \qquad
 C_M(\mathbf t):=-\operatorname{rev}(\mathbf t).
\]
L’inversion \(C_M\) change simultanément le sens de lecture et l’orientation tritique.
\end{definition}

\begin{theorem}[Involution matérielle et secteur fixe]
\label{thm:involucion-cm-censo} L’application \(C_M\) est une involution de \(\mathscr T_{12}^{\rm word}\) et
\[
 \left|\operatorname{Fix}(C_M)\right|=3^6=729.
\]
\end{theorem}

\begin{proof}
Le renversement et le changement de signe commutent et ont tous deux pour carré l’identité ; par conséquent, \(C_M^2=I\). Un mot est fixe si et seulement si
\[
 t_i=-t_{13-i},\qquad 1\leq i\leq6.
\]
Les six premières coordonnées sont libres et les six dernières sont déterminées par elles. Chaque coordonnée libre admet trois valeurs, de sorte que le secteur fixe a pour cardinal \(3^6\).
\end{proof}

Cette conjugaison finie ne s’obtient pas en changeant uniquement le signe visible : elle inverse aussi l’orientation de l’enregistrement.

\section{La bande de parcours}

Soit \(\mathscr T_{12}^{\rm word}\) l’espace de mots de la \cref{def:palabras-cm}. Nous ajoutons une orientation de section \(\chi\in\{-1,+1\}\) et définissons
\[
(\mathbf t,\chi)\sim(C_M\mathbf t,-\chi).
\]

\begin{definition}[Bande dodécaphasique de parcours]
\label{def:banda-dodecafasica-ruta} La bande de la section orientée \((\mathbf t,\chi)\) est la classe
\[
[(\mathbf t,\chi)]_M
=\{(\mathbf t,\chi),(C_M\mathbf t,-\chi)\},
\]
et l’espace des bandes est
\[
\mathcal B_{12}
:=
\bigl(\mathscr T_{12}^{\rm word}\times\{-1,+1\}\bigr)/{\sim}.
\]
\end{definition}

\begin{proposition}[Liberté de la conjugaison orientée]
\label{prop:accion-cm-orientada-libre}
L’action
\[
(\mathbf t,\chi)\longmapsto(C_M\mathbf t,-\chi)
\]
est libre. Par conséquent,
\[
|\mathcal B_{12}|=3^{12}.
\]
\end{proposition}

\begin{proof}
Un point fixe exigerait \(\chi=-\chi\), ce qui est impossible pour \(\chi\in\{-1,+1\}\). Toutes les orbites ont deux éléments et le domaine possède \(2\cdot3^{12}\).
\end{proof}

La bande n’efface pas l’orientation : elle la conserve comme revêtement double. La particule et l’antiparticule sont les deux sections orientées de la même bande.

\section{Masse paire et charge impaire}

Soit \(E_0:\mathscr T_{12}^{\rm word}\to\mathbb R\) un caractère additif de parcours. Ses composantes paire et impaire par rapport à \(C_M\) sont
\[
E_+(\mathbf t)
=\frac12\bigl(E_0(\mathbf t)+E_0(C_M\mathbf t)\bigr),
\]
\[
E_-(\mathbf t)
=\frac12\bigl(E_0(\mathbf t)-E_0(C_M\mathbf t)\bigr).
\]
Alors
\[
E_+(C_M\mathbf t)=E_+(\mathbf t),
\qquad
E_-(C_M\mathbf t)=-E_-(\mathbf t).
\]

\begin{definition}[Caractère de bande]
\label{def:caracter-banda} Pour une section orientée, nous définissons
\[
E_{\rm band}(\mathbf t,\chi)
=E_+(\mathbf t)+\chi E_-(\mathbf t),
\qquad
m(\mathbf t,\chi)
=m_\circ\exp\bigl(E_{\rm band}(\mathbf t,\chi)\bigr),
\]
où \(m_\circ\) est l’unité interne de comparaison de la famille.
\end{definition}

\begin{theorem}[Principe de conjugaison de masse]
\label{thm:principio-conjugacion-masa-banda} Le caractère et la masse passent au quotient des bandes :
\[
E_{\rm band}(C_M\mathbf t,-\chi)
=E_{\rm band}(\mathbf t,\chi),
\]
\[
\boxed{
m(C_M\mathbf t,-\chi)=m(\mathbf t,\chi).}
\]
Les invariants pairs appartiennent à la bande ; les invariants impairs distinguent l’orientation de ses deux sections.
\end{theorem}

\begin{proof}
D’après les parités précédentes,
\[
\begin{aligned}
E_{\rm band}(C_M\mathbf t,-\chi)
&=E_+(C_M\mathbf t)-\chi E_-(C_M\mathbf t)\\
&=E_+(\mathbf t)+\chi E_-(\mathbf t).
\end{aligned}
\]
L’égalité des masses s’obtient en appliquant l’exponentielle.
\end{proof}

En particulier, une forme linéaire à poids symétriques est impaire, tandis qu’une forme quadratique compatible avec le renversement est paire. La charge et l’orientation changent ; la masse, construite sur la bande, demeure.

\begin{corollary}[Électron et positron]
\label{cor:electron-positron-banda} Dans l’équation électronique de bande, la paire \((n,\chi)=(-22,+1)\) et sa conjuguée \((+22,-1)\) produisent le même terme \(\chi n=-22\). L’électron et le positron possèdent donc la même masse et des orientations de charge opposées.
\end{corollary}

\section{Clôture finie de l’atlas sous anti-section}

Dans les quatre-vingt-un parcours CT--108 apparaissent sept formes tritiques. Le dénombrement exhaustif distingue trois opérations :
\[
\mathbf t\longmapsto-\mathbf t,\qquad
\mathbf t\longmapsto\operatorname{rev}(\mathbf t),\qquad
\mathbf t\longmapsto-\operatorname{rev}(\mathbf t).
\]

\begin{proposition}[Clôture des formes CT--108]
\label{prop:cierre-antiseccion-ct108} L’ensemble des formes tritiques de CT--108 n’est stable ni sous le seul changement de signe ni sous le seul renversement, et il est stable sous \(C_M=-\operatorname{rev}\). Ses orbites possèdent les multiplicités
\[
\boxed{81=54+9+9+9.}
\]
L’orbite de multiplicité \(54\) est autoconjuguée ; les trois autres échangent des paires de formes.
\end{proposition}

\begin{proof}
La forme fixe est
\[
\texttt{+++---+++---}
\]
et possède la multiplicité \(54\). Les trois paires
\[
\begin{array}{c}
\texttt{+++--0+++--0}\ /\ \texttt{0++---0++---},\\
\texttt{+++-0-+++-0-}\ /\ \texttt{+0+---+0+---},\\
\texttt{+++0--+++0--}\ /\ \texttt{++0---++0---}
\end{array}
\]
possèdent chacune la multiplicité totale \(9\). L’application \(-\operatorname{rev}\) fixe la première et permute les membres de chaque paire. Le changement de signe et le renversement, pris séparément, produisent des mots extérieurs à cette liste. La somme des multiplicités est \(54+9+9+9=81\).
\end{proof}

Le résultat n’identifie pas le secteur \(54\) à un autre objet de même cardinal. Il démontre quelle conjugaison respecte la taxonomie concrète des parcours et pourquoi particule et antiparticule partagent une bande.

\section{Persistance finie et particule virtuelle}

Soit
\[
\cdots\xrightarrow{\rho_{M+1,M}}S_M
\xrightarrow{\rho_{M,M-1}}S_{M-1}\xrightarrow{}\cdots
\]
le système des états admissibles aux profondeurs finies. Pour \(s_N\in S_N\), nous définissons l’ensemble des prolongements
\[
\mathcal P_M(s_N)
:=
\{s_M\in S_M:\rho_{M,N}(s_M)=s_N\},
\qquad M\geq N.
\]
Sa profondeur de vie est
\[
\mathcal L(s_N)
:=\sup\{M\geq N:\mathcal P_M(s_N)\neq\varnothing\}
\in\mathbb N\cup\{\infty\}.
\]

\begin{definition}[Persistance et virtualité]
\label{def:persistencia-virtualidad} Un préfixe \(s_N\) est \emph{persistant} s’il appartient à une section globale \((s_M)_{M\geq N}\) compatible. Il est \emph{virtuel fini} si \(\mathcal L(s_N)<\infty\), s’il existe un prolongement jusqu’à \(\mathcal L(s_N)\) et s’il n’en existe aucun à la profondeur suivante.
\end{definition}

\begin{theorem}[Critère de persistance]
\label{thm:criterio-persistencia-virtual} Supposons que chaque état possède un nombre fini de prolongements immédiats et que les restrictions soient cohérentes. Alors :
\[
\mathcal L(s_N)=\infty
\quad\Longleftrightarrow\quad
s_N\text{ appartient à une section globale}.
\]
Si \(\mathcal L(s_N)=L<\infty\), le préfixe transporte le parcours, la mémoire, l’enregistrement et la signature jusqu’à \(L\), mais ne définit pas une particule persistante.
\end{theorem}

\begin{proof}
Une section globale fournit un prolongement à toute profondeur, de sorte que \(\mathcal L=\infty\). Réciproquement, les prolongements de \(s_N\) forment un arbre enraciné, infini et finiment ramifié. D’après le lemme de Kőnig, il contient une branche infinie, dont les sommets sont compatibles par construction et définissent une section globale.

Si \(\mathcal L=L<\infty\), il existe un état à la profondeur \(L\) et aucun à \(L+1\). Les lecteurs définis sur les préfixes peuvent l’évaluer jusqu’à \(L\) ; l’absence de prolongement empêche de transformer cette signature finie en une section de la limite inverse.
\end{proof}

La particule virtuelle possède ainsi une définition interne. Elle n’est ni une particule ``moins réelle'' ni une brève violation d’une loi de conservation. C’est un parcours dont la mémoire agit pendant une profondeur finie et dont la généalogie s’éteint avant de se stabiliser. La confrontation avec les propagateurs et les états intermédiaires de la physique conventionnelle intervient après l’identification du lecteur qui représente ces profondeurs.

\section{Chaîne parcours–enregistrement–signature–caractère–tour de la réalisation matérielle}

Les résultats de ce chapitre fixent la direction complète :
\[
\boxed{
\text{extension}\to
\text{route}\to
\text{registre}\to
\text{signature}\to
\text{bande orientée}\to
\text{particule}\to
\text{masse}.}
\]
L’antiparticule est la seconde section de la bande. La virtualité est une persistance finie. Aucune de ces notions n’est définie à partir d’une masse observée.

\endgroup
\begingroup
\let\chapter\section
\let\section\subsection
\let\subsection\subsubsection
\let\subsubsection\paragraph
\let\clearpage\relax
% Vista científica exacta materializada desde una rama condicional.
% Fuente: source/demostraciones_primarias/31_06c_sectores_topologicos_no_omision.tex; rama: HMTIncludeCMBlock.
% No modifica el contenido matemático de la rama de origen.
\chapter[Structure complexe \texorpdfstring{\(\widetilde J_M^{\,2}=-I\)}{J M au carré = -I} et secteurs topologiques]
{Structure complexe \texorpdfstring{\(\widetilde J_M^{\,2}=-I\)}{J M au carré = -I} sur la strate fixe, involutions dodécaphasiques et secteurs topologiques}
\label{chap:sectores-topologicos-no-omision}
\index{secteur topologique}\index{Fibonacci, anneau basé}
\index{Majorana}\index{Ising}\index{groupe de tresses}

L’involution \(C_M\), l’anneau basé associé à \(F_{\rm av}\), les caractères de tresse et les catégories pointées appartiennent à des strates algébriques distinctes. Ils sont présentés avec leurs domaines et leurs contrôles d’identification afin qu’une coïncidence de nom ne remplace pas le morphisme qui devrait les relier. Le dénombrement de l’atlas, la hiérarchie spectrale infinie dont les huit hauteurs historiques sont les témoins initiaux, et les complétions statistiques conservent les démonstrations de \cref{thm:censo-atlas-81,thm:descenso-necesario-fav,%
thm:principio-exclusion-pauli,thm:dos-cuatro-pi-estadistica}.

Chaque résultat énonce directement son domaine, sa construction et sa conclusion ; aucune coïncidence de nom ne modifie ces dépendances.

\section{Strate périodique de l’involution dodécaphasique}
\label{sec:involucion-cm-censo}

Le domaine \(\mathscr T_{12}^{\rm word}\), la conjugaison \(C_M\mathbf t=-\operatorname{rev}(\mathbf t)\), son caractère involutif et le dénombrement \(\lvert\operatorname{Fix}(C_M)\rvert=3^6=729\) ont leur emplacement principal dans les \cref{def:palabras-cm,thm:involucion-cm-censo}, au sein de la bande de particule. Seule est ajoutée ici la condition périodique nécessaire à la strate topologique. Si \(r\) désigne le décalage cyclique \((r\mathbf t)_m=t_{m-1}\), nous posons
\[
 \operatorname{Per}_6
 :=\{\mathbf t\in\mathscr T_{12}^{\rm word}:r^6\mathbf t=\mathbf t\}.
\]

\begin{proposition}[Dénombrement de la strate périodique]
\label{prop:estrato-periodico-cm-rev7} En imposant une périodicité de période divisant six,
\[
 \bigl|\operatorname{Fix}(C_M)\cap\operatorname{Per}_6\bigr|
 =3^3=27.
\]
Cette identité est exacte.
\end{proposition}

\begin{proof}
Le \cref{thm:involucion-cm-censo} donne la forme générale des points fixes. Si, de plus, \(r^6\mathbf t=\mathbf t\), la seconde moitié du mot coïncide avec la première. En écrivant les six premières coordonnées, les deux conditions équivalent à
\[
 (t_0,t_1,t_2,t_3,t_4,t_5)
 =(a,b,c,-c,-b,-a).
\]
Les paramètres \(a,b,c\) sont indépendants et prennent trois valeurs ; l’intersection possède donc \(3^3\) éléments.
\end{proof}

La forme générale d’un point fixe, avant d’exiger la période six, est
\[
 (a,b,c,d,e,f,-f,-e,-d,-c,-b,-a).
\]
Cette forme en coordonnées permet de vérifier manuellement l’involution et le dénombrement, et évite d’interpréter l’entier \(729\) au moyen d’un autre objet homonyme.

\begin{remark}[Dénombrement exhaustif de l’involution]
\label{rem:procedencia-n40-cm} L’énumération des \(3^{12}=531\,441\) mots ne relève aucune violation de l’involution et donne
\[
 |\operatorname{Fix}(C_M)|=729,
 \qquad
 |\operatorname{Fix}(C_M)\cap\operatorname{Per}_6|=27.
\]
La démonstration principale et la \cref{prop:estrato-periodico-cm-rev7} expliquent ces deux nombres indépendamment de l’énumération : les points fixes ont six coordonnées libres, et la condition supplémentaire de période six réduit le choix à trois. Le dénombrement exhaustif confirme qu’il n’existe aucune orbite exceptionnelle en dehors de cette description.
\end{remark}

\begin{criterion}[Critères de réfutation du dénombrement \(C_M\)] \label{crit:falsadores-cm} Chacun des faits suivants réfuterait le résultat : un mot tel que \(C_M^2\mathbf t\ne\mathbf t\) ; un point fixe ne possédant pas la forme en coordonnées précédente ; un total exhaustif différent de \(729\) ; ou un total différent de \(27\) après avoir imposé \(r^6\mathbf t=\mathbf t\). La coïncidence de ce \(729\) avec un autre cardinal n’identifie pas leurs domaines.
\end{criterion}

\section{Structure quadratique transportée sur la strate fixe}
\label{sec:estructura-cuadratica-estrato-fijo-cm}

Nous identifions l’alphabet des signes au corps ternaire au moyen de
\[
 \{-1,0,+1\}\xrightarrow{\sim}\mathbb F_3,
 \qquad
 -1\longmapsto2,\quad 0\longmapsto0,\quad +1\longmapsto1.
\]
Avec cette identification, \(\mathscr T_{12}^{\rm word}\cong\mathbb F_3^{12}\), et tant \(C_M=-\operatorname{rev}\) que le décalage \(r\) sont des opérateurs \(\mathbb F_3\)-linéaires. Par conséquent, \(\operatorname{Fix}(C_M)\) et \(\operatorname{Per}_6=\ker(r^6-I)\) sont des sous-espaces vectoriels.

La forme en coordonnées des points fixes définit l’isomorphisme linéaire
\[
 s_{\rm fix}:\mathbb F_3^6\xrightarrow{\sim}\operatorname{Fix}(C_M),
\]
\[
 s_{\rm fix}(a,b,c,d,e,f)
 =(a,b,c,d,e,f,-f,-e,-d,-c,-b,-a).
\]
Soit \(J_W\) l’opérateur de Paley--Witt construit dans le \cref{thm:f9-tres-desde-paley}, avec la base, la polarisation et l’orientation qui y sont déclarées.

\begin{proposition}[Structure complexe \(\widetilde J_M^{\,2}=-I\) sur la strate fixe]
\label{prop:raiz-menos-uno-estrato-fijo-cm} Le transport
\[
\widetilde J_M:=s_{\rm fix}J_Ws_{\rm fix}^{-1}
\in\operatorname{End}_{\mathbb F_3}
   \bigl(\operatorname{Fix}(C_M)\bigr)
\]
satisfait
\[
 \widetilde J_M^{\,2}=-I.
\]
\end{proposition}

\begin{proof}
Par conjugaison,
\[
 \widetilde J_M^{\,2}
 =s_{\rm fix}J_Ws_{\rm fix}^{-1}
  s_{\rm fix}J_Ws_{\rm fix}^{-1}
 =s_{\rm fix}J_W^2s_{\rm fix}^{-1}
 =-I.
\]
\end{proof}

Le résultat est démontré avec la carte \(s_{\rm fix}\), la base, la polarisation et l’orientation comme structure initiale explicite. Il n’identifie pas \(C_M\) à \(J_W\) : le premier est une involution de l’espace des mots et le second une structure quadratique sur sa strate fixe. Il ne démontre pas non plus que le transport soit naturel relativement à toute la dynamique TPK.

Posons
\[
 W_{27}
 :=\operatorname{Fix}(C_M)\cap\operatorname{Per}_6.
\]
L’inclusion
\[
 W_{27}\hookrightarrow\operatorname{Fix}(C_M)
\]
est canonique, et la \cref{prop:estrato-periodico-cm-rev7} prouve que \(\dim_{\mathbb F_3}W_{27}=3\).

\begin{corollary}[Obstruction quadratique dans la strate à vingt-sept éléments]
\label{cor:obstruccion-cuadratica-w27} Le sous-espace \(W_{27}\) ne peut être stable sous \(\widetilde J_M\).
\end{corollary}

\begin{proof}
S’il était stable, la relation \(\widetilde J_M^2=-I\) en ferait un module sur
\[
 \mathbb F_3[X]/(X^2+1)\cong\mathbb F_9.
\]
Tout espace vectoriel sur \(\mathbb F_9\) est de dimension paire lorsqu’il est considéré sur \(\mathbb F_3\), ce qui contredit \(\dim_{\mathbb F_3}W_{27}=3\).
\end{proof}

L’égalité de cardinal entre \(W_{27}\) et les vingt-sept droites de Schläfli ne fournit pas d’application entre ces deux objets. Une telle relation exigerait la construction d’une application d’incidence préservant l’orientation centrale et l’action du stabilisateur, ainsi qu’une clôture dynamique compatible. L’atlas de \(81\) cellules ne s’interpose pas entre \(W_{27}\) et la strate fixe de \(729\) états : l’atlas des particules appartient à un autre domaine.

\section[Niveaux de Majorana]{Stratification mathématique de l’objet de Majorana en quatre niveaux}
\label{sec:cuatro-niveles-majorana}

L’inversion \(C_M\), la parité extérieure, un fermion relativiste de Majorana et un mode zéro de Majorana sont des objets de types distincts :

\begin{enumerate}
 \item \(C_M=-\operatorname{rev}\) agit sur \(\mathscr T_{12}^{\rm word}\). Son involutivité et son dénombrement sont démontrés exactement ; cette construction ne fournit à elle seule ni spineur, ni équation de champ, ni terme de masse.

 \item Le caractère de signe \(q=-1\) et la complétion extérieure appartiennent à l’algèbre d’échange. Relativement aux CAR et à la graduation déclarée, ils produisent \(N_s^2=N_s\), le spectre \(\{0,1\}\) et le retour d’orientation \(2\pi/4\pi\). Ces identités se déduisent des CAR et de la graduation ; elles ne se déduisent pas du seul dénombrement de \(C_M\).

 \item Un fermion relativiste de Majorana est un objet spinoriel muni d’une représentation relativiste, d’une conjugaison de charge compatible, d’une condition d’autoconjugaison, d’un terme de masse et d’une dynamique. Ces données ne font pas partie de la définition de \(C_M\) et ne se déduisent pas de son dénombrement.

 \item Un mode zéro de Majorana, ou, de manière équivalente, le secteur non abélien d’une théorie d’Ising au sens topologique, requiert des objets simples \(\mathbf 1,\psi,\sigma\) avec
 \[
  \psi\otimes\psi\simeq\mathbf 1,\qquad
  \psi\otimes\sigma\simeq\sigma,\qquad
  \sigma\otimes\sigma\simeq\mathbf 1\oplus\psi,\qquad
  d_\sigma=\sqrt2,
 \]
 ainsi qu’un associateur et un tressage cohérents. Le gap, le dédoublement et la plateforme relèvent d’une réalisation physique supplémentaire ; les règles de fusion précédentes ne définissent pas une masse de pôle universelle.
\end{enumerate}

La séparation concerne les domaines et non la nomenclature : involution de mots, parité CAR, spineur autoconjugué et objet simple non abélien sont quatre objets mathématiques distincts. En l’absence d’un morphisme explicite entre eux, aucun n’est identifié à un autre du seul fait qu’ils portent le nom de « Majorana ».

\endgroup
\begingroup
\let\chapter\section
\let\section\subsection
\let\subsection\subsubsection
\let\subsubsection\paragraph
\let\clearpage\relax
% Vista científica exacta materializada desde una rama condicional.
% Fuente: source/demostraciones_primarias/31_06c_sectores_topologicos_no_omision.tex; rama: HMTIncludeTopologicalBlock.
% No modifica el contenido matemático de la rama de origen.
\chapter[Incidence et secteurs topologiques]{Incidence, Fibonacci et secteurs topologiques}
\label{chap:incidencia-fibonacci-topologia}
\index{secteur topologique}\index{Fibonacci, anneau basé}
\index{groupe de tresses}

L’incidence résidu--support prolonge l’atlas vers des structures de fusion et de tresse. L’ordre est essentiel : on construit d’abord l’anneau basé de Fibonacci et ses dimensions ; on étudie ensuite les représentations algébriques de tresses et le centre pointé. Chaque objet conserve son propre domaine et n’est pas identifié à une particule physique parce qu’il partage un nom ou un cardinal.

\section{Incidence résidu--support}
\label{sec:incidencia-residuo-soporte}

Nous fixons les classes de résidus et les supports dans l’ordre
\[
 \begin{aligned}
 C_1&=\{1,4,7\},&
 C_2&=\{2,5,8\},&
 C_0&=\{3,6,9\},\\
 \mathcal D_{\rm act}&=\{1,3,5,7\},&
 \mathcal D_{\rm evt}&=\{2,4,6,8\},&
 \mathcal D_{\bar0}^{(9)}&=\{9\}.
 \end{aligned}
\]
Le dernier symbole signifie « support du résidu nonadique nul ». Il ne désigne ni le domaine complet \(D_9=\{1,\ldots,9\}\) ni un secteur physique de masse nulle.

Soient \(E_{\mathcal D}\cong\mathbb Z^3\) et \(E_C\cong\mathbb Z^3\) les modules libres de bases ordonnées
\[
 (\mathcal D_{\rm act},\mathcal D_{\rm evt},
   \mathcal D_{\bar0}^{(9)})
 \quad\text{et}\quad
 (C_1,C_2,C_0),
\]
respectivement. L’application d’incidence \(B_{\rm res,sup}:E_{\mathcal D}\to E_C\) compte les intersections.

\begin{theorem}[Forme de Smith et image de l’incidence]
\label{thm:incidencia-residuo-soporte} Dans les bases précédentes,
\[
 B_{\rm res,sup}
 =\bigl(|C_i\cap\mathcal D_j|\bigr)_{i,j}
 =
 \begin{pmatrix}
 2&1&0\\
 1&2&0\\
 1&1&1
 \end{pmatrix}.
\]
De plus,
\[
 \det B_{\rm res,sup}=3,\qquad
 \operatorname{SNF}(B_{\rm res,sup})
 =\operatorname{diag}(1,1,3),
\]
\[
 \operatorname{im}B_{\rm res,sup}
 =\{(u,v,w)\in\mathbb Z^3:u+v\equiv0\pmod3\}.
\]
Si
\[
 P=\begin{pmatrix}0&1&0\\1&0&0\\0&0&1\end{pmatrix},
\]
l’échange simultané \(C_1\leftrightarrow C_2\) et \(\mathcal D_{\rm act}\leftrightarrow\mathcal D_{\rm evt}\) est équivariant :
\[
 P B_{\rm res,sup}P=B_{\rm res,sup}
 \qquad\Longleftrightarrow\qquad
 PB_{\rm res,sup}=B_{\rm res,sup}P.
\]
\end{theorem}

\begin{proof}
Les neuf entrées proviennent des intersections
\[
\begin{array}{cccc}
 &\mathcal D_{\rm act}&\mathcal D_{\rm evt}
 &\mathcal D_{\bar0}^{(9)}\\
 C_1&\{1,7\}&\{4\}&\varnothing\\
 C_2&\{5\}&\{2,8\}&\varnothing\\
 C_0&\{3\}&\{6\}&\{9\}.
\end{array}
\]
Le développement suivant la troisième colonne donne
\[
 \det B_{\rm res,sup}
 =\det\begin{pmatrix}2&1\\1&2\end{pmatrix}=3.
\]
Le plus grand commun diviseur des entrées vaut un. Le mineur d’ordre deux formé par les première et troisième lignes et les deuxième et troisième colonnes a pour déterminant un :
\[
 \det\begin{pmatrix}1&0\\1&1\end{pmatrix}=1.
\]
D’après les diviseurs de Smith, \(d_1=d_2=1\). Le produit \(d_1d_2d_3\) coïncide avec la valeur absolue du déterminant, que le calcul précédent fixe à trois. Par conséquent,

\[
 d_1d_2d_3=3.
\]
La forme normale est donc \(\operatorname{diag}(1,1,3)\).

Chaque colonne vérifie \(u+v\equiv0\pmod3\), de sorte que l’image est contenue dans le sous-module indiqué. Réciproquement, si \(u+v\equiv0\pmod3\), alors
\[
 x=\frac{2u-v}{3},\qquad
 y=\frac{2v-u}{3},\qquad
 z=w-x-y
\]
sont entiers et
\[
 B_{\rm res,sup}(x,y,z)^{\mathsf T}=(u,v,w)^{\mathsf T}.
\]
Cela prouve l’égalité des images. Enfin, la multiplication explicite par \(P\) échange les deux premières lignes et les deux premières colonnes et laisse la matrice invariante.
\end{proof}

\begin{criterion}[Contrôle de l’incidence] \label{crit:incidencia-residuo-soporte} Une seule intersection mal comptée, un mineur modifiant les diviseurs de Smith, un vecteur de l’image tel que \(u+v\not\equiv0\pmod3\), un vecteur satisfaisant la congruence mais n’admettant pas l’antécédent exhibé, ou la perte de l’équivariance \(PBP=B\) réfuteraient le théorème.
\end{criterion}

\section[Anneau basé de Fibonacci]{Anneau basé associé à l’incidence \texorpdfstring{\(F_{\rm av}\)}{F\_av}}
\label{sec:anillo-basado-fibonacci}

Le \cref{thm:descenso-necesario-fav} démontre, à partir de l’orbite de modes \((\Sigma,\Pi,\Pi)\), la matrice
\[
 F_{\rm av}=
 \begin{pmatrix}0&1\\1&1\end{pmatrix}.
\]
De cette matrice, on obtient l’anneau basé suivant.

\begin{definition}[Anneau basé de Fibonacci]
\label{def:anillo-basado-fibonacci}
Soit
\[
 K_{\rm Fib}:=\mathbb Z\mathbf 1\oplus\mathbb Z\tau
\]
le groupe abélien libre d’unité \(\mathbf 1\) et de produit bilinéaire fixé par
\[
 \tau\cdot\mathbf 1=\tau,\qquad
 \tau\cdot\tau=\mathbf 1+\tau.
\]
Dans cette strate, \(+\) et \(\cdot\) sont l’addition et la multiplication de l’anneau. Les symboles catégoriques \(\oplus\) et \(\otimes\) sont réservés à une catégorification choisie ultérieurement.
\end{definition}

\begin{theorem}[Réalisation de \(F_{\rm av}\) et dimension de Perron]
\label{thm:anillo-fibonacci-fav} Dans la base ordonnée \((\mathbf 1,\tau)\), la multiplication par \(\tau\) dans \(K_{\rm Fib}\) a pour matrice \(F_{\rm av}\). Si \(F_0=0\), \(F_1=1\) et \(F_{n+1}=F_n+F_{n-1}\) pour \(n\geq1\), alors
\[
 \tau^n=F_{n-1}\mathbf 1+F_n\tau
 \qquad(n\geq1).
\]
L’unique homomorphisme de dimension positif envoyant \(\mathbf 1\) sur un vérifie
\[
 \operatorname{FPdim}(\tau)=\varphi.
\]
L’identification des deux feuillets typés à la base \((\mathbf 1,\tau)\) détermine entièrement cet anneau et sa dimension positive.
\end{theorem}

\begin{proof}
Les colonnes de la matrice de multiplication par \(\tau\) sont les coordonnées de
\[
 \tau\cdot\mathbf 1=\tau
 \quad\text{et}\quad
 \tau\cdot\tau=\mathbf 1+\tau;
\]
et valent donc \((0,1)^{\mathsf T}\) et \((1,1)^{\mathsf T}\), respectivement. On obtient ainsi \(F_{\rm av}\).

La formule des puissances est vraie pour \(n=1\). Si elle est vraie pour \(n\), alors
\[
 \tau^{n+1}
 =F_{n-1}\tau+F_n(\mathbf 1+\tau)
 =F_n\mathbf 1+F_{n+1}\tau,
\]
ce qui achève la récurrence. Enfin, une dimension multiplicative positive \(d=\operatorname{FPdim}(\tau)\) vérifie
\[
 d^2=1+d.
\]
Ses racines sont \(\varphi\) et \(-\varphi^{-1}\) ; la positivité sélectionne \(d=\varphi\).
\end{proof}

La dénomination \emph{Fibonacci} s’explique ainsi par la règle de fusion elle-même : en itérant \(\tau\), les deux canaux possibles apparaissent avec des multiplicités consécutives \(F_{n-1}\) et \(F_n\). L’espace des états de fusion croît donc avec le rapport asymptotique \(\varphi\). C’est le fondement combinatoire de son utilisation computationnelle : l’information est codée dans des canaux globaux de fusion et non dans une masse attribuée à un point.

\begin{criterion}[Conditions d’une catégorification de Fibonacci] \label{crit:categorificacion-fibonacci} La couche annulaire serait réfutée si la matrice de multiplication n’était pas \(F_{\rm av}\), si la récurrence des puissances échouait ou si la dimension positive n’était pas \(\varphi\). Une catégorification tressée dont l’anneau de Grothendieck est celui de la \cref{def:anillo-basado-fibonacci} exige en outre un foncteur monoïdal
\[
 \mathfrak C_{\rm Fib}:\mathsf{TPK}
 \dashrightarrow\mathcal C_{\rm Fib},
\]
la vérification du pentagone et de l’hexagone dans son image, ainsi qu’une règle interne sélectionnant la jauge et la chiralité. Ces conditions ne découlent pas de la seule identité annulaire.
\end{criterion}

\section[Caractères de tresse]{Six caractères et quatre torsions de tresse}
\label{sec:rescate-trenzas-ex02}

Soit \(B_n\), \(n\geq2\), le groupe de tresses d’Artin de générateurs \(\sigma_1,\ldots,\sigma_{n-1}\) et de relations
\[
 \sigma_i\sigma_{i+1}\sigma_i
 =\sigma_{i+1}\sigma_i\sigma_{i+1},
 \qquad
 \sigma_i\sigma_j=\sigma_j\sigma_i\quad(|i-j|\geq2),
\]
et soit \(w:B_n\to\mathbb Z\) la somme des exposants. Les lecteurs
\[
 w_2:=w\bmod2,\qquad w_3:=w\bmod3
\]
se réunissent, par le théorème chinois des restes, en
\[
 w_6:B_n\longrightarrow
 \mathbb Z/2\mathbb Z\times\mathbb Z/3\mathbb Z
 \simeq\mathbb Z/6\mathbb Z.
\]

\begin{theorem}[Abélianisation, six caractères et torsions de permutation]
\label{thm:trenzas-caracteres-z6}
Pour \(n\geq2\),
\[
 B_n^{\rm ab}\simeq\mathbb Z,
\]
au moyen de \(w\). Par conséquent, \(w_6\) produit six caractères unitaires
\[
 \chi_k(b)
 :=\exp\!\left(\frac{2\pi\mathrm i k}{6}w(b)\right),
 \qquad k\in\mathbb Z/6\mathbb Z.
\]

Soit \(V\ne0\) un espace de Hilbert, soit \(P_i:V^{\otimes n}\to V^{\otimes n}\) l’échange des facteurs \(i,i+1\), et soit \(q\in\mu_6\). La règle
\[
 \rho_q(\sigma_i):=qP_i,
 \qquad
 \rho_q(b)=q^{w(b)}\pi_n(b),
\]
où \(\pi_n\) est la représentation de permutation, définit une représentation de \(B_n\). Cette représentation descend au groupe symétrique si et seulement si
\[
 q^2=1.
\]
Pour les quatre valeurs \(q\in\mu_6\setminus\{\pm1\}\),
\[
 \operatorname{Inv}(\rho_q)=0.
\]
Ces quatre valeurs définissent des torsions scalaires véritablement tressées, mais ne constituent pas à elles seules quatre espaces de Fock ni quatre secteurs anyoniques physiques.
\end{theorem}

\begin{proof}
Dans tout quotient abélien, la relation de tresse donne
\[
 2[\sigma_i]+[\sigma_{i+1}]
 =[\sigma_i]+2[\sigma_{i+1}],
\]
tous les générateurs ont donc la même classe. Aucune autre relation ne porte sur cette classe, puisque \(w(\sigma_i)=1\) ; l’abélianisation est donc \(\mathbb Z\).

Les opérateurs \(P_i\) vérifient les relations de tresse et \(P_i^2=I\). Multiplier chaque générateur par le même scalaire \(q\) conserve les deux relations, de sorte que \(\rho_q\) est une représentation. Pour se factoriser par \(S_n\), elle doit satisfaire la relation supplémentaire \(\sigma_i^2=1\), et
\[
 \rho_q(\sigma_i)^2=q^2I.
\]
Cela prouve le critère \(q^2=1\).

Si \(v\) était invariant, alors \(qP_i v=v\), c’est-à-dire \(P_i v=q^{-1}v\), pour tout \(i\). Comme chaque involution \(P_i\) n’a pour valeurs propres que \(\pm1\), il n’existe aucun tel vecteur non nul lorsque \(q\ne\pm1\). Par conséquent, \(\operatorname{Inv}(\rho_q)=0\) pour les quatre valeurs restantes.
\end{proof}

Lorsque \(\dim V>1\) et \(n\geq3\), les opérateurs \(P_1\) et \(P_2\) ne commutent pas ; la représentation complète ne devient donc pas abélienne du fait qu’elle provient d’un caractère scalaire. Le caractère \(\chi_k\) et la représentation \(\rho_q\) sont des objets distincts : le premier est unidimensionnel ; la seconde conserve l’action de permutation.

\section[Centre modulaire et réalisations topologiques]{Centre modulaire pointé et conditions de réalisation de Fibonacci et d’Ising}
\label{sec:centro-apuntado-no-go}

Soit
\[
 A=(\mathbb Z/9\mathbb Z)^2
\]
et soit \(\operatorname{Vec}_A\) la catégorie des espaces vectoriels \(A\)-gradués à associateur trivial. Fixer un bicharactère non dégénéré identifie \(\widehat A\) à \(A\), sans modifier les dénombrements suivants.

\begin{theorem}[Centre pointé et obstruction dimensionnelle]
\label{thm:centro-apuntado-no-go} Le centre de Drinfeld
\[
 \mathcal Z(\operatorname{Vec}_A)
\]
est une catégorie modulaire pointée dont les objets simples sont paramétrés par \((a,\chi)\in A\times\widehat A\). Elle possède donc
\[
 |A|\,|\widehat A|=81^2=6561
\]
objets simples, tous de dimension quantique un, et une dimension quantique totale
\[
 \mathcal D
 =\sqrt{\sum_{x\ {\rm simple}}d_x^2}
 =\sqrt{6561}=81.
\]
Dans cette catégorie, on ne peut réaliser, en préservant les objets simples et les dimensions de Frobenius--Perron, ni l’anneau de Fibonacci ni le secteur d’Ising : ils requièrent respectivement un objet simple de dimension \(\varphi\) et un objet simple de dimension \(\sqrt2\).
\end{theorem}

\begin{proof}
Pour un groupe abélien fini et un associateur trivial, chaque objet simple du centre est déterminé par un degré \(a\in A\) et un caractère \(\chi\in\widehat A\) qui spécifie son demi-tressage. L’objet gradué et le caractère sont tous deux unidimensionnels ; tous les objets simples sont donc inversibles et de dimension quantique un. Comme \(|A|=|\widehat A|=9^2=81\), le nombre de paires est \(81^2\), et la formule de \(\mathcal D\) donne \(81\).

Toute sous-catégorie de fusion d’une catégorie pointée est elle-même pointée ; ses objets simples ont en particulier dimension un. Cela contredit \(\operatorname{FPdim}(\tau)=\varphi\) pour Fibonacci et \(d_\sigma=\sqrt2\) pour Ising. L’obstruction est dimensionnelle et ne dépend pas d’une ressemblance de noms ou de cardinaux.
\end{proof}

\begin{criterion}[Conditions d’une réalisation physique des secteurs] \label{crit:promocion-sectores-topologicos} Les résultats de \cref{thm:trenzas-caracteres-z6,thm:centro-apuntado-no-go} seraient algébriquement réfutés par un défaut des relations d’Artin, du critère \(q^2=1\), du dénombrement \(6561\), du caractère pointé ou de \(\mathcal D=81\). Une réalisation physique exige en outre un Hamiltonien ou un protocole piloté, un gap spectral, des opérateurs de tresse réalisables, une protection contre les perturbations et une application issue des multisections enrichies. Aucune de ces données ne se déduit uniquement des deux théorèmes algébriques précédents.
\end{criterion}

\endgroup
\begingroup
\let\chapter\section
\let\section\subsection
\let\subsection\subsubsection
\let\subsubsection\paragraph
\let\clearpage\relax
% Las formulaciones iniciales de esta fuente coinciden byte por byte con
% las residencias REV.2 de realizaciones y ontología. Se publica aquí únicamente
% su desarrollo exclusivo posterior, sin alterar ninguna fuente de procedencia.
% Extracto editor-ready del contenido exclusivo de masas/08_composicion.tex.
% Las líneas 1--52 son byte-exactas a 75_rev2_ontologia_particula.tex:231--282
% y las líneas 53--396 son byte-exactas a 74_rev2_realizaciones_i.tex:641--1052
% en dos tramos contiguos. Sus propietarios íntegros permanecen preservados;
% esta residencia publica una sola vez el contenido exacto y conserva desde
% la línea 397 todo el desarrollo exclusivo del propietario de composición.

\subsection{Objet composé comme section persistante d'un produit cocyclique de parcours}
Une particule composée est une section du produit fibré des biographies constituantes, avec une frontière de liaison. Les enregistrements généalogiques sont composés avant l'application de l'extracteur de signature. Si \(\Gamma\) et \(\Lambda\) sont liés par la frontière \(B\) et l'action \(g\),

\[
(n,e)_\Gamma\star_{(g,B)}(n,e)_\Lambda
=
(n_\Gamma+n_\Lambda-n_B,
e_\Gamma+ge_\Lambda-e_B).
\]
La signature composée est

\[
\operatorname{Sig}_{\rm comp}
=L_{\rm tick}\circ\mathfrak C_{12}.
\]
Elle ne coïncide généralement pas avec la somme des signatures individuelles, car l'enregistrement généalogique de liaison porte énergie, retenue et mémoire. C'est là que réside l'énergie de liaison.

\Needspace{7\baselineskip}
\subsection{Résonance comme section transitoire et largeur comme taux de désintégration}
Une résonance est une section métastable dont la résolvante possède un pôle d'énergie complexe. Si \(z_X=M_Xc_{\rm int}^{2}-i\Gamma_X/2\), la masse est \(M_X=c_{\rm int}^{-2}\operatorname{Re}z_X\) ; la largeur mesure sa durée de vie moyenne :

\[
\Gamma_X=\frac{\hbar_{\rm int}}{\tau_X^{\rm life}}.
\]
La largeur n'est pas une sixième coordonnée de \(\sigma\). Elle procède de l'opérateur de survie, des canaux de désintégration et du prolongement analytique de la résolvante. L'inventaire de 384 largeurs exige donc un lecteur propre, distinct de celui de l'inventaire de 471 masses.

\Needspace{7\baselineskip}
\subsection{Section virtuelle comme extension intermédiaire non asymptotique de l'atlas}
Une section virtuelle est une extension finie qui participe à une composition ou à un propagateur sans définir de branche globale stable dans le codomaine physique. Si \(L(s)\) est son horizon,

\[
L(s)<\infty,
\qquad
\operatorname{Ext}_{L(s)}(s)\ne\varnothing,
\qquad
\operatorname{Ext}_{L(s)+1}(s)=\varnothing.
\]
La virtualité est une obstruction terminale, et non une étiquette documentaire ou une masse imaginaire. Une résonance peut être instable sans être virtuelle en ce sens ; une ligne interne peut être virtuelle sans constituer une nouvelle espèce.

\Needspace{7\baselineskip}
\subsection{Section de Majorana comme point fixe de la conjugaison particule–antiparticule}
Quatre constructions doivent rester distinctes. L'involution des mots \(C_M\) définit l'antisection ; un spineur de Majorana satisfait une condition de réalité ; un mode zéro de Majorana est une excitation localisée d'énergie nulle dans un Hamiltonien topologique ; le secteur d'Ising possède les règles de fusion

\[
\sigma\otimes\sigma=1\oplus\psi.
\]
Une particule physique est de Majorana lorsque sa représentation fermionique admet une autoconjugaison compatible avec le Hamiltonien, la localisation et le produit scalaire. La neutralité électrique et l'invariance d'un mot sous \(C_M\) ne suffisent pas. En particulier, déterminer la nature Majorana des neutrinos exige un entrelaceur entre l'opérateur neutre, la base des masses et la structure réelle fermionique.

\Needspace{7\baselineskip}
\subsection{Section anyonique et représentation non triviale du groupe de tresses}
Un anyon est une section bidimensionnelle dont la représentation du groupe de tresses ne se réduit pas à la permutation bosonique ou fermionique. Dans le secteur de Fibonacci,

\[
\tau\otimes\tau=1\oplus\tau,
\qquad
\operatorname{FPdim}(\tau)=\varphi.
\]
Dans le secteur abélien \(\mathbb Z/6\mathbb Z\), le quotient d'enroulement possède six caractères unidimensionnels. Pour obtenir une particule physique, il faut des matrices \(F\) et \(R\), un pentagone, un hexagone, un Hamiltonien à gap et des opérateurs de tressage. L'anneau de fusion exact détermine le codomaine topologique, mais ne fixe pas, à lui seul, une masse universelle.

\Needspace{7\baselineskip}
\subsection{Classes tachyonique et fantôme comme obstructions de positivité et de stabilité}
Une particule persistante possède une évolution unitaire, une énergie bornée inférieurement et une masse spectrale réelle sur la droite physique. Un tachyon exigerait un pôle avec \(m^2<0\) demeurant un état asymptotique stable. Dans l'architecture HMT--MD, une direction qui quitte le cône des réalisations positives revient à la frontière nulle, perd sa survie ou se manifeste comme l'instabilité d'un fond qui doit se condenser. Elle ne définit pas de section persistante supraluminique. L'impossibilité se formule comme une incompatibilité entre persistance globale, positivité spectrale et \(m^2<0\), et non comme une interdiction des paramètres tachyoniques transitoires dans un développement perturbatif.

Un fantôme est un mode de norme négative. La positivité du produit scalaire et la conservation unitaire excluent sa réalisation physique. Les fantômes de Faddeev--Popov appartiennent au complexe de jauge et non à l'espace physique des états ; cette distinction doit être conservée.

\Needspace{7\baselineskip}
\section{Composition, résonances et largeurs}
\Needspace{7\baselineskip}
\subsection{Composition non additive des masses par la loi cocyclique de liaison}
La Loi Générale des Masses agit sur des biographies généalogiques. Lorsque plusieurs constituants forment un état lié, l'objet d'entrée n'est pas la liste de leurs masses isolées, mais un enregistrement composite qui conserve l'ordre du couplage, la frontière commune, la retenue, l'orientation et la mémoire. Soient

\[
\Gamma_1,\ldots,\Gamma_n
\]
des parcours constituants, \(\mathcal T\) un arbre temporel de composition et \(\partial\) la frontière de liaison. L'opérateur dodécaphasique de composition s'écrit

\[
\mathfrak C_{12}^{(n)}:
(\Gamma_1,\ldots,\Gamma_n;\mathcal T,\partial)
\longmapsto \mathcal L_{\rm comp}.
\]
Ce n'est qu'ensuite que s'appliquent le lecteur de l'enregistrement, l'extracteur de signature, le caractère et les tours :

\[
\mathcal L_{\rm comp}
\longmapsto \sigma_{\rm comp}
\longmapsto \chi_{\rm comp}
\longmapsto \widehat M_{\rm comp}.
\]
Cette priorité est essentielle. Si les masses isolées étaient d'abord évaluées puis additionnées, les données qui distinguent un état lié d'un seuil de constituants libres disparaîtraient précisément. L'énergie de liaison n'est pas un terme ajouté à la fin : elle est l'image dimensionnelle d'une différence d'enregistrement produite par la liaison.

\Needspace{7\baselineskip}
\subsection{Loi cocyclique de liaison}
Pour deux enregistrements \(L_1,L_2\), la composition prend la forme

\[
L_1\star L_2
=L_1+L_2+\mathfrak b(L_1,L_2),
\]
où \(\mathfrak b\) enregistre frontière, retenue et mémoire partagées. L'indépendance à l'égard du parenthésage exige l'identité cocyclique

\[
\mathfrak b(L_1,L_2)
+\mathfrak b(L_1\star L_2,L_3)
=
\mathfrak b(L_2,L_3)
+\mathfrak b(L_1,L_2\star L_3).
\]
Si cette égalité échoue, les deux façons de construire un même arbre triple produisent des enregistrements généalogiques distincts et il n'existe pas d'espèce composée bien définie. Si elle est satisfaite, le cocycle ne disparaît pas : il garantit que sa contribution est une propriété de l'état et non un artefact de notation.

Une présentation plus explicite conserve en outre l'action de bord. Soit \(e\in\mathbb Z^{12}\) le vecteur ordonné des événements, et \(n\) le compteur d'action ou d'événements d'un enregistrement. Si \((n,e)_\Gamma\) et \((n,e)_\Lambda\) sont des enregistrements, si \(B\) est le sous-enregistrement commun de données \((n_B,e_B)\), et si \(g\in\operatorname{Aut}(\mathbb Z^{12})\) transporte l'orientation du second enregistrement vers celle du premier en préservant le lecteur,

\[
(n,e)_\Gamma\star_{(g,B)}(n,e)_\Lambda
=
(n_\Gamma+n_\Lambda-n_B,
e_\Gamma+g e_\Lambda-e_B).
\]
Le terme \(B\) évite de compter deux fois la frontière identifiée ; l'action \(g\) transporte le second enregistrement dans l'orientation du premier. La signature composée est donc

\[
\operatorname{Sig}_{\rm comp}
=L_{\rm tick}\!\left(
\mathfrak C_{12}^{(n)}
(\Gamma_1,\ldots,\Gamma_n;\mathcal T,\partial)
\right),
\]
et non la somme automatique des signatures élémentaires.

\Needspace{7\baselineskip}
\subsection{Clôture cocyclique de la composition mésonique quark–antiquark}
Une expression mésonique orientée prend la forme

\[
M(q,q')=q\otimes(q')^\vee.
\]
Le domaine brut contient \(6\times6\times3\times4=432\) expressions lorsque l'on considère saveur, couleur compatible et orientation. Cette cardinalité n'est pas le nombre de mésons physiques. Avant d'individualiser une espèce, on applique :

\begin{enumerate}
\item le projecteur de singulet de couleur ;
\item le projecteur de charge et d'isospin ;
\item la représentation de spin et de parité ;
\item l'arbre temporel de liaison ;
\item la condition radiale ou orbitale ;
\item le canal de préparation et de désintégration lorsque l'état est résonant.
\end{enumerate}
Le descripteur physique complet est

\[
\mathcal I_M=
(q\bar q',J^{PC},I,I_3,S,C,B,\mathcal T,\partial,
\Gamma_M,\mathcal L_M,\sigma_M,\eta_M).
\]
Deux états de mêmes saveurs peuvent différer par \(J^{PC}\), l'excitation radiale, le moment orbital, le feuillet ou le canal, et posséder ainsi des pôles distincts. La masse ne se déduit pas du seul nom des constituants.

Les évaluations publiées pour \(\pi^0,\pi^\pm,K^\pm,K^0\) sont exactes une fois fixés les enregistrements historiques déclarés. L'affirmation plus forte --- la génération prospective de chaque enregistrement nominal à partir de son arbre de liaison --- exige que \(\mathfrak C_{12}\) reconstruise ces enregistrements sans consulter la valeur qui sera comparée. Les deux couches sont conservées ensemble : le calcul conditionnel n'est pas supprimé, mais n'est pas non plus présenté comme un sélecteur d'identité.

\Needspace{7\baselineskip}
\subsection{Clôture trilinéaire de la composition baryonique et neutralisation du caractère de couleur}
Une expression baryonique orientée utilise trois constituants et trois couleurs. Le domaine brut du dénombrement adopté contient

\[
\lvert\mathcal B_{\rm RGB}^{\rm raw}\rvert=1728
\]
expressions. La projection physique doit imposer simultanément :

\[
P_{\rm singlete\ de\ color},
\qquad
P_{J^P},
\qquad
P_{I,I_3,S,C,B},
\qquad
P_{\rm Pauli}.
\]
L'arbre triple conserve la paire qui s'est couplée en premier et la façon dont le troisième constituant a traversé la frontière. L'identité cocyclique garantit que les parenthésages équivalents représentent le même état lorsque l'entrelaceur physique les identifie ; les canaux non équivalents demeurent des états distincts.

Les signatures historiques

\[
\sigma_p=(59,0,9,1,0),
\qquad
\sigma_n=(59,1,-34,0,1)
\]
produisent les évaluations publiées du proton et du neutron lorsque l'on applique le caractère complet. Ce sont des contrôles exacts du lecteur pour ces enregistrements. L'explication physique de la différence neutron--proton exige en outre de dériver, avant la confrontation, le signe de l'isospin, l'orientation des lacunes, le bord de liaison et la contribution électromagnétique dans une carte commune.

\Needspace{7\baselineskip}
\subsection{Réalisations composées exotiques et conditions de stabilité de la liaison}
La même grammaire admet davantage de constituants sans modifier la loi :

\[
\mathfrak C_{12}^{(n)}:
\prod_{j=1}^{n}\Gamma_j\times\mathcal T_n\times\partial_n
\longrightarrow\mathcal L_{\rm comp}.
\]
Un tétraquark correspond à \(qq\bar q\bar q\) ; un pentaquark, à \(qqqq\bar q\) ; une boule de glu, à des parcours de la représentation adjointe de couleur projetés sur le singulet ; un hybride, à des parcours fermioniques et de jauge réunis dans un même arbre. La liste des constituants ne détermine pas, à elle seule, l'espèce. Une classification physique complète doit inclure couleur, spin, parité, conjugaison, isospin, topologie de l'arbre et pôle de l'opérateur couplé.

Cette formulation distingue aussi deux descriptions concurrentes d'un même candidat --- molécule hadronique et état compact, par exemple --- : ce sont deux arbres ou deux projecteurs sur un même secteur de constituants. Elles ne deviennent équivalentes que s'il existe un opérateur unitaire entrelaçant leurs opérateurs, leurs canaux et leurs pôles.

\Needspace{7\baselineskip}
\subsection{État stable, état lié et résonance}
Il convient de distinguer trois notions.

Un état stable est un vecteur propre normalisable du Hamiltonien physique, d'énergie réelle et sans canal de désintégration ouvert. Un état lié est un vecteur propre situé sous le seuil pertinent ; il peut être stable ou se désintégrer par un autre canal. Une résonance est un pôle du prolongement de la résolvante ou de la matrice de diffusion. Elle n'est pas un vecteur propre normalisable du Hamiltonien autoadjoint sur l'axe réel.

Pour un pôle d'énergie

\[
z_X=E_X-\frac{i}{2}\Gamma_X
=M_Xc_{\rm int}^{2}-\frac{i}{2}\Gamma_X,
\qquad \Gamma_X>0,
\]
l'énergie du pôle est \(E_X=\operatorname{Re}z_X\), sa masse est \(M_X=c_{\rm int}^{-2}\operatorname{Re}z_X\) et sa largeur énergétique est \(\Gamma_X=-2\operatorname{Im}z_X\). La Loi Générale des Masses fournit la branche réelle conditionnée par l'identité et l'enregistrement ; le prolongement ouvert des canaux détermine le déplacement complexe et la largeur. Dans une carte d'unités naturelles \(c_{\rm int}=1\), les coordonnées numériques d'énergie et de masse coïncident, mais leurs types ne sont pas identifiés avant le choix de cette carte.

\Needspace{7\baselineskip}
\subsection{Correction du signe dans la lecture discrète du pôle}
Soit \(t_0>0\) un pas temporel et supposons que le mode d'un pas ait pour valeur propre

\[
\lambda=r e^{-i\theta},
\qquad 0<r<1,
\]
avec la convention

\[
\lambda=\exp\!\left(-\frac{i}{\hbar_{\rm int}}z t_0\right).
\]
Alors

\[
z=\frac{\hbar_{\rm int}}{t_0}
\bigl(\theta+i\log r\bigr)
=E-\frac{i}{2}\Gamma
=Mc_{\rm int}^{2}-\frac{i}{2}\Gamma,
\]
de sorte que

\[
E=\frac{\hbar_{\rm int}}{t_0}\theta,
\qquad
M=\frac{\hbar_{\rm int}}{t_0c_{\rm int}^{2}}\theta,
\qquad
\Gamma=-\frac{2\hbar_{\rm int}}{t_0}\log r>0.
\]
Le signe positif devant \(i\log r\) est imposé : puisque \(\log r<0\), la partie imaginaire de \(z\) est négative. La forme \(\theta-i\log r\) placerait le pôle dans le demi-plan supérieur et décrirait une croissance, et non une décroissance.

La probabilité de survie après \(N\) pas est

\[
|\lambda|^{2N}=r^{2N}
=\exp\!\left(-\frac{\Gamma}{\hbar_{\rm int}}Nt_0\right).
\]
La durée de vie associée est donc

\[
\tau_X^{\rm life}=\frac{\hbar_{\rm int}}{\Gamma_X}.
\]
Cette identité démontre aussi pourquoi la largeur ne doit pas être introduite comme sixième coordonnée de la signature \(\mathbb Z^5\) : elle naît du module du mode ouvert, et non du caractère positif qui organise la masse.

\Needspace{7\baselineskip}
\subsection{Opérateur effectif des canaux}
Soit \(P\) le sous-espace préparé et \(Q=I-P\) l'ensemble des canaux éliminés. Le complément est incorporé au moyen de l'opérateur de Feshbach

\[
H_{\rm eff}(z)
=PHP+PHQ(z-QHQ)^{-1}QHP.
\]
Les pôles satisfont

\[
\det\bigl(z-H_{\rm eff}(z)\bigr)=0.
\]
Le premier terme conserve l'opérateur interne préparé. Le second est l'autoénergie des canaux : il déplace la partie réelle et produit une partie imaginaire lorsque le continu s'ouvre. La masse du caractère et la largeur de survie convergent vers le même pôle, mais procèdent d'applications distinctes.

Une formulation équivalente utilise un projecteur de survie distinct du complément de Feshbach. Soit

\[
S_X=S_X^*=S_X^2
\]
le projecteur sur le sous-espace survivant de l'état préparé. L'opérateur comprimé d'un pas est

\[
T_X=S_XUS_X.
\]
Supposons que \(T_X\) possède une valeur propre dominante simple \(\lambda_X\), que l'état préparé ait un recouvrement non nul avec son sous-espace spectral et que le reste du spectre soit séparé en module. Si \(r_X=|\lambda_X|=\rho(T_X)<1\), alors

\[
\Gamma_X=-\frac{2\hbar_{\rm int}}{t_0}\log r_X.
\]
Si l'opérateur est normalisé dès le départ sur les probabilités plutôt que sur les amplitudes, le facteur deux disparaît. Le texte doit déclarer laquelle des deux conventions il utilise ; les mélanger produit une largeur erronée d'un facteur deux.

\Needspace{7\baselineskip}
\subsection{Conjugaison de la particule et de l'antiparticule}
Soit \(C_M\) l'involution qui inverse orientation et feuillet, et \(U_C\) sa réalisation unitaire sur l'espace physique. L'égalité des masses de pôle et des largeurs exige la covariance de l'opérateur comme des canaux :

\[
U_CH_{\rm eff}^{X}(z)U_C^{-1}
=H_{\rm eff}^{\bar X}(z).
\]
Alors les ensembles de pôles retardés coïncident après transport du secteur et de la prescription causale, et

\[
M_{\bar X}=M_X,
\qquad
\Gamma_{\bar X}=\Gamma_X.
\]
La seule équivalence unitaire ne transforme pas le pôle retardé en son conjugué complexe ; elle transporte le même problème spectral vers le secteur antiparticule. La parité de la signature massique ne suffit pas à démontrer la seconde égalité : les canaux et leurs mesures doivent eux aussi être covariants. Si la symétrie est représentée antiunitairement, il faut également transporter la prescription retardée ; la seule conjugaison \(z\mapsto\bar z\) échange pôles retardés et avancés et ne constitue pas, à elle seule, l'affirmation physique de l'égalité des durées de vie moyennes.

\Needspace{7\baselineskip}
\subsection{Inventaire typé des réalisations et des observables de comparaison}
La coupe externe conservée contient

\[
471\ \text{observables de masse},
\qquad
384\ \text{observables de largeur}.
\]
La partition des masses est

\Needspace{7\baselineskip}
\begingroup
\footnotesize
\setlength{\tabcolsep}{1.9pt}
\begin{longtable}{L{5.04cm}L{5.04cm}L{5.04cm}}
\toprule
\rowcolor{HMTNavy}\HMTtablehead{Domaine physique d'arrivée} & \HMTtablehead{Masses} & \HMTtablehead{Largeurs} \\
\midrule
\endfirsthead
\multicolumn{3}{l}{\sffamily\scriptsize\color{HMTGray}Suite du tableau}\\[-1pt]
\toprule
\rowcolor{HMTNavy}\HMTtablehead{Dominio físico de llegada} & \HMTtablehead{Masas} & \HMTtablehead{Anchuras} \\
\midrule
\endhead
\midrule
\endfoot
\bottomrule
\endlastfoot
partículas elementales, leptónicas y extensiones & 53 & 9 \\
\rowcolor{HMTLight}quarks y extensiones quark & 8 & 1 \\
mesones & 243 & 225 \\
\rowcolor{HMTLight}bariones & 166 & 149 \\
registro gravitatorio de contraste & 1 & 0 \\
\rowcolor{HMTLight}\textbf{Total} & \textbf{471} & \textbf{384} \\
\end{longtable}
\endgroup
Estas filas son observables, no especies elementales. Varias filas pueden referirse a diferentes cargas de un mismo multiplete, estimadores de una misma resonancia, masas Breit--Wigner, masas de polo o estados excitados. La individualización interna precede siempre a la fila externa.

El inventario conserva además una propiedad metodológica fuerte: ninguna de sus 855 magnitudes fue autorizada como selector de ruta. Los 471 registros de masa y los 384 de anchura sólo se abren después de congelar la identidad HMT, el proyector, el operador y la carta de comparación.

\iffalse
\Needspace{7\baselineskip}
\subsection{Malla histórica extendida de pares y resonancias}
Una formulación histórica asoció a la malla angular extendida la expresión

\[
\Theta(n,k,t)
=nA_{\rm rad}-k\frac{C^*_{\rm rad}}6+t\Delta_4.
\]
Avec les définitions actuelles de \(A\), \(C^*\), \(\Delta_4\) et de la carte exponentielle, cette expression isolée ne reproduit pas les valeurs archivées. Elle est conservée comme hypothèse historique de paramétrisation, et non comme leur générateur exécutable. Nous conservons les douze évaluations importées parce qu'elles font partie de la généalogie numérique de la loi. Cinq noms --- \(\eta,\rho^0,\phi(1020),J/\psi\) et \(\Upsilon(1S)\) --- élargissent concrètement les étiquettes visibles par rapport aux dix-sept évaluations de référence :

\Needspace{7\baselineskip}
\begingroup
\fontsize{7.2}{8.2}\selectfont
\setlength{\tabcolsep}{1.9pt}
\begin{longtable}{L{1.25cm}L{0.78cm}L{0.78cm}L{0.78cm}L{3.43cm}L{3.90cm}L{4.06cm}}
\toprule
\rowcolor{HMTNavy}\HMTtablehead{État} & \HMTtablehead{\(n\)} & \HMTtablehead{\(k\)} & \HMTtablehead{\(t\)} & \HMTtablehead{Référence ultérieure (MeV/\allowbreak{}c\({}^{2}\))} & \HMTtablehead{Évaluation HMT historique (MeV/\allowbreak{}c\({}^{2}\))} & \HMTtablehead{Diferencia (ppm)} \\
\midrule
\endfirsthead
\multicolumn{7}{l}{\sffamily\scriptsize\color{HMTGray}Suite du tableau}\\[-1pt]
\toprule
\rowcolor{HMTNavy}\HMTtablehead{Estado} & \HMTtablehead{\(n\)} & \HMTtablehead{\(k\)} & \HMTtablehead{\(t\)} & \HMTtablehead{Referencia posterior (MeV/\allowbreak{}c\({}^{2}\))} & \HMTtablehead{Evaluación HMT histórica (MeV/\allowbreak{}c\({}^{2}\))} & \HMTtablehead{Diferencia (ppm)} \\
\midrule
\endhead
\midrule
\endfoot
\bottomrule
\endlastfoot
\(\mu\) & 64 & 146 & 1 & 105.658374 & 105.656529 & -17.461938\allowbreak{}2274 \\
\rowcolor{HMTLight}\(\tau\) & 64 & 349 & -1 & 1776.86 & 1776.835 & -14.069763\allowbreak{}5154 \\
\(\pi^\pm\) & 37 & 125 & 0 & 139.57039 & 139.570582 & 1.375649\allowbreak{}94982 \\
\rowcolor{HMTLight}\(K^\pm\) & 52 & 177 & 1 & 493.677 & 493.679 & 4.051231\allowbreak{}87833 \\
\(K^0\) & 53 & 178 & 1 & 497.611 & 497.616 & 10.048009\allowbreak{}3889 \\
\rowcolor{HMTLight}\(p\) & 55 & 195 & 1 & 938.272081 & 938.27244 & 0.382618\allowbreak{}226919 \\
\(n\) & 55 & 197 & 0 & 939.565413 & 939.538 & -29.176254\allowbreak{}9161 \\
\rowcolor{HMTLight}\(\eta\) & 46 & 158 & 0 & 547.862 & 547.836 & -47.457206\allowbreak{}3768 \\
\(\rho^0\) & 50 & 176 & 0 & 775.26 & 775.267 & 9.029228\allowbreak{}90385 \\
\rowcolor{HMTLight}\(\phi(1020)\) & 54 & 190 & 0 & 1019.461 & 1019.475 & 13.732747\allowbreak{}0104 \\
\(J/\allowbreak{}\psi\) & 64 & 237 & 0 & 3096.916 & 3096.979 & 20.342818\allowbreak{}4684 \\
\rowcolor{HMTLight}\(\Upsilon(1S)\) & 89 & 334 & 0 & 9460.3 & 9460.34 & 4.228195\allowbreak{}72318 \\
\end{longtable}
\endgroup
Las doce cifras archivadas son evaluaciones históricas importadas; no se atribuyen a \(\Theta\) hasta localizar la carta que las reproduzca. Además, la malla no construye el morfismo

\[
\text{état composé}
\longrightarrow\text{route et registre généalogique}
\longrightarrow(n,k,t).
\]
Elle conserve en outre la référence externe dans le même tableau historique. Sa force probante est donc celle d'une preuve numérique conditionnelle, et non d'une sélection physique antérieure à la confrontation. La présenter élargit le dénombrement visible sans transformer cinq résonances en cinq prédictions antérieures à la confrontation, ni supplanter les évaluations E3 en vigueur.
\fi

\Needspace{7\baselineskip}
\subsection{Normalisation des identités nominales en enregistrements comparables}
Le raccordement prospectif complet prend la forme

\[
\begin{aligned}
\text{identité physique}
&\longrightarrow
\text{secteur et représentation}
\longrightarrow
\text{projecteur de fibre ou d'arbre composé}\\
&\longrightarrow
\text{route enrichie}
\longrightarrow
\text{registre généalogique}
\longrightarrow
\text{signature et tours}\\
&\longrightarrow
\text{opérateur préparé}
\longrightarrow
\text{pôle ou valeur propre}
\longrightarrow
\text{carte dimensionnelle}\\
&\longrightarrow
\text{comparaison externe}.
\end{aligned}
\]
Pour un état stable, le résultat opératoriel est une valeur propre réelle. Pour une résonance, c'est un pôle complexe. Pour un quark, l'application inclut schéma et échelle. Pour une limite expérimentale, elle inclut le niveau de confiance. Pour un observable corrélé, elle inclut la matrice de covariance. Ces différences sont structurelles : elles déterminent quelles grandeurs peuvent être comparées.

\Needspace{7\baselineskip}
\subsection{Conditions de clôture individuelle des réalisations matérielles}
Une ligne de l'inventaire est individualisée à partir de HMT--MD lorsque sont construits, sans consulter sa grandeur externe :

\begin{enumerate}
\item le secteur et la représentation ;
\item le parcours, la multisection ou l'arbre de composition ;
\item l'enregistrement généalogique et la signature ;
\item l'opérateur de masse et, le cas échéant, l'opérateur des canaux ;
\item le projecteur qui sélectionne l'état nommé ;
\item la convention de masse, de largeur, de schéma et d'échelle ;
\item la carte dimensionnelle ;
\item le critère interne de réfutation.
\end{enumerate}
La loi universelle fournit déjà la grammaire commune. Le travail par état ne consiste pas à inventer une formule différente, mais à construire l'entrelaceur qui transporte une identité physique jusqu'à l'opérateur approprié.

\Needspace{7\baselineskip}
\subsection{obstructions et conditions spécifiques de non-dégénérescence}
La composition est réfutée si deux parenthésages physiquement équivalents donnent des signatures ou des spectres distincts. Le sélecteur nominal est réfuté s'il change lorsque la colonne externe de masse est masquée. Le lecteur de largeur est réfuté s'il produit \(\Gamma<0\) pour une contraction, confond amplitude et probabilité ou attribue une largeur à une section stable sans canal ouvert. La carte de résonance est réfutée si elle ne déclare pas masse de pôle ou de Breit--Wigner. La comparaison des quarks reste indéterminée si elle omet le schéma ou l'échelle. L'égalité des largeurs particule--antiparticule est réfutée si l'opérateur des canaux n'est pas covariant par conjugaison.

\Needspace{7\baselineskip}
\subsection{Transport des classes matérielles vers le corridor exceptionnel au moyen des parcours, signatures et représentations}
La grammaire de composition, les domaines mésonique et baryonique, l'opérateur de Feshbach, le dénombrement 471/384 et le principe de comparaison sont conservés dans le développement intégral antérieur, sections 16 et 18. Les enregistrements individualisés sont présentés intégralement dans les inventaires exhaustifs de masses et de largeurs de l'appendice. La correction du signe du pôle découle directement de \(\lambda=\exp(-izt_0/\hbar_{\rm int})\) et renforce la formulation sans modifier aucun de ces enregistrements.

\Needspace{7\baselineskip}
\section{Secteurs topologiques, extensions hypothétiques et gravitation}
\Needspace{7\baselineskip}
\subsection{Séparation typée entre secteurs topologiques, extensions hypothétiques et réalisation gravitationnelle}
L'architecture distingue précisément :

\begin{enumerate}
\item une identité algébrique interne ;
\item un codomaine de fusion ou de tressage ;
\item une particule physique avec Hamiltonien, gap, localisation et canal de
mesure.

\end{enumerate}
Une règle de fusion exacte n'est pas encore une masse. Un mode de bord n'est pas encore une espèce. Une involution de parcours n'est pas encore une condition de Majorana. Le passage physique exige un foncteur qui préserve produit, involution, orientation, énergie et composition des chemins.

\Needspace{7\baselineskip}
\subsection{Secteur de Fibonacci : anneau basé, fusion et représentation de tresses}
L'incidence surfacique--volumétrique est régie par

\[
F_{\rm av}=
\begin{pmatrix}
0&1\\[2pt]1&1
\end{pmatrix}.
\]
Son polynôme caractéristique est

\[
x^2-x-1,
\]
et sa valeur propre dominante est \(\varphi\). Dans la base \(\{\mathbf1,\tau\}\), la multiplication par \(\tau\) est représentée par \(F_{\rm av}\) ; ainsi

\[
\tau\otimes\tau\cong\mathbf1\oplus\tau,
\qquad
\operatorname{FPdim}(\tau)=\varphi.
\]
C'est une réalisation exacte de l'anneau basé de Fibonacci. Pour obtenir une catégorie anyonique physique, il faut en outre construire l'associateur \(F\), le tressage \(R\), les équations du pentagone et de l'hexagone, la chiralité, un Hamiltonien local à gap et l'application de préparation et de lecture. Le résultat de masse, lorsqu'il existe, est un gap d'énergie converti au moyen d'une vitesse effective :

\[
m_{\rm brecha}=\frac{\Delta E}{c_{\rm eff}^{2}}.
\]
Il n'existe pas de masse universelle de l'anyon de Fibonacci déductible uniquement de \(\varphi\).

\Needspace{7\baselineskip}
\subsection{Séparation entre le secteur topologique et le centre pointé de l'atlas}
Le centre tressé du secteur pointé fondé sur \(\mathbb Z/9\mathbb Z\) possède 81 objets simples inversibles, chacun de dimension quantique un. Un objet simple de Fibonacci exigerait la dimension \(\varphi\), et un objet simple d'Ising exigerait \(\sqrt2\). Aucun des deux n'apparaît donc comme objet simple interne de ce centre pointé.

Ce contrôle n'élimine pas les réalisations topologiques. Il détermine leur place : elles doivent s'obtenir par extension, condensation, quotient, défaut ou foncteur monoïdal vers un codomaine non pointé. Présenter l'anneau de Fibonacci comme s'il constituait déjà toute la catégorie effacerait précisément le morphisme physique à construire.

\Needspace{7\baselineskip}
\subsection{Secteur d'Ising et modes de Majorana}
La catégorie d'Ising a pour objets simples

\[
\mathbf1,\quad\psi,\quad\sigma
\]
et règles

\[
\psi\otimes\psi=\mathbf1,
\qquad
\psi\otimes\sigma=\sigma,
\qquad
\sigma\otimes\sigma=\mathbf1\oplus\psi,
\qquad
d_\sigma=\sqrt2.
\]
Quatre objets distincts reçoivent souvent des noms proches et doivent demeurer séparés :

\begin{itemize}
\item l'involution des parcours \(C_M=-\operatorname{rev}\) ;
\item la parité de la complétion CAR ;
\item un spineur relativiste autoconjugué ;
\item un mode zéro topologique de Majorana dans une phase d'Ising.
\end{itemize}
Un fermion relativiste est de Majorana lorsqu'il existe une conjugaison réelle \(\mathcal C\) compatible avec sa représentation et sa dynamique,

\[
\psi=\mathcal C\bar\psi^{T},
\]
et que le terme de masse respecte cette condition. Un mode zéro de Majorana se décrit par un opérateur autoadjoint \(\gamma=\gamma^\dagger\) qui satisfait

\[
\{\gamma_i,\gamma_j\}=2\delta_{ij}
\]
et reste séparé du continu par un gap. L'égalité d'un mot avec son image par \(C_M\) ne démontre aucune de ces deux réalisations.

Pour les neutrinos, l'alternative Dirac/Majorana ne se décide qu'après construction de l'opérateur neutre, de sa base de masses, de l'entrelaceur de mélange et de la structure réelle fermionique. La neutralité électrique est nécessaire à une autoconjugaison non triviale, mais ne suffit pas.

\Needspace{7\baselineskip}
\subsection{Tressage abélien d'ordre \(6\)}
Soit

\[
\operatorname{wind}_6:B_n\longrightarrow\mathbb Z/6\mathbb Z
\]
le lecteur d'enroulement et

\[
\chi_k(b)=\exp\!\left(
\frac{2\pi i k}{6}\operatorname{wind}_6(b)\right).
\]
Les réalisations de phase \(q\) se réduisent à de simples permutations seulement lorsque

\[
q^2=1.
\]
Les quatre caractères restants distinguent des classes algébriques du lecteur d'enroulement. Le résultat exact est la classification des six caractères unidimensionnels du quotient d'enroulement ; ce n'est pas encore une classification de six phases matérielles. Affirmer qu'ils conservent un tressage physique exige des opérateurs de tresse et une représentation du groupe de tresses. Leur existence n'implique pas quatre nouvelles particules : une espèce anyonique exige un Hamiltonien, des opérateurs de tresse, un gap, une stabilité et un couplage à une multisection physique. Leur promotion au statut de phases matérielles dépend de cette réalisation.

\Needspace{7\baselineskip}
\subsection{Positivité et exclusion du tachyon persistant}
La droite d'action possède une orientation positive,

\[
\hbar_{\rm int}>0,
\]
le caractère de masse prend des valeurs positives,

\[
\chi_{\rm full}(\sigma,\eta)>0,
\]
et le raffinement bidirectionnel est réalisé par conjugaison. Supposons explicitement

\[
\widehat M^{(0)}\ge0,
\qquad
\widehat K=\widehat K^*,
\qquad
R_{\rm act}>0.
\]
Alors

\[
\widehat M^{\beta}
=R_{\rm act}^{\widehat K/2}
\widehat M^{(0)}
R_{\rm act}^{\widehat K/2}\ge0.
\]
En conséquence,

\[
\operatorname{Spec}(\widehat M^\beta)\subset[0,\infty).
\]
Ainsi, l'opérateur de masse au carré

\[
\widehat{M^2}_{\beta}:=(\widehat M^\beta)^2
\]
est positif et satisfait

\[
\operatorname{Spec}(\widehat{M^2}_{\beta})\subset[0,\infty).
\]
Une particule asymptotique persistante satisfaisant les prémisses de positivité et d'unitarité ne peut avoir \(m^2<0\). Si la linéarisation d'un fond produit une fréquence imaginaire, l'objet est une direction instable : le fond doit se relaxer ou se condenser. Si le parcours cesse de se prolonger, l'objet est une section virtuelle à obstruction terminale. Aucun des deux cas ne définit une particule stable supraluminique.

La portée de la proposition est fixée par ses prémisses. Un paramètre « tachyonique » dans un potentiel perturbatif peut signaler une transition de phase ; ce qui est exclu est sa promotion au statut d'espèce persistante de norme positive.

\Needspace{7\baselineskip}
\subsection{Exclusion indépendante des fantômes}
Un fantôme physique serait un mode de norme négative. Soit \(\mathcal H_{\rm phys}=\ker G/\operatorname{im}G_0\) le quotient de jauge muni d'un produit scalaire positif, et \(\widehat M_{\rm phys}\) autoadjoint sur ce quotient. Alors la mesure spectrale d'un état préparé satisfait

\[
\mu_{P,a}^{\rm mass}(B)\ge0.
\]
Une direction négative indique que la positivité du quotient a échoué ou que le mode n'appartient pas à l'espace physique ; elle ne disparaît pas par une affirmation nominale. Les fantômes de Faddeev--Popov sont des variables du complexe de jauge, et non des particules asymptotiques. Cette exclusion est distincte de l'exclusion tachyonique : un fantôme viole la positivité de la norme ; une direction tachyonique viole la stabilité du fond.

\Needspace{7\baselineskip}
\subsection{Photon, biais unidirectionnel et lien Proca--Stückelberg}
Les canaux harmoniques produisent le biais

\[
\varepsilon=q_-^{90}-q_+^{90}.
\]
Une réalisation cinématique bidirectionnelle peut s'écrire

\[
c_+=\frac{c_{\rm int}}{1+\varepsilon},
\qquad
c_-=\frac{c_{\rm int}}{1-\varepsilon},
\]
pour laquelle la mesure d'aller et retour conserve exactement

\[
\frac{2}{c_+^{-1}+c_-^{-1}}=c_{\rm int}.
\]
Le biais demeure dans la lecture unidirectionnelle et s'annule dans la lecture réciproque. Cette identité appartient à l'architecture à deux orientations ; elle ne démontre pas, à elle seule, une masse photonique.

Si la clôture de frontière produit une fréquence de gap \(\mu_T\in\mathcal L_T^{-1}\) dans l'enregistrement,

\[
\omega^2=c_{\rm int}^{2}k^2+\mu_T^2,
\]
un entrelaceur dimensionnel peut la transporter vers une réalisation Proca--Stückelberg de masse effective

\[
m_\gamma=\frac{\hbar_{\rm int}\mu_T}{c_{\rm int}^{2}}.
\]
Pour une propagation de fréquence \(\omega\), la correction temporelle du premier ordre satisfait

\[
\frac{\Delta T}{T}
\simeq
\frac12
\left(
\frac{m_\gamma c_{\rm int}^{2}}
     {\hbar_{\rm int}\omega}
\right)^2.
\]
La chaîne requise est

\[
\text{aller/retour TPK}
\longrightarrow
\text{mémoire}
\longrightarrow
\mu_T
\longrightarrow
m_\gamma.
\]
On n'identifie pas automatiquement \(\varepsilon\), le quotient d'action \(R_{\rm act}\) et \(\mu_T\) : ils appartiennent à des codomaines distincts. La carte nulle du photon ordinaire demeure exacte tant que le gap et son entrelaceur physique ne sont pas construits.

\Needspace{7\baselineskip}
\subsection{Classification opératorielle des secteurs sombres par domaine, lecteur et codomaine physique}
Un secteur sombre ne se définit pas par une masse arbitraire, mais par le noyau de ses couplages. Si \(P_X\) est son projecteur,

\[
P_{\rm em}P_X=0,
\qquad
P_{\rm color}P_X=0,
\]
tandis qu'il peut se produire

\[
P_{\rm grav}P_X\ne0
\]
ou exister un couplage faible. La loi fournit d'abord un opérateur de fibre ; le canal détermine ensuite si son spectre contient un mode stable, métastable ou résonant.

Un photon sombre exige une seconde connexion \(A'\), sa courbure \(F'\) et un entrelaceur cinétique

\[
\mathcal L_{\rm mix}
=-\frac{\varepsilon_{\rm kin}}2F_{\mu\nu}F'^{\mu\nu}.
\]
Si le second secteur acquiert une masse par Stückelberg, Higgs caché ou frontière, son pôle est distinct de celui du photon nul. La grandeur \(\varepsilon_{\rm kin}\) doit être dérivée d'une incidence commune avant la comparaison des observables.

Soient \(P_{\rm weak},P_{\rm em},P_{\rm color}\) les projecteurs sur les canaux de représentation faible, électromagnétique et chromatique, et \(P_X\) le support de la section sombre. Un WIMP exige un recouvrement faible non nul et l'absence de support électromagnétique et chromatique :

\[
P_{\rm weak}P_X\ne0,
\qquad
P_{\rm em}P_X=P_{\rm color}P_X=0.
\]
Le premier produit exprime une incidence de sous-espaces ; un commutateur non nul ne constituerait pas un critère de couplage, puisque deux projecteurs peuvent commuter tout en partageant un support. L'intensité de l'interaction exige ensuite un opérateur de vertex ou un Hamiltonien de couplage.

Sa stabilité exige une valeur propre unitaire du transport ou une parité protectrice. « WIMP » est une condition de représentation et de persistance, et non une masse prédéterminée.

\Needspace{7\baselineskip}
\subsection{Classification de la section neutrinique stérile par fibre invariante, mélange et charge de jauge nulle}
Un neutrino stérile est un vecteur propre du bloc neutre qui satisfait

\[
P_{\rm weak}n_s=0,
\qquad
P_{\rm mass}n_s=n_s.
\]
Il n'est pas synonyme de la quatrième profondeur structurelle du support neutrinique. Son existence exige un noyau faible non nul, une valeur propre de l'opérateur de masse dans ce noyau et un transport de mélange avec les modes actifs. Le critère est opératoriel et peut être réfuté avant l'introduction d'une échelle externe.

\Needspace{7\baselineskip}
\subsection{Classification des sections axionique et pseudoscalaires par caractère topologique et couplage de bord}
Un candidat axionique doit être un mode pseudoscalaire à terme cinétique normalisé,

\[
CP\,a\,CP^{-1}=-a,
\qquad
\mathcal L_{\rm kin}=\frac{Z_a}{2}\,\partial_\mu a\,\partial^\mu a,
\qquad Z_a>0,
\]
couplé à la densité topologique de couleur. Sa masse procède de la courbure du potentiel effectif,

\[
m_a^2=\frac1{Z_a}
\left.
\frac{\partial^2V_{\rm eff}}{\partial a^2}
\right|_{a=a_0}.
\]
Le canal \(C^*\mapsto-C^*\) fournit une orientation impaire, mais ne suffit pas à construire \(V_{\rm eff}\), le terme cinétique ou le projecteur physique. L'espèce n'est définie que lorsque ces applications sont présentes.

\Needspace{7\baselineskip}
\subsection{Réalisations monopolaires et défauts topologiques : classes de charge et stabilité}
Soit \(F\) la courbure d'une connexion \(U(1)\) dans la normalisation intégrale choisie. Un monopôle magnétique se caractérise par la classe

\[
\frac1{2\pi}\int_{S^2}F=n\ne0.
\]
HMT doit réaliser \(n\) comme holonomie nette d'une famille fermée de parcours et démontrer qu'elle survit au raffinement. Une lacune ponctuelle ou une retenue isolée ne constitue pas un monopôle. Si le défaut est dynamique et stable, sa masse est une valeur propre de l'opérateur linéarisé autour de cette classe.

\Needspace{7\baselineskip}
\subsection{Classification des extensions supersymétriques par appariement des représentations et des caractères}
La coexistence de complétions fermionique et bosonique ne démontre pas la supersymétrie. En quatre dimensions, il faut une famille de charges impaires satisfaisant

\[
Q_\alpha:\mathcal F_-\longleftrightarrow\mathcal F_+,
\qquad
\{Q_\alpha,\bar Q_{\dot\beta}\}
=2\sigma^\mu_{\alpha\dot\beta}P_\mu,
\]
qui entrelace les opérateurs de masse :

\[
Q_\alpha\widehat M_-=\widehat M_+Q_\alpha.
\]
L'identité \(Q^2=H\) est réservée à une réduction hermitienne unidimensionnelle explicite ; elle ne remplace pas l'algèbre relativiste précédente.

Ce n'est qu'alors qu'apparaissent des paires spectrales. La brisure doit dériver la séparation des masses sans utiliser la valeur du partenaire recherché. Sélectrons, neutralinos, gluinos ou autres noms demeurent des réalisations possibles, et non des lignes prédites par la seule existence de CAR et de CCR.

\Needspace{7\baselineskip}
\subsection{Classification des sections de dilaton et de radion par modes scalaires de volume et de compactification}
Le dilaton est une fluctuation de la section d'échelle ; le radion, d'une longueur ou d'un volume interne. Leurs actions élémentaires sont

\[
t_0\mapsto e^{\phi_d}t_0,
\qquad
\mathcal V_{\rm int}\mapsto e^{\phi_r}\mathcal V_{\rm int}.
\]
Pour constituer des particules, ils nécessitent un terme cinétique, un potentiel, un projecteur persistant et un pôle. Auparavant, ce sont des degrés collectifs de la carte dimensionnelle. Cette distinction évite de transformer une liberté de normalisation en nouvelle espèce scalaire.

\Needspace{7\baselineskip}
\subsection{Composés et énergie de liaison}
Un état composé \(Y\) de constituants \(X_1,\ldots,X_k\) se construit dans l'ordre

\[
(\Gamma_{X_1},\ldots,\Gamma_{X_k})
\longrightarrow
\mathscr L_{\rm bind}
\longrightarrow
\sigma_Y
\longrightarrow
\eta_Y
\longrightarrow
\widehat M_Y.
\]
L'enregistrement de liaison \(\mathscr L_{\rm bind}\) contient l'ordre des parcours, la frontière commune, la retenue transférée, le canal de couleur et l'holonomie de liaison. Il ne peut être remplacé par \(\sum_i m_{X_i}\). Mésons et baryons exigent des lecteurs distincts parce que leurs nombres de constituants, leurs représentations et leurs symétries d'échange diffèrent.

Les six évaluations composées publiées prouvent l'évaluation du caractère une fois la signature fixée. La clôture état par état exige en outre que l'enregistrement de liaison régénère cette signature. Cette condition s'applique de la même manière aux noyaux, atomes, molécules et candidats tétraquarks ou pentaquarks, avec leurs Hamiltoniens de frontière respectifs.

\Needspace{7\baselineskip}
\subsection{Largeur comme 2e observable}
Les 384 largeurs ne s'obtiennent pas en ajoutant une sixième coordonnée à \(\sigma\). Pour chaque état instable, on construit les canaux \(\mathcal H_{\rm in}\oplus\mathcal H_{\rm out}\), l'opérateur effectif \(H_{\rm eff}(z)\) et le pôle

\[
z_X=M_Xc_{\rm int}^{2}-\frac{i}{2}\Gamma_X.
\]
La partie réelle relève de la loi de masse ; la partie imaginaire, du lecteur de désintégration. La relation \(\tau_X=\hbar_{\rm int}/\Gamma_X\) s'interprète après fixation de la carte temporelle. Une théorie exhaustive de 471 masses et de 384 largeurs contient donc deux lecteurs coordonnés, et non une seule colonne de 855 nombres hétérogènes.

\Needspace{7\baselineskip}

\endgroup
\begingroup
\let\chapter\section
\let\section\subsection
\let\subsection\subsubsection
\let\subsubsection\paragraph
\let\clearpage\relax
% La subsección inicial de gluones coincide byte por byte con la residencia
% REV.2 anterior; la continuación conjunta fotón--gluón se conserva íntegra.
% Extracto editor-ready del contenido exclusivo de masas/16_hidrodinamica_y_gauge.tex.
% Las líneas 1--14 están publicadas byte-exactamente desde
% 74_rev2_realizaciones_i.tex:544--557. El propietario íntegro permanece
% preservado; este extracto conserva la continuación exclusiva.

\subsection{Comparaison des réalisations nulles électromagnétique et chromodynamique}
Le photon et les huit gluons appartiennent à des cartes nulles distinctes. Pour le photon,

\[
m_\gamma=0
\]
dans la branche électromagnétique. Pour les gluons,

\[
m_g=0
\]
dans la représentation adjointe de couleur. La nullité ne s'obtient pas en prenant la limite d'un caractère positif ; le domaine bifurque auparavant. La réalisation photonique exige Maxwell, l'invariance de jauge et deux polarisations transversales. La réalisation gluonique exige la connexion de couleur, ses huit générateurs et la restriction à l'espace physique de jauge.

La propagation bidirectionnelle peut s'écrire

\[
c_\pm=\frac{c}{1\pm\epsilon},
\qquad
\epsilon=q_-^{90}-q_+^{90},
\]
avec pour moyenne harmonique

\[
\frac{2}{1/c_++1/c_-}=c.
\]
La vitesse bidirectionnelle conserve \(c\) ; l'orientation unidirectionnelle conserve le biais. Un gap supplémentaire de l'enregistrement généalogique peut se réaliser par Proca--Stückelberg, mais ne doit pas être identifié automatiquement à \(\epsilon\) ni au dédoublement de l'action.

\Needspace{7\baselineskip}

\endgroup
\begingroup
\let\chapter\section
\let\section\subsection
\let\subsection\subsubsection
\let\subsubsection\paragraph
\let\clearpage\relax
\section{Chaîne exceptionnelle, théorie conforme, cordes et classification des masses}
\Needspace{7\baselineskip}
\subsection{Antécédent commun et coordination des \(2\) lectures de la chaîne exceptionnelle}
La relation entre masses et symétries exceptionnelles ne commence pas par une action du Monstre sur une table de particules. Elle commence dans le même état enrichi du TPK. Soit

\[
x\in\mathcal X_{\rm TPK}^{\rm enr}.
\]
La lecture archimédienne contracte des cylindres compatibles et conserve la valeur réelle avec sa généalogie :

\[
\mathcal P_{\rm arq}(x)\in\mathbb R.
\]
La lecture exceptionnelle conserve les incidences de frontière et produit une classe dans la limite projective :

\[
\mathcal P_{\rm exc}(x)=Q_\infty(x).
\]
La naturalité exige, pour tout raffinement \(n\ge m\),

\[
r_{n,m}\circ Q_n=Q_m\circ p_{n,m}.
\]
L’application conjointe

\[
\mathcal P_{\rm joint}:x\longmapsto
\bigl(\mathcal P_{\rm arq}(x),\mathcal P_{\rm exc}(x)\bigr)
\]
est la forme positive de la double projection. Elle ne nie pas la connexion Monstre--masse : elle la type. La branche de masse et la branche exceptionnelle sont les images d’un même antécédent doté d’une orientation, d’une graduation, d’une retenue, d’une mémoire et d’une holonomie communes. Elle n’affirme pas non plus que le Monstre agit littéralement sur les chiffres décimaux ou qu’un coefficient de Moonshine est à lui seul une masse.

\Needspace{7\baselineskip}
\subsection{Construction de la chaîne exceptionnelle}
L’incidence ternaire prolonge l’état par la chaîne

\[
C_3
\longrightarrow A_2^{12}
\longrightarrow N(A_2^{12})
\longrightarrow \Lambda_{24}
\longrightarrow V^\natural
\longrightarrow \mathbb M,
\]
où \(\Lambda_{24}\) est le réseau de Leech, \(V^\natural\) l’algèbre de vertex du Moonshine et \(\mathbb M\) le groupe Monstre. Le passage ne consiste pas à juxtaposer des noms classiques : chaque flèche doit préserver l’incidence, la graduation et la compatibilité de raffinement héritées de l’état.

La trace ternaire s’écrit

\[
t_3(q)=q^{-1}
\prod_{n\ge1}(1+q^n+q^{2n})^{-12}.
\]
Le premier défaut non trivial vérifie

\[
D_{\rm APP}=[q]t_3=v_\alpha^2=54,
\]
et le premier coefficient du caractère du Monstre conserve la décomposition

\[
196884=54+10\cdot3^9.
\]
Ces identités transportent la graduation et la multiplicité. Elles n’attribuent pas de masse en MeV. Leur fonction dans la théorie des masses est de classifier l’espace sur lequel agit l’opérateur avant de le scalariser.

\Needspace{7\baselineskip}
\subsection{Appariement fibre à fibre avec la loi des masses}
La branche matérielle du même état produit

\[
x\longmapsto
\bigl(F_X,\sigma_X,\eta_X,\widehat M_X\bigr),
\]
où \(F_X\) est un point, une orbite, une fibre ou une multisection ; \(\sigma_X\in\mathbb Z^5\) est la signature ; \(\eta_X\in\mathcal E_{90,120}\) est le raffinement ; et \(\widehat M_X\) l’opérateur de masse. La lecture conjointe est alors

\[
\mathfrak J_{\rm mass,exc}(x)
=\bigl(F_X,\sigma_X,\eta_X,\widehat M_X,Q_\infty(x)\bigr).
\]
Deux états ayant la même valeur scalaire de masse peuvent rester distincts par leur composante exceptionnelle ; deux états ayant le même caractère exceptionnel peuvent posséder des opérateurs de masse différents. La classification correcte est l’appariement des fibres, et non l’égalité isolée des nombres.

Si une représentation exceptionnelle \(G\) agit dans une fibre physique \(\mathcal H_X\), la compatibilité avec la masse exige

\[
U_g\widehat M_XU_g^{-1}=\widehat M_X,
\qquad g\in G.
\]
La décomposition

\[
\mathcal H_X=\bigoplus_{\lambda}
\mathcal H_{X,\lambda}
\]
sépare les multiplets exceptionnels. Sur un facteur direct irréductible, la commutation impose que la compression de \(\widehat M_X\) soit scalaire. C’est une voie légitime de scalarisation par symétrie. La valeur du scalaire continue de provenir du caractère, des tours, de la section dimensionnelle et du couplage ; la dimension de la représentation ne détermine que la dégénérescence.

\Needspace{7\baselineskip}
\subsection{Fock et algèbre de vertex}
Le prolongement d’un parcours à un espace à plusieurs particules possède deux complétions principales. La complétion extérieure produit les CAR, Pauli et la statistique fermionique ; la complétion symétrique produit les CCR et la statistique bosonique. Une algèbre de vertex ajoute un produit d’opérateurs, un vide, une translation et une graduation :

\[
V=\bigoplus_{n\ge0}V_n,
\qquad
L_0|_{V_n}=nI.
\]
Le caractère

\[
\operatorname{ch}_V(q)
=\operatorname{Tr}_V q^{L_0-c/24}
\]
compte les états par degré. Une multiplicité dans \(V_n\) n’est pas une masse. Pour transformer la graduation en énergie, il faut un Hamiltonien et une section de longueur ou de temps. Dans une réalisation conforme sur un cylindre de circonférence \(L\),

\[
H_{\rm CFT}
=\frac{2\pi\hbar_{\rm int}v}{L}
\left(L_0+\bar L_0-\frac{c}{12}\right).
\]
La masse ou le gap physique n’apparaît qu’après fixation de \(v\), de \(L\), des conditions de bord et de l’application du parcours TPK vers les modules conformes. La fonction de la chaîne exceptionnelle consiste à organiser les dégénérescences et les règles de composition de ces modules.

\Needspace{7\baselineskip}
\subsection{Spectre de corde comme réalisation dimensionnelle ultérieure}
Soit \(\alpha'\in\mathcal L_L^2\) la longueur quadratique représentant la tension inverse dans la carte naturalisée de corde déclarée. Dans une compactification circulaire de rayon \(R\), les impulsions gauche et droite sont

\[
p_L=\frac{m}{R}+\frac{wR}{\alpha'},
\qquad
p_R=\frac{m}{R}-\frac{wR}{\alpha'},
\]
avec \(m,w\in\mathbb Z\). Dans les unités naturelles de la réalisation, le spectre fermé vérifie

\[
M^2=p_L^2+\frac{4}{\alpha'}(N_L-a_L)
=p_R^2+\frac{4}{\alpha'}(N_R-a_R),
\]
ainsi que la condition de niveau

\[
N_L-N_R=a_L-a_R-mw.
\]
L'involution

\[
(m,w,R)\longleftrightarrow
\left(w,m,\frac{\alpha'}{R}\right)
\]
conserve le spectre : c’est la réalisation de la dualité T. HMT apporte au préalable l’orientation bidirectionnelle, la mémoire et l’appariement des feuillets ; la théorie des cordes ajoute la tension, le rayon, les oscillateurs, le niveau et les conditions aux limites. On n’obtient pas une masse de corde en remplaçant \(N_L,N_R\) par des hauteurs de la tour \(\langle90,120\rangle\). Toute identification entre les deux graduations exige un opérateur d’entrelacement qui conserve la somme, l’orientation et le caractère.

\Needspace{7\baselineskip}
\subsection{Écrans réticulaires de théorie M : graduation, compactification et transport de la mémoire résiduelle}
Les écrans internes

\[
12\longrightarrow11\longrightarrow10,
\qquad
26\longrightarrow24,
\]
conservent la descente dodécaphasique, l’élimination d’une direction redondante et l’écran transversal. Dans une réalisation de théorie M, les champs

\[
g_{MN},\qquad C_{MNP},\qquad\Psi_M
\]
doivent être obtenus au moyen d’un opérateur d’entrelacement préservant la jauge, la supersymétrie, les contraintes et le produit scalaire. Les écrans numériques ne remplacent pas cette application.

La classification des masses est transportée vers cette branche au moyen de

\[
\text{route}
\longrightarrow\text{registre généalogique}
\longrightarrow\text{signature}
\longrightarrow\text{caractère}
\longrightarrow\text{opérateur}
\longrightarrow\text{spectre compactifié}.
\]
Les modes de Kaluza--Klein, d’enroulement, de membrane ou d’excitation oscillatoire sont des états distincts même s’ils partagent une masse. Le graviton d’une réalisation de corde fermée, s’il apparaît, est un mode collectif sans masse du secteur géométrique ; il ne devient pas pour autant une graine élémentaire du TPK et n’invalide pas la formulation de la gravité en termes d’holonomie et de torsion.

\Needspace{7\baselineskip}
\subsection{Réalisations ordinaires et hypothétiques induites par la chaîne exceptionnelle}
La branche exceptionnelle permet de distinguer quatre sources de multiplicité :

\begin{enumerate}
\item multiplicité cellulaire ou de parcours dans l’atlas ;
\item dégénérescence d’une représentation exceptionnelle ;
\item occupation de Fock ;
\item niveau conforme, oscillatoire ou compactifié.
\end{enumerate}
Les confondre produit de faux partenaires, de fausses générations ou des masses dupliquées. Un état jumeau, supersymétrique, sombre ou exotique ne constitue une nouvelle espèce que lorsque son quintuplet

\[
(F_X,Q_\infty(x),\rho_X,\widehat M_X,P_X)
\]
diffère physiquement et que l’Hamiltonien possède un pôle ou une section persistante supplémentaire. Une coïncidence de dimension, de caractère ou de masse ne suffit pas.

Inversement, une dégénérescence protégée par la chaîne exceptionnelle est une prédiction structurelle : les perturbations préservant la symétrie ne peuvent scinder le multiplet. Une scission observée exige l’identification de la brisure de symétrie, du couplage ou de la carte qui la produit.

\Needspace{7\baselineskip}
\subsection{Conditions de positivité, de clôture et de stabilité des réalisations topologiques}
La connexion entre la chaîne exceptionnelle et les masses est réfutée dans une réalisation concrète si l’un des faits suivants se produit :

\begin{itemize}
\item la naturalité \(r_{n,m}Q_n=Q_mp_{n,m}\) échoue ;
\item l’étiquette exceptionnelle change lorsqu’on permute les masses externes ;
\item une multiplicité de Moonshine est utilisée comme valeur de masse sans Hamiltonien ;
\item l’opérateur de masse ne commute pas avec la symétrie déclarée préservée ;
\item une hauteur de tour est identifiée à un niveau de corde sans opérateur d’entrelacement ;
\item la condition de niveau ou l’invariance par dualité T échoue ;
\item un écran dimensionnel est déclaré théorie M sans construction de ses champs ;
\item l’existence d’un mode de spin deux est utilisée pour redéfinir la gravité interne antérieure.

\end{itemize}
Ces contrôles maintiennent la connexion positive : il existe un état commun et une classification conjointe ; ce qui est interdit, c’est de remplacer les flèches par une analogie numérique.

\Needspace{7\baselineskip}
\subsection{Construction mathématique des classes matérielles à partir des parcours, des signatures et des représentations}
La double réalisation, la naturalité projective, la chaîne \(A_2^{12}\)--Leech--\(V^\natural\)--Monstre, la trace \(t_3\), les identités \(D_{\rm APP}=54\) et \(196884=54+10\cdot3^9\), ainsi que les écrans dimensionnels appartiennent à l’architecture mathématique APP--TRIT--TPK développée précédemment. La formulation de \(\mathfrak J_{\rm mass,exc}\), la séparation des quatre multiplicités et les contrôles précédents explicitent leur fonction au sein d’une théorie autonome des masses. Les formules conformes et de corde sont des réalisations dimensionnelles ultérieures, avec leurs paramètres déclarés ; elles ne sont pas employées pour engendrer rétroactivement APP, TRIT, TPK, le caractère ou les tours.

\endgroup

\section{Résultats opératoriels en vigueur et comparaison métrologique ultérieure}
La comparaison ne s’ouvre qu’après fixation du parcours, de la signature, de l’opérateur, du schéma, de l’échelle et de l’unité. Le tableau en vigueur dispose en colonnes adjacentes la sortie HMT, la valeur externe, l’incertitude, le résidu et le statut. Les valeurs des quarks n’admettent pas de résidu publiable si le schéma et l’échelle ne coïncident pas.

\begingroup
\let\chapter\section
\let\section\subsection
\let\subsection\subsubsection
\let\subsubsection\paragraph
\let\clearpage\relax
\section{Opérateurs de masse dans HMT--MD : échelle absolue, atlas des parcours et réalisation par secteurs}
\label{sec:tabla-completa-masas-hmt-sm-revfinal}

La masse apparaît en Mécanique Dimensionnelle comme la réalisation de deux branches
mathématiques qui convergent sans se confondre. La première conserve la généalogie de
l'état : une extension survivante détermine un parcours observable ; le parcours écrit
un enregistrement chronologique dodécaphasique ; l'enregistrement produit une signature entière, et la signature
alimente le caractère et la hiérarchie harmonique. La seconde construit la droite
dimensionnelle : le lecteur déterminantal fixe l'échelle d'action, l'horloge interne la
transporte vers l'énergie et la carte cinématique la mène à la masse. L'opérateur d'espèce
agit alors par homothéties positives sur cette droite. La coordonnée expérimentale
n'est introduite qu'après fixation des deux branches.

Le domaine interne complet possède quatre dénombrements distincts :
\begin{equation}
 81=25_{\ell^-}+20_{\nu}+20_d+16_u,\qquad
 56=41\!\times\!1+10\!\times\!2+2\!\times\!3+1\!\times\!4+2\!\times\!5,
 \label{eq:censos-81-56-masas-rectoras}
\end{equation}
\begin{equation}
 13\ \text{multisections},\qquad
 324=81\times4\ \text{routes orientées}.
 \label{eq:censos-13-324-masas-rectoras}
\end{equation}
Les réalisations physiques comparées plus loin constituent une
classification de réalisations ; elles ne sont pas un autre nom pour les 81 cellules, les 56
familles, les 13 multisections ou les 324 histoires orientées. Sur ce domaine,
les deux lecteurs de parcours produisent quatorze profils \(K_D\), neuf profils
\(K_\Omega\), quarante paires conjointes et dix-huit coïncidences. Le secteur nul
demeure distinct du caractère positif de masse.

\subsection{De l'incompatibilité APP à la signature de masse}

La différence primordiale entre les feuillets somme et produit apparaît dans la
décomposition spectrale du bloc central d'APP :
\begin{equation}
 B_c=7I+2R=9P_++5P_-,\qquad
 \operatorname{Spec}(B_c)=\{9,5\},\qquad
 \Delta_{\rm APP}=9-5=4.
 \label{eq:salto-app-cuatro-ley-masas}
\end{equation}
Les entiers qui réapparaissent ensuite conservent des sources distinctes. Le coefficient
local d'échange est \(2\) dans \(B_c=7I+2R\) ; la compression modulaire donne
\(51-45=6\), et son déploiement sur neuf blocs produit \(9\cdot6=54\). Ces
égalités appartiennent à des niveaux typés d'une même généalogie et ne constituent pas
une chaîne de définitions du saut \(4\). Dans la section électronique, \(54\) est
en outre \(D_1(\Gamma_e)\) et la périodicité cinématique d'ordre \(54\) de
l'endomorphisme temporel CT108 : il mesure une composante de l'histoire déjà transportée
par TPK, et non la séparation spectrale primordiale.

Pour chaque parcours enrichi \(\Gamma\), l'enregistrement dodécaphasique détermine
\begin{equation}
 \sigma(\Gamma)
 =(n_A,s_C,k_{\Delta_4},b_{120},b_\zeta)\in\mathbb Z^5,
 \qquad \zeta_{270}=135\Delta_4^2.
 \label{eq:firma-z5-masas-rectoras}
\end{equation}
Le caractère central et son raffinement harmonique sont
\begin{align}
 \log\chi_0(\sigma)
 &=n_A A_{\rm rad}+\frac{s_C}{6}C^*_{\rm rad}
   +k_{\Delta_4}\Delta_4+b_{120}s_{120}+b_\zeta\zeta_{270},
 \label{eq:caracter-central-masas-rectoras}\\
 \log\chi_{\rm full}(\sigma,\eta)
 &=n_A A_{\rm rad}+\frac{s_C}{6}C^*_{\rm rad}
   +k_{\Delta_4}\Delta_4+b_\zeta\zeta_{270}
   +(b_{120}+\eta_{120})s_{120}
   +\sum_{\substack{h\in\mathfrak S\\h\ne120}}\eta_hs_h.
 \label{eq:caracter-completo-masas-rectoras}
\end{align}
La division \(s_C/6\) procède de l'intégralité dodécaphasique. La paire étiquetée
\((b_{120},\eta_{120})\) conserve sa généalogie même lorsque son évaluation est
regroupée en un seul coefficient.

Les moments qui interviennent dans
\eqref{eq:caracter-completo-masas-rectoras} appartiennent à une hiérarchie infinie :
\begin{equation}
 q_\pm=\exp\!\left[-\frac{\pi}{180}(A\pm C^*)\right],
 \qquad s_h=q_-^h-q_+^h,\qquad c_h=q_-^h+q_+^h,
 \label{eq:momentos-torres-masas-rectoras}
\end{equation}
\begin{equation}
 \mathfrak S=\langle90,120\rangle
 =\{90,120\}\cup\{30r:r\ge6\}.
 \label{eq:semigrupo-global-torres-masas-rectoras}
\end{equation}
\(s_n\) comme \(c_n\) satisfont
\begin{equation}
 X_{n+2}=c_1X_{n+1}-(q_-q_+)X_n.
 \label{eq:recurrencia-torres-masas-rectoras}
\end{equation}
Les hauteurs \(90,120,180,210,240,270,360,420\) sont des témoins généalogiques
particulièrement visibles ; le domaine de la loi est le semi-groupe complet. En
particulier, \(360\) conserve deux histoires avant le quotient, quatre pas de
\(90\) et trois de \(120\), même lorsque les deux produisent le même moment scalaire.

\subsection{La section absolue de masse et la décennie \texorpdfstring{\(10^{-34}\)}{10⁻³⁴}}

Le bloc local de Hadamard
\[
 H_4=
 \begin{pmatrix}
  1&1&1&1\\
  1&1&-1&-1\\
  1&-1&1&-1\\
  1&-1&-1&1
 \end{pmatrix},
 \qquad
 \operatorname{SNF}(H_4)=\operatorname{diag}(1,2,2,4),
\]
a un conoyau d'ordre \(16\). Son lecteur déterminantal produit
\begin{equation}
 \Sigma_H=\varphi\alpha_{\rm HMT}^{16},
 \qquad
 10^{-34}<\Sigma_H<10^{-33},
 \qquad
 \left\lfloor\log_{10}\Sigma_H\right\rfloor=-34.
 \label{eq:orden-menos-34-tabla-masas}
\end{equation}
La conversion en \(\mathrm{J\,s}\) constitue une carte ultérieure : la puissance
\(10^{-34}\) naît dans la structure discrète et fixe l'échelle d'une unique droite
d'action. Ses sections antérieure et postérieure au retour sont
\begin{equation}
 \hbar_{\rm pre}=H_5\,10^{-34}\mathcal U_S,\qquad
 \hbar_{\rm ret}=\eta_{\rm ret}\,10^{-34}\mathcal U_S,
 \qquad
 R_{\rm act}=\frac{H_5}{\eta_{\rm ret}}>0.
 \label{eq:accion-absoluta-y-razon-tabla-masas}
\end{equation}
Le dédoublement possède une carte additive,
\(\delta_{2w}=H_5-\eta_{\rm ret}\), et la carte multiplicative adimensionnelle
\(R_{\rm act}\). Une durée interne \(t_0\) étant déclarée, la périodicité de
108 itérations de l'endomorphisme temporel CT108 construit
\begin{align}
 \omega_0&=\frac{2\pi}{108t_0},
 &\varepsilon_{\rm clk}&=\hbar_{\rm ret}\omega_0,
 \label{eq:cuanto-reloj-masa-rectora}\\
 \mathbf m_\star^{\rm pre}
 &=\left(H_5\,10^{-34}\mathcal U_S\right)
   \frac{2\pi}{108t_0}c_{\rm int}^{-2},
 &\mathbf m_\star^{\rm ret}
 &=\left(\eta_{\rm ret}\,10^{-34}\mathcal U_S\right)
   \frac{2\pi}{108t_0}c_{\rm int}^{-2}\in\mathcal L_M.
 \label{eq:seccion-reloj-masa-rectora}
\end{align}
Les deux sections satisfont
\(\mathbf m_\star^{\rm pre}/\mathbf m_\star^{\rm ret}=R_{\rm act}\).
La décennie \(10^{-34}\), l'horloge commune et \(c_{\rm int}^{-2}\) fixent l'échelle
absolue. Dans un rapport de masses, cette section commune se simplifie exactement ; le
raffinement relatif relève du parcours, et non d'une seconde introduction de la
puissance décimale.

\subsection{Lecteurs bidirectionnels de trajectoire et loi opératorielle sur les \(13\) multisections}

Les lecteurs qui agissent sur un parcours \(\Gamma\) sont
\begin{equation}
 K_D(\Gamma)
 =\left(D_{27}+D_{36},\,D_1,\,\frac{D_4-D_3}{2}\right),
 \qquad
 K_\Omega(\Gamma)=(\Omega,54,P_6).
 \label{eq:lectores-kd-komega-masas-rectoras}
\end{equation}
Chaque lecteur induit sur une multisection \(F\) un opérateur diagonal
\(\widehat K_F^r\), \(r\in\{D,\Omega\}\). L'indice \(s\) désigne une section ou un
bloc physique au sein de \(F\), et \(\widehat B_{F,s}^r\) est la restriction
autoadjointe de \(\widehat K_F^r\) à ce bloc. Si
\(\widehat{\mathcal A}^{(0)}_{F,s}\) est la restriction de l'opérateur réunissant la
signature, le caractère complet et la hiérarchie globale des tours, la loi générale prend
la forme
\begin{equation}
 \boxed{
 \widehat{\mathbf M}^{\,\beta,r}_{F,s}
 =\mathbf m_\star^{\rm ret}\otimes
  R_{\rm act}^{\widehat B_{F,s}^r/2}
  \widehat{\mathcal A}^{(0)}_{F,s}
  R_{\rm act}^{\widehat B_{F,s}^r/2},
 \qquad r\in\{D,\Omega\}.}
 \label{eq:ley-operatorial-general-masas-rectoras}
\end{equation}
L'écriture symétrique préserve la positivité sans imposer la commutation de
\(\widehat{\mathcal A}^{(0)}_{F,s}\) et \(\widehat B_{F,s}^r\). Dans la base
simultanément diagonale des certificats en vigueur, elle coïncide avec le produit
\(\mathbf m_\star^{\rm ret}\otimes\widehat{\mathcal A}^{(0)}_F
\exp((\log R_{\rm act})\widehat K_F^r)\).
Cette équation régit les 81 cellules, les 56 familles, les 13 multisections et les
324 parcours orientés. Hors d'une fibre de rang un, le résultat primaire est
l'opérateur complet et son spectre. La réalisation scalaire utilise le morphisme établi
\(\mathcal S_{\rm phys}:(\sigma_{\rm ud},g,\chi_{\rm fina},[\gamma]_{\rm col})
\mapsto\gamma^{\rm surv}_{P,\mathfrak M}\) de
\cref{eq:realizacion-fisica-sabor-ruta-cerrada-rev2} ; cette sélection est
indépendante des masses cibles et précède toute comparaison métrologique.

\subsection{Réduction de rang \(1\) dans la construction de l'électron central}

La section électronique est l'unique point fixe central de l'atlas :
\begin{equation}
 (r,c)=(5,5),\qquad
 \Gamma_e=(5,5,++),\qquad
 \tau_e=\mathtt{+++---+++---}.
 \label{eq:ruta-electronica-central-masas-rectoras}
\end{equation}
Les deux lecteurs y coïncident,
\begin{equation}
 K_D(\Gamma_e)=K_\Omega(\Gamma_e)=(80,54,6),\qquad
 \kappa_e=80-54\alpha_{\rm HMT}+6\Delta_4.
 \label{eq:kappa-electron-rector}
\end{equation}
La réduction de rang un de
\eqref{eq:ley-operatorial-general-masas-rectoras} fournit
\begin{equation}
 \mathbf m_e^\beta
 =\mathbf m_\star^{\rm ret}\,
   \mu_e^{(0)}R_{\rm act}^{\,80-54\alpha_{\rm HMT}+6\Delta_4},
 \label{eq:electron-rank1-ley-masas}
\end{equation}
dont la coordonnée dans la carte comparative est
\begin{equation}
 \boxed{m_e^\beta
 =0.5109989506900008086082571648\ldots\ \mathrm{MeV}/c^2.}
 \label{eq:masa-electron-tabla-completa-revfinal}
\end{equation}
La valeur de base
\(0.5109989224880660725\ldots\ \mathrm{MeV}/c^2\) conserve sa fonction
d'étape intermédiaire ; le résultat bêta de
\eqref{eq:masa-electron-tabla-completa-revfinal} régit la comparaison finale.

\subsection{Classification physique et comparaison quantitative par secteurs}

L'exposé s'organise en quatre strates coordonnées. La première présente le
résultat interne et son codomaine ; la deuxième déclare le morphisme de réalisation
physique ; la troisième montre la meilleure coordonnée numérique déjà disponible et conserve
ses hypothèses et son domaine ; la quatrième introduit la donnée externe comme confrontation ultérieure.
Chaque colonne conserve ainsi son type mathématique : résultat HMT, morphisme de
réalisation, coordonnée calculée et confrontation externe. Le tableau présente
simultanément les évaluations disponibles et les applications de construction qui les
produisent.

L'électron donne sa clôture bêta scalaire. Le photon et les gluons donnent
des sections nulles exactes dans leurs cartes de masse, ainsi que des extensions de jauge non
triviales. Les autres leptons chargés et les six quarks donnent simultanément
leurs parcours nominaux, caractères scalaires, opérateurs diagonaux et spectres sur
leurs fibres ; les neutrinos donnent l'opérateur neutre et son mélange. \(W\), \(Z\) et
\(H\) possèdent des parcours CT--108 individualisés, des empreintes complètes et des caractères
\((n_A,s_C)=(94,-1),(95,-1),(97,9)\). La paire \(K_D/K_\Omega\) appartient au
raffinement des treize multisections de l'atlas et n'est pas une prémisse de la
réalisation bosonique par parcours nominal. Les composés, les secteurs topologiques et les états
virtuels maintiennent respectivement leurs codomaines de liaison, de fusion, de tressage
et d'obstruction.

\Needspace{15\baselineskip}
\subsubsection{Correspondance typée entre objets internes et codomaines physiques}


\endgroup
\begingroup
\let\chapter\section
\let\section\subsection
\let\subsection\subsubsection
\let\subsubsection\paragraph
\let\clearpage\relax

\Needspace{12\baselineskip}
\subsubsection{Opérateurs des \(13\) multisections}

Les nombres de familles indiqués dans chaque fiche sont des cardinalités locales du
support de cette multisection. Ils ne s'additionnent pas entre fiches : le quotient des
56 familles de parcours et le quotient des 13 multisections sont deux partitions
distinctes de l'atlas de 81 cellules, reliées par les applications de projection, et non
par une décomposition disjointe commune.

\begingroup
\small
\setlength{\parskip}{.12\baselineskip}
\setlist[itemize]{leftmargin=1.65em,itemsep=0pt,topsep=.1em,
  parsep=0pt,partopsep=0pt}
\let\texttt\textnormal
Les fiches suivantes présentent l’opérateur dans son codomaine natif. Chaque spectre est reproduit intégralement ; aucune liste de valeurs propres n’est remplacée par une coordonnée de masse.
\Needspace{10\baselineskip}
\noindent\textbf{secteur leptonique chargé, profondeur G0 (centre).} \(\mathcal B_{\ell,\mathrm{G0}}\).\par
\noindent\textbf{Domaine de fibre.} Rang 1; cellule 41; sans orbite chromatique; 1 famille de parcours.
\noindent\textbf{Spectre complet de \(K_D\).}
\begin{itemize}
\item \texttt{\detokenize{(80,54,6)x1}}
\end{itemize}
\noindent\textbf{Spectre complet de \(K_\Omega\).}
\begin{itemize}
\item \texttt{\detokenize{(80,54,6)x1}}
\end{itemize}
\noindent\textbf{Conclusion dans la fibre.} Opérateur diagonal exact sur la fibre \(B_1\) ; le parcours nominal et le morphisme déclaré expriment le caractère scalaire sans remplacer le spectre de la fibre.
\medskip
\Needspace{10\baselineskip}
\noindent\textbf{secteur leptonique chargé, profondeur G2.} \(\mathcal B_{\ell,\mathrm{G2}}\).\par
\noindent\textbf{Domaine de fibre.} Rang 8; cellules 21,23,25,39,43,57,59,61; sans orbite chromatique; 8 familles de parcours.
\noindent\textbf{Spectre complet de \(K_D\).}
\begin{itemize}
\item \texttt{\detokenize{(0,24,0)x1}}
\item \texttt{\detokenize{(40,78,27)x2}}
\item \texttt{\detokenize{(80,54,7)x1}}
\item \texttt{\detokenize{(80,66,2)x1}}
\item \texttt{\detokenize{(80,66,3)x1}}
\item \texttt{\detokenize{(100,66,0)x1}}
\item \texttt{\detokenize{(100,66,2)x1}}
\end{itemize}
\noindent\textbf{Spectre complet de \(K_\Omega\).}
\begin{itemize}
\item \texttt{\detokenize{(80,54,6)x8}}
\end{itemize}
\noindent\textbf{Conclusion dans la fibre.} Opérateur diagonal exact sur la fibre \(B_1\) ; le parcours nominal et le morphisme déclaré expriment le caractère scalaire sans remplacer le spectre de la fibre.
\medskip
\Needspace{22\baselineskip}
\noindent\textbf{secteur leptonique chargé, profondeur G4.} \(\mathcal B_{\ell,\mathrm{G4}}\).\par
\noindent\textbf{Domaine de fibre.} Rang 16; cellules 1,3,5,7,9,19,27,37,45,55,63,73,75,77,79,81; sans orbite chromatique; 16 familles de parcours.
\noindent\textbf{Spectre complet de \(K_D\).}
\begin{itemize}
\item \texttt{\detokenize{(0,24,0)x1}}
\item \texttt{\detokenize{(40,24,2)x2}}
\item \texttt{\detokenize{(40,54,6)x2}}
\item \texttt{\detokenize{(40,84,30)x2}}
\item \texttt{\detokenize{(60,78,25)x3}}
\item \texttt{\detokenize{(80,66,2)x2}}
\item \texttt{\detokenize{(80,66,3)x3}}
\item \texttt{\detokenize{(80,66,4)x1}}
\end{itemize}
\noindent\textbf{Spectre complet de \(K_\Omega\).}
\begin{itemize}
\item \texttt{\detokenize{(24,54,2)x1}}
\item \texttt{\detokenize{(56,54,5)x4}}
\item \texttt{\detokenize{(80,54,6)x11}}
\end{itemize}
\noindent\textbf{Conclusion dans la fibre.} Opérateur diagonal exact sur la fibre \(B_1\) ; le parcours nominal et le morphisme déclaré expriment le caractère scalaire sans remplacer le spectre de la fibre.
\medskip
\Needspace{10\baselineskip}
\noindent\textbf{secteur neutrinique, profondeur G1.} \(\mathcal B_{\nu,\mathrm{G1}}\).\par
\noindent\textbf{Domaine de fibre.} Rang 2; cellules 40,42; sans orbite chromatique; 2 familles de parcours.
\noindent\textbf{Spectre complet de \(K_D\).}
\begin{itemize}
\item \texttt{\detokenize{(40,24,2)x1}}
\item \texttt{\detokenize{(40,78,27)x1}}
\end{itemize}
\noindent\textbf{Spectre complet de \(K_\Omega\).}
\begin{itemize}
\item \texttt{\detokenize{(40,54,4)x1}}
\item \texttt{\detokenize{(80,54,6)x1}}
\end{itemize}
\noindent\textbf{Conclusion dans la fibre.} Opérateur diagonal exact sur la fibre \(B_1\) ; le parcours nominal et le morphisme déclaré expriment le caractère scalaire sans remplacer le spectre de la fibre.
\medskip
\Needspace{10\baselineskip}
\noindent\textbf{secteur neutrinique, profondeur G2.} \(\mathcal B_{\nu,\mathrm{G2}}\).\par
\noindent\textbf{Domaine de fibre.} Rang 4; cellules 22,24,58,60; sans orbite chromatique; 4 familles de parcours.
\noindent\textbf{Spectre complet de \(K_D\).}
\begin{itemize}
\item \texttt{\detokenize{(0,24,0)x1}}
\item \texttt{\detokenize{(40,84,30)x1}}
\item \texttt{\detokenize{(60,78,25)x1}}
\item \texttt{\detokenize{(80,66,3)x1}}
\end{itemize}
\noindent\textbf{Spectre complet de \(K_\Omega\).}
\begin{itemize}
\item \texttt{\detokenize{(4,54,2)x1}}
\item \texttt{\detokenize{(80,54,6)x3}}
\end{itemize}
\noindent\textbf{Conclusion dans la fibre.} Opérateur diagonal exact sur la fibre \(B_1\) ; le parcours nominal et le morphisme déclaré expriment le caractère scalaire sans remplacer le spectre de la fibre.
\medskip
\Needspace{18\baselineskip}
\noindent\textbf{secteur neutrinique, profondeur G3.} \(\mathcal B_{\nu,\mathrm{G3}}\).\par
\noindent\textbf{Domaine de fibre.} Rang 6; cellules 20,26,38,44,56,62; sans orbite chromatique; 6 familles de parcours.
\noindent\textbf{Spectre complet de \(K_D\).}
\begin{itemize}
\item \texttt{\detokenize{(40,24,2)x2}}
\item \texttt{\detokenize{(40,78,27)x1}}
\item \texttt{\detokenize{(60,84,28)x1}}
\item \texttt{\detokenize{(80,54,6)x1}}
\item \texttt{\detokenize{(80,66,3)x1}}
\end{itemize}
\noindent\textbf{Spectre complet de \(K_\Omega\).}
\begin{itemize}
\item \texttt{\detokenize{(36,54,4)x1}}
\item \texttt{\detokenize{(56,54,5)x1}}
\item \texttt{\detokenize{(80,54,6)x4}}
\end{itemize}
\noindent\textbf{Conclusion dans la fibre.} Opérateur diagonal exact sur la fibre \(B_1\) ; le parcours nominal et le morphisme déclaré expriment le caractère scalaire sans remplacer le spectre de la fibre.
\medskip
\Needspace{10\baselineskip}
\noindent\textbf{secteur neutrinique, profondeur G4.} \(\mathcal B_{\nu,\mathrm{G4}}\).\par
\noindent\textbf{Domaine de fibre.} Rang 8; cellules 2,4,6,8,74,76,78,80; sans orbite chromatique; 8 familles de parcours.
\noindent\textbf{Spectre complet de \(K_D\).}
\begin{itemize}
\item \texttt{\detokenize{(0,24,0)x1}}
\item \texttt{\detokenize{(40,24,2)x1}}
\item \texttt{\detokenize{(40,78,27)x2}}
\item \texttt{\detokenize{(40,84,30)x1}}
\item \texttt{\detokenize{(80,54,7)x1}}
\item \texttt{\detokenize{(80,66,2)x1}}
\item \texttt{\detokenize{(100,66,0)x1}}
\end{itemize}
\noindent\textbf{Spectre complet de \(K_\Omega\).}
\begin{itemize}
\item \texttt{\detokenize{(36,54,4)x1}}
\item \texttt{\detokenize{(56,54,5)x1}}
\item \texttt{\detokenize{(80,54,6)x6}}
\end{itemize}
\noindent\textbf{Conclusion dans la fibre.} Opérateur diagonal exact sur la fibre \(B_1\) ; le parcours nominal et le morphisme déclaré expriment le caractère scalaire sans remplacer le spectre de la fibre.
\medskip
\Needspace{10\baselineskip}
\noindent\textbf{branche de quarks de type bas, profondeur G1.} \(\mathcal B_{d,\mathrm{G1}}\).\par
\noindent\textbf{Domaine de fibre.} Rang 2; cellules 32,50; B | G; 2 familles de parcours.
\noindent\textbf{Spectre complet de \(K_D\).}
\begin{itemize}
\item \texttt{\detokenize{(40,24,2)x1}}
\item \texttt{\detokenize{(40,78,27)x1}}
\end{itemize}
\noindent\textbf{Spectre complet de \(K_\Omega\).}
\begin{itemize}
\item \texttt{\detokenize{(40,54,4)x1}}
\item \texttt{\detokenize{(80,54,6)x1}}
\end{itemize}
\noindent\textbf{Conclusion dans la fibre.} Opérateur diagonal exact sur la fibre \(B_1\) ; le parcours nominal et le morphisme déclaré expriment le caractère scalaire sans remplacer le spectre de la fibre.
\medskip
\Needspace{10\baselineskip}
\noindent\textbf{branche de quarks de type bas, profondeur G2.} \(\mathcal B_{d,\mathrm{G2}}\).\par
\noindent\textbf{Domaine de fibre.} Rang 4; cellules 30,34,48,52; B | G | R; 4 familles de parcours.
\noindent\textbf{Spectre complet de \(K_D\).}
\begin{itemize}
\item \texttt{\detokenize{(0,24,0)x1}}
\item \texttt{\detokenize{(40,84,30)x1}}
\item \texttt{\detokenize{(60,78,25)x1}}
\item \texttt{\detokenize{(80,66,2)x1}}
\end{itemize}
\noindent\textbf{Spectre complet de \(K_\Omega\).}
\begin{itemize}
\item \texttt{\detokenize{(4,54,2)x1}}
\item \texttt{\detokenize{(80,54,6)x3}}
\end{itemize}
\noindent\textbf{Conclusion dans la fibre.} Opérateur diagonal exact sur la fibre \(B_1\) ; le parcours nominal et le morphisme déclaré expriment le caractère scalaire sans remplacer le spectre de la fibre.
\medskip
\Needspace{10\baselineskip}
\noindent\textbf{branche de quarks de type bas, profondeur G3.} \(\mathcal B_{d,\mathrm{G3}}\).\par
\noindent\textbf{Domaine de fibre.} Rang 6; cellules 12,14,16,66,68,70; B | G | R; 6 familles de parcours.
\noindent\textbf{Spectre complet de \(K_D\).}
\begin{itemize}
\item \texttt{\detokenize{(40,24,2)x2}}
\item \texttt{\detokenize{(40,78,27)x1}}
\item \texttt{\detokenize{(60,84,28)x1}}
\item \texttt{\detokenize{(80,54,6)x1}}
\item \texttt{\detokenize{(80,66,4)x1}}
\end{itemize}
\noindent\textbf{Spectre complet de \(K_\Omega\).}
\begin{itemize}
\item \texttt{\detokenize{(36,54,4)x1}}
\item \texttt{\detokenize{(80,54,6)x5}}
\end{itemize}
\noindent\textbf{Conclusion dans la fibre.} Opérateur diagonal exact sur la fibre \(B_1\) ; le parcours nominal et le morphisme déclaré expriment le caractère scalaire sans remplacer le spectre de la fibre.
\medskip
\Needspace{10\baselineskip}
\noindent\textbf{branche de quarks de type bas, profondeur G4.} \(\mathcal B_{d,\mathrm{G4}}\).\par
\noindent\textbf{Domaine de fibre.} Rang 8; cellules 10,18,28,36,46,54,64,72; B | G | R; 8 familles de parcours.
\noindent\textbf{Spectre complet de \(K_D\).}
\begin{itemize}
\item \texttt{\detokenize{(0,24,0)x1}}
\item \texttt{\detokenize{(40,24,2)x1}}
\item \texttt{\detokenize{(40,78,27)x2}}
\item \texttt{\detokenize{(40,84,30)x1}}
\item \texttt{\detokenize{(80,54,7)x1}}
\item \texttt{\detokenize{(80,66,3)x1}}
\item \texttt{\detokenize{(100,66,2)x1}}
\end{itemize}
\noindent\textbf{Spectre complet de \(K_\Omega\).}
\begin{itemize}
\item \texttt{\detokenize{(36,54,4)x1}}
\item \texttt{\detokenize{(80,54,6)x7}}
\end{itemize}
\noindent\textbf{Conclusion dans la fibre.} Opérateur diagonal exact sur la fibre \(B_1\) ; le parcours nominal et le morphisme déclaré expriment le caractère scalaire sans remplacer le spectre de la fibre.
\medskip
\Needspace{10\baselineskip}
\noindent\textbf{branche de quarks de type haut, profondeur G1.} \(\mathcal B_{u,\mathrm{G1}}\).\par
\noindent\textbf{Domaine de fibre.} Rang 4; cellules 31,33,49,51; B | G | R; 3 familles de parcours.
\noindent\textbf{Spectre complet de \(K_D\).}
\begin{itemize}
\item \texttt{\detokenize{(40,54,6)x1}}
\item \texttt{\detokenize{(60,78,25)x1}}
\item \texttt{\detokenize{(80,66,2)x1}}
\item \texttt{\detokenize{(80,66,3)x1}}
\end{itemize}
\noindent\textbf{Spectre complet de \(K_\Omega\).}
\begin{itemize}
\item \texttt{\detokenize{(80,54,6)x4}}
\end{itemize}
\noindent\textbf{Conclusion dans la fibre.} Opérateur diagonal exact sur la fibre \(B_1\) ; le parcours nominal et le morphisme déclaré expriment le caractère scalaire sans remplacer le spectre de la fibre.
\medskip
\Needspace{29\baselineskip}
\noindent\textbf{branche de quarks de type haut, profondeur G3.} \(\mathcal B_{u,\mathrm{G3}}\).\par
\noindent\textbf{Domaine de fibre.} Rang 12; cellules 11,13,15,17,29,35,47,53,65,67,69,71; B | G | R; 12 familles de parcours.
\noindent\textbf{Spectre complet de \(K_D\).}
\begin{itemize}
\item \texttt{\detokenize{(40,24,2)x2}}
\item \texttt{\detokenize{(40,78,27)x2}}
\item \texttt{\detokenize{(60,84,28)x1}}
\item \texttt{\detokenize{(80,54,6)x3}}
\item \texttt{\detokenize{(80,66,3)x1}}
\item \texttt{\detokenize{(80,66,4)x1}}
\item \texttt{\detokenize{(100,66,0)x1}}
\item \texttt{\detokenize{(100,66,2)x1}}
\end{itemize}
\noindent\textbf{Spectre complet de \(K_\Omega\).}
\begin{itemize}
\item \texttt{\detokenize{(56,54,5)x3}}
\item \texttt{\detokenize{(80,54,6)x9}}
\end{itemize}
\noindent\textbf{Conclusion dans la fibre.} Opérateur diagonal exact sur la fibre \(B_1\) ; le parcours nominal et le morphisme déclaré expriment le caractère scalaire sans remplacer le spectre de la fibre.
\medskip
\endgroup


\Needspace{15\baselineskip}
\subsubsection{Morphismes de réalisation physique}

\begingroup
\setlength{\tabcolsep}{4.5pt}
\renewcommand{\arraystretch}{1.28}
\begin{longtable}{@{}>{\raggedright\arraybackslash}p{.14\textwidth}>{\raggedright\arraybackslash}p{.33\textwidth}>{\raggedright\arraybackslash}p{.25\textwidth}>{\raggedright\arraybackslash}p{.22\textwidth}@{}}
\caption{Réalisation physique des leptons chargés.}\label{tab:masas-realizacion-leptones}\\
\toprule
\textbf{Classe} & \textbf{Objeto HMT} & \textbf{Application de réalisation} & \textbf{Conclusion et conditions} \\
\midrule
\endfirsthead
\toprule
\textbf{Classe} & \textbf{Objeto HMT} & \textbf{Application de réalisation} & \textbf{Conclusion et conditions} \\
\midrule
\endhead
\midrule
\multicolumn{4}{r}{\footnotesize Suite à la page suivante.}\\
\endfoot
\bottomrule
\endlastfoot
$e^-/e^+$ & Cellule centrale $(5,5)$, parcours $\Gamma_e=(5,5,++)$ et $K_D(\Gamma_e)=K_\Omega(\Gamma_e)=(80,54,6)$. & Réduction de rang un fermée. & Clôture bêta de rang un ; sélecteur indépendant des grandeurs métrologiques utilisées comme cibles. \\
$\mu^-/\mu^+$ & famille leptonique chargée ; niveaux G2 ; 8 cellules ; $K_D$ : 7 profils ; $K_\Omega$ : 1 profil & Génération et signature fine $\to$ parcours persistant $\to\chi_{\rm mass}\to\mathcal L_M$. & Parcours nominal et caractère scalaire fermés ; la paire d’opérateurs conserve le raffinement complet de la fibre. \\
$\tau^-/\tau^+$ & famille leptonique chargée ; niveaux G4 ; 16 cellules ; $K_D$ : 8 profils ; $K_\Omega$ : 3 profils & Génération et signature fine $\to$ parcours persistant $\to\chi_{\rm mass}\to\mathcal L_M$. & Parcours nominal et caractère scalaire fermés ; la paire d’opérateurs conserve le raffinement complet de la fibre. \\
\end{longtable}
\endgroup

\begingroup
\setlength{\tabcolsep}{4.5pt}
\renewcommand{\arraystretch}{1.28}
\begin{longtable}{@{}>{\raggedright\arraybackslash}p{.14\textwidth}>{\raggedright\arraybackslash}p{.33\textwidth}>{\raggedright\arraybackslash}p{.25\textwidth}>{\raggedright\arraybackslash}p{.22\textwidth}@{}}
\caption{Réalisation physique des neutrinos et transport de bases.}\label{tab:masas-realizacion-neutrinos}\\
\toprule
\textbf{Classe} & \textbf{Objeto HMT} & \textbf{Application de réalisation} & \textbf{Conclusion et conditions} \\
\midrule
\endfirsthead
\toprule
\textbf{Classe} & \textbf{Objeto HMT} & \textbf{Application de réalisation} & \textbf{Conclusion et conditions} \\
\midrule
\endhead
\midrule
\multicolumn{4}{r}{\footnotesize Suite à la page suivante.}\\
\endfoot
\bottomrule
\endlastfoot
$\nu_e/\bar\nu_e$ & famille neutre ; niveaux G1, G2, G3, G4 ; 20 cellules ; $K_D$ : 11 profils ; $K_\Omega$ : 5 profils & Enregistrement neutre de rang plein $\to$ vecteurs propres $\to$ base PMNS. & La structure extérieure/CAR fixe le caractère fermionique ; PMNS transporte les bases et ne crée pas de valeurs propres de masse. \\
$\nu_\mu/\bar\nu_\mu$ & famille neutre ; niveaux G1, G2, G3, G4 ; 20 cellules ; $K_D$ : 11 profils ; $K_\Omega$ : 5 profils & Enregistrement neutre de rang plein $\to$ vecteurs propres $\to$ base PMNS. & La structure extérieure/CAR fixe le caractère fermionique ; PMNS transporte les bases et ne crée pas de valeurs propres de masse. \\
$\nu_\tau/\bar\nu_\tau$ & famille neutre ; niveaux G1, G2, G3, G4 ; 20 cellules ; $K_D$ : 11 profils ; $K_\Omega$ : 5 profils & Enregistrement neutre de rang plein $\to$ vecteurs propres $\to$ base PMNS. & La structure extérieure/CAR fixe le caractère fermionique ; PMNS transporte les bases et ne crée pas de valeurs propres de masse. \\
\end{longtable}
\endgroup

\begingroup
\setlength{\tabcolsep}{4.5pt}
\renewcommand{\arraystretch}{1.28}
\begin{longtable}{@{}>{\raggedright\arraybackslash}p{.14\textwidth}>{\raggedright\arraybackslash}p{.33\textwidth}>{\raggedright\arraybackslash}p{.25\textwidth}>{\raggedright\arraybackslash}p{.22\textwidth}@{}}
\caption{Réalisation physique des quarks : génération, saveur et couleur.}\label{tab:masas-realizacion-quarks}\\
\toprule
\textbf{Classe} & \textbf{Objeto HMT} & \textbf{Application de réalisation} & \textbf{Conclusion et conditions} \\
\midrule
\endfirsthead
\toprule
\textbf{Classe} & \textbf{Objeto HMT} & \textbf{Application de réalisation} & \textbf{Conclusion et conditions} \\
\midrule
\endhead
\midrule
\multicolumn{4}{r}{\footnotesize Suite à la page suivante.}\\
\endfoot
\bottomrule
\endlastfoot
$u/\bar u$ & branche de quarks \emph{up} ; niveaux G1 et G3 ; couleur $\mathbb Z/3$ ; 16 cellules ; $K_D$ : 11 profils ; $K_\Omega$ : 2 profils & Saveur--couleur--génération $\to$ parcours persistant $\to\chi_{\rm mass}\to\mathcal L_M$. & La saveur, la couleur et la génération déterminent la classe de parcours ; caractère scalaire et paire d’opérateurs fermés. \\
$d/\bar d$ & branche de quarks \emph{down} ; niveaux G1, G2, G3, G4 ; couleur $\mathbb Z/3$ ; 20 cellules ; $K_D$ : 12 profils ; $K_\Omega$ : 4 profils & Saveur--couleur--génération $\to$ parcours persistant $\to\chi_{\rm mass}\to\mathcal L_M$. & La saveur, la couleur et la génération déterminent la classe de parcours ; caractère scalaire et paire d’opérateurs fermés. \\
$c/\bar c$ & branche de quarks \emph{up} ; niveaux G1 et G3 ; couleur $\mathbb Z/3$ ; 16 cellules ; $K_D$ : 11 profils ; $K_\Omega$ : 2 profils & Saveur--couleur--génération $\to$ parcours persistant $\to\chi_{\rm mass}\to\mathcal L_M$. & La saveur, la couleur et la génération déterminent la classe de parcours ; caractère scalaire et paire d’opérateurs fermés. \\
$s/\bar s$ & branche de quarks \emph{down} ; niveaux G1, G2, G3, G4 ; couleur $\mathbb Z/3$ ; 20 cellules ; $K_D$ : 12 profils ; $K_\Omega$ : 4 profils & Saveur--couleur--génération $\to$ parcours persistant $\to\chi_{\rm mass}\to\mathcal L_M$. & La saveur, la couleur et la génération déterminent la classe de parcours ; caractère scalaire et paire d’opérateurs fermés. \\
$t/\bar t$ & branche de quarks \emph{up} ; niveaux G1 et G3 ; couleur $\mathbb Z/3$ ; 16 cellules ; $K_D$ : 11 profils ; $K_\Omega$ : 2 profils & Saveur--couleur--génération $\to$ parcours persistant $\to\chi_{\rm mass}\to\mathcal L_M$. & La saveur, la couleur et la génération déterminent la classe de parcours ; caractère scalaire et paire d’opérateurs fermés. \\
$b/\bar b$ & branche de quarks \emph{down} ; niveaux G1, G2, G3, G4 ; couleur $\mathbb Z/3$ ; 20 cellules ; $K_D$ : 12 profils ; $K_\Omega$ : 4 profils & Saveur--couleur--génération $\to$ parcours persistant $\to\chi_{\rm mass}\to\mathcal L_M$. & La saveur, la couleur et la génération déterminent la classe de parcours ; caractère scalaire et paire d’opérateurs fermés. \\
\end{longtable}
\endgroup

\Needspace{12\baselineskip}
\paragraph{Localisation structurelle des \(6\) saveurs de quarks}
\begingroup
\setlength{\parskip}{.3\baselineskip}
La position d’une saveur de quark est spécifiée par quatre données simultanées : branche, génération, support cellulaire et famille de parcours orientés. Les trois générations d’une même branche partagent le support de l’atlas ; le classificateur fin de génération conserve la distinction physique qui doit encore être réalisée comme vecteur propre et section de masse.
\begingroup
\setlength{\tabcolsep}{4.5pt}
\renewcommand{\arraystretch}{1.28}
\begin{longtable}{@{}>{\raggedright\arraybackslash}p{.12\textwidth}>{\raggedright\arraybackslash}p{.17\textwidth}>{\raggedright\arraybackslash}p{.19\textwidth}>{\raggedright\arraybackslash}p{.46\textwidth}@{}}
\caption{Classification générationnelle des six saveurs de quarks.}\label{tab:masas-posicion-quarks-seis-sabores}\\
\toprule
\textbf{Saveur} & \textbf{Génération} & \textbf{Branche} & \textbf{Classificateur de branche et de génération} \\
\midrule
\endfirsthead
\toprule
\textbf{Saveur} & \textbf{Génération} & \textbf{Branche} & \textbf{Classificateur de branche et de génération} \\
\midrule
\endhead
\midrule
\multicolumn{4}{r}{\footnotesize Suite à la page suivante.}\\
\endfoot
\bottomrule
\endlastfoot
$u/\bar u$ & première & supérieur & $(\operatorname{up},1,\chi_{\rm fina})$ ; paire nominale de parcours $(21,30)$ sur l’axe $\mathrm{N\!-\!S}$ \\
$d/\bar d$ & première & inférieur & $(\operatorname{down},1,\chi_{\rm fina})$ ; paire nominale de parcours $(20,32)$ sur l’axe $\mathrm{E\!-\!O}$ \\
$c/\bar c$ & deuxième & supérieur & $(\operatorname{up},2,\chi_{\rm fina})$ ; paire nominale de parcours $(21,29)$ sur l’axe $\mathrm{E\!-\!O}$ \\
$s/\bar s$ & deuxième & inférieur & $(\operatorname{down},2,\chi_{\rm fina})$ ; paire nominale de parcours $(17,27)$ sur l’axe $\mathrm{N\!-\!S}$ \\
$t/\bar t$ & troisième & supérieur & $(\operatorname{up},3,\chi_{\rm fina})$ ; paire nominale de parcours $(21,26)$ sur l’axe $\mathrm{N\!-\!S}$ \\
$b/\bar b$ & troisième & inférieur & $(\operatorname{down},3,\chi_{\rm fina})$ ; paire nominale de parcours $(21,30)$ sur l’axe $\mathrm{E\!-\!O}$ \\
\end{longtable}
\endgroup

\Needspace{18\baselineskip}
\noindent\textbf{Support structurel de la branche supérieure.}\par
\noindent\textit{Multisections.} \(\mathcal B_{u,\mathrm{G1}}\), \(\mathcal B_{u,\mathrm{G3}}\).\par
\noindent\textit{Profondeurs et cardinaux.} $\mathrm{G1}$, $\mathrm{G3}$ ; 16 cellules et 14 familles de parcours distinctes.\par
\noindent\textit{Orbite chromatique et orientations.} $\mathcal C=\{B,G,R\}$ et $\mathcal O=\{(-1,-1),(-1,1),(1,-1),(1,1)\}$.\par
\noindent\textit{Cellules et familles orientées $B_1$.} Chaque chaîne $k\leftrightarrow(r,c)\leftrightarrow\Gamma_{r,c}^{(h,v)}$ conserve l’indice cellulaire, sa coordonnée dans la carte $9\times9$ et les deux signes d’orientation de l’enregistrement de parcours enrichi.\par
\noindent Fila $r=2$: $11\leftrightarrow(2,2)\leftrightarrow\Gamma_{2,2}^{(-,+)}$, $13\leftrightarrow(2,4)\leftrightarrow\Gamma_{2,4}^{(+,-)}$, $15\leftrightarrow(2,6)\leftrightarrow\Gamma_{2,6}^{(-,+)}$, $17\leftrightarrow(2,8)\leftrightarrow\Gamma_{2,8}^{(-,-)}$.\par
\noindent Fila $r=4$: $29\leftrightarrow(4,2)\leftrightarrow\Gamma_{4,2}^{(+,-)}$, $31\leftrightarrow(4,4)\leftrightarrow\Gamma_{4,4}^{(-,+)}$, $33\leftrightarrow(4,6)\leftrightarrow\Gamma_{4,6}^{(+,+)}$, $35\leftrightarrow(4,8)\leftrightarrow\Gamma_{4,8}^{(-,-)}$.\par
\noindent Fila $r=6$: $47\leftrightarrow(6,2)\leftrightarrow\Gamma_{6,2}^{(-,+)}$, $49\leftrightarrow(6,4)\leftrightarrow\Gamma_{6,4}^{(-,-)}$, $51\leftrightarrow(6,6)\leftrightarrow\Gamma_{6,6}^{(+,+)}$, $53\leftrightarrow(6,8)\leftrightarrow\Gamma_{6,8}^{(+,+)}$.\par
\noindent Fila $r=8$: $65\leftrightarrow(8,2)\leftrightarrow\Gamma_{8,2}^{(+,+)}$, $67\leftrightarrow(8,4)\leftrightarrow\Gamma_{8,4}^{(+,+)}$, $69\leftrightarrow(8,6)\leftrightarrow\Gamma_{8,6}^{(-,-)}$, $71\leftrightarrow(8,8)\leftrightarrow\Gamma_{8,8}^{(+,+)}$.\par
\noindent\textit{Morphisme de réalisation fermé.} Pour $g\in\{1,2,3\}$, la branche, la génération, l’orbite de couleur et la signature fine déterminent
\[\chi_{\rm fina}(\operatorname{up},g,\mathcal C)\longrightarrow\Gamma_{\rm enr}\longrightarrow\bigl(\operatorname{quot}_{\mathcal C},S(B_D,B_\Omega)\bigr).\]
Cette flèche sélectionne le vecteur propre physique, puis exprime la coordonnée scalaire dans $\mathcal L_M$ sans utiliser la valeur de comparaison.\par
\medskip
\Needspace{18\baselineskip}
\noindent\textbf{Support structurel de la branche inférieure.}\par
\noindent\textit{Multisections.} \(\mathcal B_{d,\mathrm{G1}}\), \(\mathcal B_{d,\mathrm{G2}}\), \(\mathcal B_{d,\mathrm{G3}}\), \(\mathcal B_{d,\mathrm{G4}}\).\par
\noindent\textit{Profondeurs et cardinaux.} $\mathrm{G1}$, $\mathrm{G2}$, $\mathrm{G3}$, $\mathrm{G4}$ ; 20 cellules et 18 familles de parcours distinctes.\par
\noindent\textit{Orbite chromatique et orientations.} $\mathcal C=\{B,G,R\}$ et $\mathcal O=\{(-1,-1),(-1,1),(1,-1),(1,1)\}$.\par
\noindent\textit{Cellules et familles orientées $B_1$.} Chaque chaîne $k\leftrightarrow(r,c)\leftrightarrow\Gamma_{r,c}^{(h,v)}$ conserve l’indice cellulaire, sa coordonnée dans la carte $9\times9$ et les deux signes d’orientation de l’enregistrement de parcours enrichi.\par
\noindent Fila $r=2$: $10\leftrightarrow(2,1)\leftrightarrow\Gamma_{2,1}^{(+,+)}$, $12\leftrightarrow(2,3)\leftrightarrow\Gamma_{2,3}^{(-,-)}$, $14\leftrightarrow(2,5)\leftrightarrow\Gamma_{2,5}^{(+,+)}$, $16\leftrightarrow(2,7)\leftrightarrow\Gamma_{2,7}^{(-,+)}$, $18\leftrightarrow(2,9)\leftrightarrow\Gamma_{2,9}^{(-,+)}$.\par
\noindent Fila $r=4$: $28\leftrightarrow(4,1)\leftrightarrow\Gamma_{4,1}^{(-,-)}$, $30\leftrightarrow(4,3)\leftrightarrow\Gamma_{4,3}^{(+,+)}$, $32\leftrightarrow(4,5)\leftrightarrow\Gamma_{4,5}^{(+,-)}$, $34\leftrightarrow(4,7)\leftrightarrow\Gamma_{4,7}^{(-,+)}$, $36\leftrightarrow(4,9)\leftrightarrow\Gamma_{4,9}^{(-,+)}$.\par
\noindent Fila $r=6$: $46\leftrightarrow(6,1)\leftrightarrow\Gamma_{6,1}^{(-,-)}$, $48\leftrightarrow(6,3)\leftrightarrow\Gamma_{6,3}^{(-,+)}$, $50\leftrightarrow(6,5)\leftrightarrow\Gamma_{6,5}^{(-,+)}$, $52\leftrightarrow(6,7)\leftrightarrow\Gamma_{6,7}^{(-,+)}$, $54\leftrightarrow(6,9)\leftrightarrow\Gamma_{6,9}^{(+,-)}$.\par
\noindent Fila $r=8$: $64\leftrightarrow(8,1)\leftrightarrow\Gamma_{8,1}^{(+,-)}$, $66\leftrightarrow(8,3)\leftrightarrow\Gamma_{8,3}^{(-,+)}$, $68\leftrightarrow(8,5)\leftrightarrow\Gamma_{8,5}^{(-,+)}$, $70\leftrightarrow(8,7)\leftrightarrow\Gamma_{8,7}^{(+,-)}$, $72\leftrightarrow(8,9)\leftrightarrow\Gamma_{8,9}^{(-,+)}$.\par
\noindent\textit{Morphisme de réalisation fermé.} Pour $g\in\{1,2,3\}$, la branche, la génération, l’orbite de couleur et la signature fine déterminent
\[\chi_{\rm fina}(\operatorname{down},g,\mathcal C)\longrightarrow\Gamma_{\rm enr}\longrightarrow\bigl(\operatorname{quot}_{\mathcal C},S(B_D,B_\Omega)\bigr).\]
Cette flèche sélectionne le vecteur propre physique, puis exprime la coordonnée scalaire dans $\mathcal L_M$ sans utiliser la valeur de comparaison.\par
\medskip
\endgroup

\begingroup
\setlength{\tabcolsep}{4.5pt}
\renewcommand{\arraystretch}{1.28}
\begin{longtable}{@{}>{\raggedright\arraybackslash}p{.14\textwidth}>{\raggedright\arraybackslash}p{.33\textwidth}>{\raggedright\arraybackslash}p{.25\textwidth}>{\raggedright\arraybackslash}p{.22\textwidth}@{}}
\caption{Réalisation physique des secteurs de jauge et scalaire.}\label{tab:masas-realizacion-gauge-escalar}\\
\toprule
\textbf{Classe} & \textbf{Objeto HMT} & \textbf{Application de réalisation} & \textbf{Conclusion et conditions} \\
\midrule
\endfirsthead
\toprule
\textbf{Classe} & \textbf{Objeto HMT} & \textbf{Application de réalisation} & \textbf{Conclusion et conditions} \\
\midrule
\endhead
\midrule
\multicolumn{4}{r}{\footnotesize Suite à la page suivante.}\\
\endfoot
\bottomrule
\endlastfoot
photon $\gamma$ & Carte nulle sectorielle et représentation de jauge correspondante. & Inclusion exacte de la section nulle dans la droite de masse. & Nullité exacte dans sa carte sectorielle. \\
gluones $g$ & Carte nulle sectorielle et représentation de jauge correspondante. & Inclusion exacte de la section nulle dans la droite de masse. & Nullité exacte dans sa carte sectorielle. \\
$W^\pm$ & parcours CT--108 W : 15/37, axe E-O, $(N,E,S,O)=(23,44,11,30)$, feuillet/résidu $(41,8)$, $(n_A,s_C)=(94,-1)$, $(W_+,W_-)=(563,565)$ ; caractère scalaire exact & Parcours CT--108 nominal $\to(n_A,s_C)\to(W_+,W_-)\to\chi_{\rm mass}\to\mathcal L_M$. & Parcours de jauge CT--108, signature et caractère scalaire déterminés avant la comparaison. \\
$Z^0$ & parcours CT--108 Z : 20/23, axe E-O, $(N,E,S,O)=(33,45,10,20)$, feuillet/résidu $(57,7)$, $(n_A,s_C)=(95,-1)$, $(W_+,W_-)=(569,571)$ ; caractère scalaire exact & Parcours CT--108 nominal $\to(n_A,s_C)\to(W_+,W_-)\to\chi_{\rm mass}\to\mathcal L_M$. & Parcours de jauge CT--108, signature et caractère scalaire déterminés avant la comparaison. \\
Higgs $H$ & parcours CT--108 H : 24/29, axe N-S, $(N,E,S,O)=(40,34,22,12)$, feuillet/résidu $(47,10)$, $(n_A,s_C)=(97,9)$, $(W_+,W_-)=(591,573)$ ; caractère scalaire exact & Parcours CT--108 nominal $\to(n_A,s_C)\to(W_+,W_-)\to\chi_{\rm mass}\to\mathcal L_M$. & Parcours scalaire CT--108, signature et caractère déterminés avant la comparaison. \\
\end{longtable}
\endgroup

\begingroup
\setlength{\tabcolsep}{4.5pt}
\renewcommand{\arraystretch}{1.28}
\begin{longtable}{@{}>{\raggedright\arraybackslash}p{.14\textwidth}>{\raggedright\arraybackslash}p{.33\textwidth}>{\raggedright\arraybackslash}p{.25\textwidth}>{\raggedright\arraybackslash}p{.22\textwidth}@{}}
\caption{Réalisation des composés, des secteurs topologiques et des états virtuels.}\label{tab:masas-realizacion-no-escalar}\\
\toprule
\textbf{Classe} & \textbf{Objeto HMT} & \textbf{Application de réalisation} & \textbf{Conclusion et conditions} \\
\midrule
\endfirsthead
\toprule
\textbf{Classe} & \textbf{Objeto HMT} & \textbf{Application de réalisation} & \textbf{Conclusion et conditions} \\
\midrule
\endhead
\midrule
\multicolumn{4}{r}{\footnotesize Suite à la page suivante.}\\
\endfoot
\bottomrule
\endlastfoot
mésons & Enregistrement double, frontière de liaison et signature composée. & Enregistrement composé et liaison $\to$ carte scalaire. & Composition des enregistrements ; la liaison précède la carte scalaire. \\
baryons & Enregistrement triple, retenue et signature composée. & Enregistrement composé et liaison $\to$ carte scalaire. & Composition des enregistrements ; la liaison précède la carte scalaire. \\
Fibonacci $(1,\tau)$ & Anneau de fusion de Fibonacci : $\tau\otimes\tau=1+\tau$. & Fusion et tressage $\to$ observable physique du secteur. & Codomaine exact de fusion et de tressage ; la comparaison utilise ses invariants propres. \\
Ising $(1,\psi,\sigma)$ ; réalisation de Majorana séparée & Anneau de fusion d’Ising : $(1,\psi,\sigma)$. & Fusion et tressage $\to$ observable physique du secteur. & Codomaine exact de fusion et de tressage ; la comparaison utilise ses invariants propres. \\
sectores $\mathbb Z/6$ & Caractère de tressage sur $\mathbb Z/6$. & Fusion et tressage $\to$ observable physique du secteur. & Codomaine exact de fusion et de tressage ; la comparaison utilise ses invariants propres. \\
secciones virtuales & Système inverse d’extensions et horizon de persistance. & Prolongement enregistré $\to$ horizon physique de persistance. & Objet exact de prolongement et d’obstruction ; sa sortie est l’horizon de persistance. \\
\end{longtable}
\endgroup


\Needspace{15\baselineskip}
\subsubsection{Comparabilité physique dans les codomaines natifs}

La comparaison externe se définit sur le codomaine natif de chaque classe. Dans les
classes à section scalaire établie, on compare des grandeurs homogènes, avec unité,
schéma et échelle déclarés. Dans les fibres opératorielles, on compare le spectre
interne et l'on conserve explicitement le morphisme de réalisation établi qui
sélectionne le parcours sans utiliser la valeur externe. Les enregistrements mésoniques et baryoniques n'admettent une comparaison état par
état qu'après incorporation de la frontière de liaison. Dans le secteur de Fibonacci,
l'anneau basé et \(\operatorname{FPdim}(\tau)=\varphi\) sont établis ; les matrices
\(F/R\), le pentagone, l'hexagone, le Hamiltonien et le gap relèvent de la
réalisation ultérieure. Ising et \(\mathbb Z/6\) conservent séparément leurs invariants
de fusion et de tressage, et la virtualité se compare par la persistance et la structure
analytique. Dans aucun de ces six
cas il n'existe de masse scalaire universelle de classe ; en attribuer une remplacerait l'objet
démontré par une valeur numérique sans morphisme.

\endgroup
\begingroup
\let\chapter\section
\let\section\subsection
\let\subsection\subsubsection
\let\subsubsection\paragraph
\let\clearpage\relax

\Needspace{15\baselineskip}
\subsubsection{Évaluations internes de signatures explicites et confrontation ultérieure}

Le tableau suivant présente lisiblement les dix-sept coordonnées HMT--MD déjà
calculées --- onze élémentaires et six composées --- et les confronte à
l'édition 2026 du jeu de données du Particle Data Group, avec une date de clôture
éditoriale au 15 janvier 2026 \cite{PDG2026}. Elles ne constituent pas une autre ontologie
distincte de la couche opératorielle précédente. Le symbole \(\dagger\)
identifie une évaluation obtenue en appliquant l'application de construction à une signature
explicite. Dans les onze lignes élémentaires, la signature procède du parcours nominal
déterminé sans masses ; dans les six lignes composées, elle procède de l'enregistrement de liaison
déclaré. Chaque valeur dépend de cette signature, du lecteur bidirectionnel et de la carte
dimensionnelle commune ; le jeu de données externe intervient seulement dans la comparaison
et ne sélectionne ni saveur, ni génération, ni couleur, ni parcours.


Pour les leptons chargés et les bosons massifs, on présente la différence relative entre
les deux coordonnées centrales et, lorsque le jeu de données fournit une incertitude
bilatérale, son quotient par cette incertitude. Pour les quarks, on présente la
différence arithmétique entre valeurs centrales, mais l'emploi du terme
\emph{résidu métrologique} reste réservé : il exige que la masse interne et la référence externe
utilisent la même définition, le même schéma de renormalisation et la même échelle.
Les masses légères se lisent dans la convention \(\overline{\mathrm{MS}}\) à
\(2\,\mathrm{GeV}\) ; celles de \(c\) et de \(b\), dans
\(\mu=\overline m_q\) ; la masse supérieure procède d'une reconstruction cinématique.

\begingroup
\setlength{\tabcolsep}{4.5pt}
\renewcommand{\arraystretch}{1.28}
\begin{longtable}{@{}>{\raggedright\arraybackslash}p{.12\textwidth}>{\raggedright\arraybackslash}p{.27\textwidth}>{\raggedright\arraybackslash}p{.27\textwidth}>{\raggedright\arraybackslash}p{.28\textwidth}@{}}
\caption{Évaluations internes de signatures explicites et comparaison métrologique ultérieure.}\label{tab:masas-evaluaciones-firmas-explicitas}\\
\toprule
\textbf{État} & \textbf{Coordonnée de la signature} & \textbf{Référence physique externe} & \textbf{Portée} \\
\midrule
\endfirsthead
\toprule
\textbf{État} & \textbf{Coordonnée de la signature} & \textbf{Référence physique externe} & \textbf{Portée} \\
\midrule
\endhead
\midrule
\multicolumn{4}{r}{\footnotesize Suite à la page suivante.}\\
\endfoot
\bottomrule
\endlastfoot
% HMT-MASS-EXPLICIT-SIGNATURE:SM-LC-2
$\mu^-/\mu^+$ & $105.658360\allowbreak{}294435\allowbreak{}829303\allowbreak{}4\;\mathrm{MeV}/c^2$\textsuperscript{\(\dagger\)}; tour E3, signature (42,-3,8,0,2) & $105.6583755\pm0.0000023\;\mathrm{MeV}/c^2$ ; masse leptonique au repos ; moyenne PDG et conversion métrologique déclarée & Différence relative : $-0.143912\allowbreak{}530\;\mathrm{ppm}$\textsuperscript{\(\dagger\)} ; $-6.611114\,\sigma$. La valeur HMT est l’évaluation exacte de la signature déclarée. \\
% HMT-MASS-EXPLICIT-SIGNATURE:SM-LC-3
$\tau^-/\tau^+$ & $1776.860917\allowbreak{}631432\allowbreak{}295595\;\mathrm{MeV}/c^2$\textsuperscript{\(\dagger\)}; tour E3, signature (64,0,25,-1,6) & $1776.93\pm0.09\;\mathrm{MeV}/c^2$ ; moyenne directe PDG & Différence relative : $-38.877371\allowbreak{}966\;\mathrm{ppm}$\textsuperscript{\(\dagger\)} ; $-0.813547\,\sigma$. La valeur HMT est l’évaluation exacte de la signature déclarée. \\
% HMT-MASS-EXPLICIT-SIGNATURE:SM-Q-U1
$u/\bar u$ & $2.165924\allowbreak{}536173\allowbreak{}907\;\mathrm{MeV}/c^2$\textsuperscript{\(\dagger\)}; carte angulaire, signature nominale de parcours (11,7) & $2.16\pm0.07\;\mathrm{MeV}/c^2$ ; masse $\overline{\mathrm{MS}}$ à $\mu=2\,\mathrm{GeV}$ ; 90\% C.L. & Différence entre les valeurs centrales : $2742.840821\allowbreak{}253\;\mathrm{ppm}$\textsuperscript{\(\dagger\)}. Il ne s’agit pas d’un résidu métrologique tant que le même schéma et la même échelle ne sont pas fixés. \\
% HMT-MASS-EXPLICIT-SIGNATURE:SM-Q-D1
$d/\bar d$ & $4.679550\allowbreak{}162948\allowbreak{}874\;\mathrm{MeV}/c^2$\textsuperscript{\(\dagger\)}; carte angulaire, signature nominale de parcours (17,8) & $4.70\pm0.07\;\mathrm{MeV}/c^2$ ; masse $\overline{\mathrm{MS}}$ à $\mu=2\,\mathrm{GeV}$ ; 90\% C.L. & Différence entre les valeurs centrales : $-4351.029159\allowbreak{}814\;\mathrm{ppm}$\textsuperscript{\(\dagger\)}. Il ne s’agit pas d’un résidu métrologique tant que le même schéma et la même échelle ne sont pas fixés. \\
% HMT-MASS-EXPLICIT-SIGNATURE:SM-Q-U2
$c/\bar c$ & $1.270500\allowbreak{}838641\allowbreak{}911\;\mathrm{GeV}/c^2$\textsuperscript{\(\dagger\)}; carte angulaire, signature nominale de parcours (61,8) & $1.2729\pm0.0045\;\mathrm{GeV}/c^2$ ; masse $\overline{\mathrm{MS}}$ évaluée à $\mu=\overline{m}_q$ ; 90\% C.L. & Différence entre les valeurs centrales : $-1884.799558\allowbreak{}558\;\mathrm{ppm}$\textsuperscript{\(\dagger\)}. Il ne s’agit pas d’un résidu métrologique tant que le même schéma et la même échelle ne sont pas fixés. \\
% HMT-MASS-EXPLICIT-SIGNATURE:SM-Q-D2
$s/\bar s$ & $92.940093\allowbreak{}907845\allowbreak{}64\;\mathrm{MeV}/c^2$\textsuperscript{\(\dagger\)}; carte angulaire, signature nominale de parcours (41,-3) & $92.9\pm0.7\;\mathrm{MeV}/c^2$ ; masse $\overline{\mathrm{MS}}$ à $\mu=2\,\mathrm{GeV}$ ; 90\% C.L. & Différence entre les valeurs centrales : $431.581354\allowbreak{}635\;\mathrm{ppm}$\textsuperscript{\(\dagger\)}. Il ne s’agit pas d’un résidu métrologique tant que le même schéma et la même échelle ne sont pas fixés. \\
% HMT-MASS-EXPLICIT-SIGNATURE:SM-Q-U3
$t/\bar t$ & $172.597479\allowbreak{}544019\allowbreak{}1\;\mathrm{GeV}/c^2$\textsuperscript{\(\dagger\)}; carte angulaire, signature nominale de parcours (100,-1) & $172.60\pm0.27\;\mathrm{GeV}/c^2$ ; masse reconstruite directement à partir de la cinématique des événements & Différence entre les valeurs centrales : $-14.602873\allowbreak{}585\;\mathrm{ppm}$\textsuperscript{\(\dagger\)}. Il ne s’agit pas d’un résidu métrologique tant que le même schéma et la même échelle ne sont pas fixés. \\
% HMT-MASS-EXPLICIT-SIGNATURE:SM-Q-D3
$b/\bar b$ & $4.190119\allowbreak{}763299\allowbreak{}790\;\mathrm{GeV}/c^2$\textsuperscript{\(\dagger\)}; carte angulaire, signature nominale de parcours (71,-5) & $4.186\pm0.006\;\mathrm{GeV}/c^2$ ; masse $\overline{\mathrm{MS}}$ évaluée à $\mu=\overline{m}_q$ ; 90\% C.L. & Différence entre les valeurs centrales : $984.176612\allowbreak{}467\;\mathrm{ppm}$\textsuperscript{\(\dagger\)}. Il ne s’agit pas d’un résidu métrologique tant que le même schéma et la même échelle ne sont pas fixés. \\
% HMT-MASS-EXPLICIT-SIGNATURE:SM-GA-W
$W^\pm$ & $80.378964\allowbreak{}909704\allowbreak{}547024\allowbreak{}86\;\mathrm{GeV}/c^2$\textsuperscript{\(\dagger\)}; tour E3, signature (94,-1,0,-2,4) & $80.3625\pm0.0077\;\mathrm{GeV}/c^2$ ; moyenne de masse de l’instantané PDG scellé & Différence relative : $204.882995\allowbreak{}234\;\mathrm{ppm}$\textsuperscript{\(\dagger\)} ; $2.138299\,\sigma$. La valeur HMT est l’évaluation exacte de la signature déclarée. \\
% HMT-MASS-EXPLICIT-SIGNATURE:SM-GA-Z
$Z^0$ & $91.187533\allowbreak{}444313\allowbreak{}407844\allowbreak{}85\;\mathrm{GeV}/c^2$\textsuperscript{\(\dagger\)}; tour E3, signature (95,-1,-11,0,-4) & $91.1879\pm0.0020\;\mathrm{GeV}/c^2$ ; moyenne de masse de l’instantané PDG scellé & Différence relative : $-4.019784\allowbreak{}276\;\mathrm{ppm}$\textsuperscript{\(\dagger\)} ; $-0.186362\,\sigma$. La valeur HMT est l’évaluation exacte de la signature déclarée. \\
% HMT-MASS-EXPLICIT-SIGNATURE:SM-H-1
Higgs $H$ & $125.292466\allowbreak{}042536\allowbreak{}133\;\mathrm{GeV}/c^2$\textsuperscript{\(\dagger\)}; carte angulaire, signature nominale de parcours (97,9) & $125.13\pm0.11\;\mathrm{GeV}/c^2$ ; moyenne de masse de l’instantané PDG scellé & Différence relative : $1298.378027\allowbreak{}140\;\mathrm{ppm}$\textsuperscript{\(\dagger\)} ; $1.454206\,\sigma$. La valeur HMT est l’évaluation exacte de la signature déclarée. \\
% HMT-MASS-EXPLICIT-SIGNATURE:E3B-002
$\pi^0$\par signature $(44,\allowbreak -4,\allowbreak -25,\allowbreak 1,\allowbreak -2)$ & $134.976724\allowbreak{}327872\allowbreak{}125739\allowbreak{}384878\ldots\;\mathrm{MeV}/c^2$\textsuperscript{\(\dagger\)} & $134.976827\allowbreak{}767684\allowbreak{}71\pm 0.000485\allowbreak{}679282\allowbreak{}284233\allowbreak{}51\;\mathrm{MeV}/c^2$\par\footnotesize Catalogue PDG 2026. & Différence relative $-0.766352\allowbreak{}375\,\mathrm{ppm}$ ; $-0.212979\,\sigma$. Pour la signature composée déclarée, l’évaluation se déduit exactement du caractère bêta E3. Le sélecteur prospectif de cette signature est une application antérieure à la comparaison numérique et indépendante de celle-ci. \\
% HMT-MASS-EXPLICIT-SIGNATURE:E3B-003
$\pi^\pm$\par signature $(44,\allowbreak 1,\allowbreak -2,\allowbreak 2,\allowbreak -1)$ & $139.570485\allowbreak{}171572\allowbreak{}191374\allowbreak{}734364\ldots\;\mathrm{MeV}/c^2$\textsuperscript{\(\dagger\)} & $139.570390\allowbreak{}983681\allowbreak{}31\pm 0.000182\allowbreak{}007160\allowbreak{}408260\allowbreak{}09\;\mathrm{MeV}/c^2$\par\footnotesize Catalogue PDG 2026. & Différence relative $0.674841\allowbreak{}491\,\mathrm{ppm}$ ; $0.517495\,\sigma$. Pour la signature composée déclarée, l’évaluation se déduit exactement du caractère bêta E3. Le sélecteur prospectif de cette signature est une application antérieure à la comparaison numérique et indépendante de celle-ci. \\
% HMT-MASS-EXPLICIT-SIGNATURE:E3B-004
$K^\pm$\par signature $(54,\allowbreak -1,\allowbreak 17,\allowbreak -2,\allowbreak 2)$ & $493.677085\allowbreak{}649530\allowbreak{}612475\allowbreak{}723054\ldots\;\mathrm{MeV}/c^2$\textsuperscript{\(\dagger\)} & $493.676599\allowbreak{}458040\allowbreak{}58\pm 0.015405\allowbreak{}665226\allowbreak{}071509\;\mathrm{MeV}/c^2$\par\footnotesize Catalogue PDG 2026. & Différence relative $0.984838\allowbreak{}030\,\mathrm{ppm}$ ; $0.031559\,\sigma$. Pour la signature composée déclarée, l’évaluation se déduit exactement du caractère bêta E3. Le sélecteur prospectif de cette signature est une application antérieure à la comparaison numérique et indépendante de celle-ci. \\
% HMT-MASS-EXPLICIT-SIGNATURE:E3B-005
$K^0$\par signature $(54,\allowbreak 0,\allowbreak 32,\allowbreak 3,\allowbreak -2)$ & $497.611186\allowbreak{}497750\allowbreak{}457358\allowbreak{}558672\ldots\;\mathrm{MeV}/c^2$\textsuperscript{\(\dagger\)} & $497.610767\allowbreak{}531257\allowbreak{}52\pm 0.012990\allowbreak{}199756\allowbreak{}4242\;\mathrm{MeV}/c^2$\par\footnotesize Catalogue PDG 2026. & Différence relative $0.841956\allowbreak{}244\,\mathrm{ppm}$ ; $0.032252\,\sigma$. Pour la signature composée déclarée, l’évaluation se déduit exactement du caractère bêta E3. Le sélecteur prospectif de cette signature est une application antérieure à la comparaison numérique et indépendante de celle-ci. \\
% HMT-MASS-EXPLICIT-SIGNATURE:E3B-006
proton $p$\par signature $(59,\allowbreak 0,\allowbreak 9,\allowbreak 1,\allowbreak 0)$ & $938.271751\allowbreak{}546984\allowbreak{}898298\allowbreak{}936787\ldots\;\mathrm{MeV}/c^2$\textsuperscript{\(\dagger\)} & $938.272089\allowbreak{}430000\allowbreak{}05\pm 0.000000\allowbreak{}290000\allowbreak{}000000\allowbreak{}00\;\mathrm{MeV}/c^2$\par\footnotesize Catalogue PDG 2026. & Différence relative $-0.360111\allowbreak{}974\,\mathrm{ppm}$ ; $-1165.113845\,\sigma$. Pour la signature composée déclarée, l’évaluation se déduit exactement du caractère bêta E3. Le sélecteur prospectif de cette signature est une application antérieure à la comparaison numérique et indépendante de celle-ci. \\
% HMT-MASS-EXPLICIT-SIGNATURE:E3B-007
neutron $n$\par signature $(59,\allowbreak 1,\allowbreak -34,\allowbreak 0,\allowbreak 1)$ & $939.565810\allowbreak{}472507\allowbreak{}023312\allowbreak{}664431\ldots\;\mathrm{MeV}/c^2$\textsuperscript{\(\dagger\)} & $939.565421\allowbreak{}939999\allowbreak{}96\pm 0.000000\allowbreak{}480000\allowbreak{}000000\allowbreak{}00\;\mathrm{MeV}/c^2$\par\footnotesize Catalogue PDG 2026. & Différence relative $0.413523\allowbreak{}633\,\mathrm{ppm}$ ; $809.442723\,\sigma$. Pour la signature composée déclarée, l’évaluation se déduit exactement du caractère bêta E3. Le sélecteur prospectif de cette signature est une application antérieure à la comparaison numérique et indépendante de celle-ci. \\
\end{longtable}
\endgroup


\Needspace{12\baselineskip}
\paragraph{Invariants natifs des secteurs topologiques et virtuels.}
Les secteurs Fibonacci, Ising, \(\mathbb Z/6\) et les sections virtuelles sont comparés
dans leurs codomaines propres. Fibonacci donne l'anneau basé et sa dimension de
Frobenius--Perron ; les données \(F/R\) demeurent dans la réalisation ultérieure. Ising
conserve son anneau de fusion et distingue la réalisation de Majorana ; \(\mathbb Z/6\)
donne ses caractères de tressage ; la virtualité donne des horizons de
persistance. Ce
tableau maintient ces invariants dans la même application des réalisations sans les traduire
artificiellement en unité de masse.

\begingroup
\setlength{\tabcolsep}{4.5pt}
\renewcommand{\arraystretch}{1.28}
\begin{longtable}{@{}>{\raggedright\arraybackslash}p{.15\textwidth}>{\raggedright\arraybackslash}p{.31\textwidth}>{\raggedright\arraybackslash}p{.25\textwidth}>{\raggedright\arraybackslash}p{.23\textwidth}@{}}
\caption{Secteurs topologiques et virtuels : invariants natifs de fusion, de tressage et de persistance.}\label{tab:masas-invariantes-topologicos-virtuales-rev3}\\
\toprule
\textbf{Secteur} & \textbf{Objet et relations} & \textbf{Invariant natif} & \textbf{Portée physique} \\
\midrule
\endfirsthead
\toprule
\textbf{Secteur} & \textbf{Objet et relations} & \textbf{Invariant natif} & \textbf{Portée physique} \\
\midrule
\endhead
\midrule
\multicolumn{4}{r}{\footnotesize Suite à la page suivante.}\\
\endfoot
\bottomrule
\endlastfoot
Fibonacci & $K=\mathbb Z1\oplus\mathbb Z\tau$, $\tau^2=1+\tau$. & $\operatorname{FPdim}(\tau)=\varphi$ dans l’anneau basé. & Les matrices $F/R$, le pentagone, l’hexagone, l’Hamiltonien et le gap relèvent de la réalisation ultérieure. \\
Ising & $K=\mathbb Z\{1,\psi,\sigma\}$, $\psi^2=1$, $\psi\sigma=\sigma$, $\sigma^2=1+\psi$. & $\operatorname{FPdim}(\sigma)=\sqrt{2}$ ; représentation de tressage. & La réalisation de Majorana conserve son propre morphisme physique. \\
$\mathbb Z/6$ & $q_{\epsilon,t}=(-1)^\epsilon\omega_3^t=\zeta_6^k$. & Caractère quadratique et représentation de tressage sur les six secteurs. & L’observation pertinente est topologique, non une masse universelle. \\
Secciones virtuales & Système inverse de prolongements d’horizon $L(s_N)$. & Enregistrement fini, signature, retenue et classe d’obstruction. & L’observable natif est la persistance du prolongement. \\
\end{longtable}
\endgroup


Les neutrinos demeurent dans le codomaine des valeurs propres et du mélange décrit par leurs
fiches opératorielles ; PMNS réalise le transport entre les bases d'interaction et de
masse, tandis que l'échelle appartient au spectre de l'opérateur neutrinique. Les six
évaluations composées s'obtiennent à partir de signatures explicites au moyen du morphisme
général de liaison vers la carte scalaire. Les secteurs topologiques et les sections virtuelles maintiennent leurs
observables propres de fusion, de tressage et de persistance.

\subsection{Construction causale et couverture des classes de masse}

L'architecture APP--TRIT--TPK, la loi des masses, le lecteur bidirectionnel et
les tours déterminent conjointement l'atlas, les enregistrements chronologiques
enrichis, les signatures, le caractère et la clôture de base. Leur convergence dans
\eqref{eq:ley-operatorial-general-masas-rectoras} fournit une unique
équation de sections de masse. Le domaine complet est l'atlas \(81/56/13/324\), et chaque réalisation spécifie, selon son codomaine,
et spécifie, selon le codomaine de chaque secteur, classificateurs, cellules, parcours,
orientations, mots dodécaphasiques, signatures et morphismes de compatibilité. Les
spectres \(K_D/K_\Omega\) raffinent les fibres de l'atlas, tandis que les parcours
chronologiques de \(108\) états associés à \(W\), \(Z\) et \(H\) construisent
directement leurs signatures et caractères bosoniques.

\endgroup
\begingroup
\let\chapter\section
\let\section\subsection
\let\subsection\subsubsection
\let\subsubsection\paragraph
% Comparación trasladada detrás de la ley y de la ontología sectorial.
% Se excluyen únicamente la portada autónoma y el metalenguaje de edición.
\section{Évaluation quantitative postérieure à la fixation du graphe génératif}
Le tableau suivant place dans des colonnes contiguës toute coordonnée scalaire que la théorie permet d’évaluer de manière reproductible. La référence expérimentale n’est consultée qu’après fixation de la sortie HMT--MD. Une ligne intitulée évaluation conditionnelle énonce exactement la valeur de la fonction une fois la signature fixée ; elle n’attribue pas à cette signature un sélecteur physique qui n’aurait pas été construit. Les quarks sont présentés sans résidu, car une comparaison numérique exige la concordance du schéma et de l’échelle de renormalisation.

Le dénombrement numérique de ce tableau est exhaustif relativement aux coordonnées scalaires actuellement reproductibles : un électron bêta fermé, dix-sept évaluations conditionnelles --- onze élémentaires et six composées --- et deux sorties nulles exactes. Il contient donc vingt lignes numériques. Ce nombre n’est pas le domaine de la loi. Les classes dont la sortie naturelle demeure un opérateur, un spectre, une orbite ou une multisection apparaissent ensuite avec cet objet complet, sans remplir leur colonne par la référence externe. L’inventaire expérimental de 471 masses et de 384 largeurs est conservé intégralement comme comparaison état par état ; son cardinal n’est pas présenté comme un dénombrement des prédictions HMT.

\Needspace{7\baselineskip}
\begingroup
\footnotesize
\setlength{\tabcolsep}{2.1pt}
\begin{longtable}{L{1.17cm}L{2.65cm}L{3.12cm}L{2.18cm}L{5.69cm}}
\toprule
\rowcolor{HMTNavy}\HMTtablehead{État} & \HMTtablehead{Valeur HMT reproductible (MeV/\allowbreak{}c\({}^{2}\))} & \HMTtablehead{Referencia posterior} & \HMTtablehead{Différence relative (ppm)} & \HMTtablehead{Nature et portée} \\
\midrule
\endfirsthead
\multicolumn{5}{l}{\sffamily\scriptsize\color{HMTGray}Suite du tableau}\\[-1pt]
\toprule
\rowcolor{HMTNavy}\HMTtablehead{État} & \HMTtablehead{Valeur HMT reproductible (MeV/\allowbreak{}c\({}^{2}\))} & \HMTtablehead{Referencia posterior} & \HMTtablehead{Différence relative (ppm)} & \HMTtablehead{Nature et portée} \\
\midrule
\endhead
\midrule
\endfoot
\bottomrule
\endlastfoot
e-/\allowbreak{}e+ & 0.510998\allowbreak{}950690\allowbreak{}000809 & 0.510998\allowbreak{}95069±\allowbreak{}0.000000\allowbreak{}00016 MeV/\allowbreak{}c\({}^{2}\) & 1.582406\allowbreak{}884×10\({}^{-9}\) & coordonnée bêta fermée dans la fibre centrale \\
\rowcolor{HMTLight}μ-/\allowbreak{}μ+ & 105.658360\allowbreak{}294435\allowbreak{}829303 & 105.658375\allowbreak{}5±\allowbreak{}0.000002\allowbreak{}3 MeV/\allowbreak{}c\({}^{2}\) & -0.143912\allowbreak{}53 & raffinement E3 conditionné par la signature \\
τ-/\allowbreak{}τ+ & 1776.860917\allowbreak{}631432\allowbreak{}295595 & 1776.932465\allowbreak{}134090\allowbreak{}1 MeV/\allowbreak{}c\({}^{2}\) (valeur centrale utilisée ; présentation arrondie : 1776.93±\allowbreak{}0.09 MeV/\allowbreak{}c\({}^{2}\)) & -40.264615\allowbreak{}601 & raffinement E3 conditionné par la signature \\
\rowcolor{HMTLight}u/\allowbreak{}\(\bar{\mathrm{u}}\) & 2.165924\allowbreak{}536173\allowbreak{}907663 & 2.16±\allowbreak{}0.07 MeV/\allowbreak{}c\({}^{2}\) & --- & carte angulaire ; aucune signature Z n’a été déclarée\({}^{5}\) ni aucun raffinement E3 \\
d/\allowbreak{}\(\bar{\mathrm{d}}\) & 4.679550\allowbreak{}162948\allowbreak{}874172 & 4.70±\allowbreak{}0.07 MeV/\allowbreak{}c\({}^{2}\) & --- & carte angulaire ; aucune signature Z n’a été déclarée\({}^{5}\) ni aucun raffinement E3 \\
\rowcolor{HMTLight}c/\allowbreak{}\(\bar{\mathrm{c}}\) & 1270.500838\allowbreak{}641911\allowbreak{}135286 & 1.2729±\allowbreak{}0.0045 GeV/\allowbreak{}c\({}^{2}\) & --- & carte angulaire ; aucune signature Z n’a été déclarée\({}^{5}\) ni aucun raffinement E3 \\
s/\allowbreak{}\(\bar{\mathrm{s}}\) & 92.940093\allowbreak{}907845\allowbreak{}640641 & 92.9 ±\allowbreak{}0.7 MeV/\allowbreak{}c\({}^{2}\) & --- & carte angulaire ; aucune signature Z n’a été déclarée\({}^{5}\) ni aucun raffinement E3 \\
\rowcolor{HMTLight}t/\allowbreak{}\(\bar{\mathrm{t}}\) & 172597.479544\allowbreak{}019136\allowbreak{}261367 & 172.60±\allowbreak{}0.27 GeV/\allowbreak{}c\({}^{2}\) & --- & carte angulaire ; aucune signature Z n’a été déclarée\({}^{5}\) ni aucun raffinement E3 \\
b/\allowbreak{}\(\bar{\mathrm{b}}\) & 4190.119763\allowbreak{}299790\allowbreak{}218075 & 4.186±\allowbreak{}0.006 GeV/\allowbreak{}c\({}^{2}\) & --- & carte angulaire ; aucune signature Z n’a été déclarée\({}^{5}\) ni aucun raffinement E3 \\
\rowcolor{HMTLight}γ & 0 & <1×10\({}^{-18}\) eV/\allowbreak{}c\({}^{2}\) & --- & sortie nulle exacte \\
g & 0 & --- & --- & sortie nulle exacte \\
\rowcolor{HMTLight}W+/\allowbreak{}W- & 80378.964909\allowbreak{}704547\allowbreak{}024865 & 80.3625 ±\allowbreak{}0.0077 GeV/\allowbreak{}c\({}^{2}\) & 204.882995\allowbreak{}235 & raffinement E3 conditionné par la signature \\
Z0 & 91187.533444\allowbreak{}313407\allowbreak{}844855 & 91.187873\allowbreak{}297221\allowbreak{}344 GeV/\allowbreak{}c\({}^{2}\) (valeur centrale utilisée ; présentation arrondie : 91.1879±\allowbreak{}0.0020 GeV/\allowbreak{}c\({}^{2}\)) & -3.726952\allowbreak{}89 & raffinement E3 conditionné par la signature \\
\rowcolor{HMTLight}H & 125292.466042\allowbreak{}536133\allowbreak{}000433 & 125.130943\allowbreak{}828161\allowbreak{}51 GeV/\allowbreak{}c\({}^{2}\) (valeur centrale utilisée ; présentation arrondie : 125.13±\allowbreak{}0.11 GeV/\allowbreak{}c\({}^{2}\)) & 1290.825509\allowbreak{}927 & carte angulaire ; aucune signature Z n’a été déclarée\({}^{5}\) ni aucun raffinement E3 \\
π\({}^{0}\) & 134.976724\allowbreak{}327872\allowbreak{}125739 & 134.976827\allowbreak{}767685 MeV/\allowbreak{}c\({}^{2}\) & -0.766352\allowbreak{}378 & raffinement E3 conditionné par la signature composée \\
\rowcolor{HMTLight}π±\allowbreak{} & 139.570485\allowbreak{}171572\allowbreak{}191375 & 139.570390\allowbreak{}983681 MeV/\allowbreak{}c\({}^{2}\) & 0.674841\allowbreak{}494 & raffinement E3 conditionné par la signature composée \\
K±\allowbreak{} & 493.677085\allowbreak{}649530\allowbreak{}612476 & 493.676599\allowbreak{}458041 MeV/\allowbreak{}c\({}^{2}\) & 0.984838\allowbreak{}03 & raffinement E3 conditionné par la signature composée \\
\rowcolor{HMTLight}K\({}^{0}\) & 497.611186\allowbreak{}497750\allowbreak{}457359 & 497.610767\allowbreak{}531258 MeV/\allowbreak{}c\({}^{2}\) & 0.841956\allowbreak{}243 & raffinement E3 conditionné par la signature composée \\
proton & 938.271751\allowbreak{}546984\allowbreak{}898299 & 938.272089\allowbreak{}43 MeV/\allowbreak{}c\({}^{2}\) & -0.360111\allowbreak{}975 & raffinement E3 conditionné par la signature composée \\
\rowcolor{HMTLight}neutron & 939.565810\allowbreak{}472507\allowbreak{}023313 & 939.565421\allowbreak{}94 MeV/\allowbreak{}c\({}^{2}\) & 0.413523\allowbreak{}633 & raffinement E3 conditionné par la signature composée \\
\end{longtable}
\endgroup
\Needspace{7\baselineskip}
\section{Évaluation individuelle de \(20\) sections matérielles et traçabilité de leurs coordonnées métrologiques}
\Needspace{7\baselineskip}
\subsection{Paire électron–positron : parcours central, conjugaison des feuillets et caractère de masse}
\begin{itemize}
\item \textbf{Identifiant du catalogue :} SM-LC-1.
\item \textbf{Type de descripteur :} sélecteur basal et lecteurs bidirectionnels.
\item \textbf{Descripteur :} sélecteur (-22,+15); K\_D=\allowbreak{}K\_Ω=\allowbreak{}(80,54,6).
\item \textbf{Étape :} bêta bidirectionnel ; coordonnée initiale 0.510998\allowbreak{}922488\allowbreak{}066073 MeV/c\({}^{2}\).
\item \textbf{Évaluation HMT--MD reproductible :} 0.510998\allowbreak{}950690\allowbreak{}000809 MeV/c\({}^{2}\).
\item \textbf{Nature mathématique :} coordonnée bêta fermée dans la fibre centrale.
\item \textbf{Référence ultérieure :} 0.510998\allowbreak{}95069±\allowbreak{}0.000000\allowbreak{}00016 MeV/c\({}^{2}\).
\item \textbf{Identité et métrologie externes :} Particle Data Group (PDG); PDG2026:\allowbreak{}observable:\allowbreak{}S003:\allowbreak{}S003M:\allowbreak{}332838; schéma non conservé dans l’extrait documentaire; échelle non conservée dans l’extrait documentaire; niveau de confiance non conservé dans l’extrait documentaire.
\item \textbf{Différence relative:} 1.582406\allowbreak{}884×10\({}^{-9}\) ppm.
\item \textbf{Portée :} dérivation scalaire fermée.
\end{itemize}
\Needspace{7\baselineskip}
\subsection{Paire muon–antimuon : conjugaison des parcours et caractère leptonique}
\begin{itemize}
\item \textbf{Identifiant du catalogue :} SM-LC-2.
\item \textbf{Type de descripteur :} signature complète σ\(\in\)Z\({}^{5}\).
\item \textbf{Descripteur :} (42,-3,8,0,2).
\item \textbf{Étape :} tour bêta E3 ; coordonnée initiale 105.564059\allowbreak{}012859\allowbreak{}26472 MeV/c\({}^{2}\).
\item \textbf{Évaluation HMT--MD reproductible :} 105.658360\allowbreak{}294435\allowbreak{}829303 MeV/c\({}^{2}\).
\item \textbf{Nature mathématique :} raffinement E3 conditionné par la signature.
\item \textbf{Référence ultérieure :} 105.658375\allowbreak{}5±\allowbreak{}0.000002\allowbreak{}3 MeV/c\({}^{2}\).
\item \textbf{Identité et métrologie externes :} Particle Data Group (PDG); PDG2026:\allowbreak{}observable:\allowbreak{}S004:\allowbreak{}S004M:\allowbreak{}332849; schéma non conservé dans l’extrait documentaire; échelle non conservée dans l’extrait documentaire; niveau de confiance non conservé dans l’extrait documentaire.
\item \textbf{Différence relative:} -0.143912\allowbreak{}53 ppm.
\item \textbf{Portée :} évaluation exacte une fois la signature fixée ; la sélection physique de cette signature constitue une opération indépendante.
\end{itemize}
\Needspace{7\baselineskip}
\subsection{Paire tau–antitau : conjugaison des parcours et caractère leptonique}
\begin{itemize}
\item \textbf{Identifiant du catalogue :} SM-LC-3.
\item \textbf{Type de descripteur :} signature complète σ\(\in\)Z\({}^{5}\).
\item \textbf{Descripteur :} (64,0,25,-1,6).
\item \textbf{Étape :} tour bêta E3 ; coordonnée initiale 1771.945393\allowbreak{}886234\allowbreak{}110148 MeV/c\({}^{2}\).
\item \textbf{Évaluation HMT--MD reproductible :} 1776.860917\allowbreak{}631432\allowbreak{}295595 MeV/c\({}^{2}\).
\item \textbf{Nature mathématique :} raffinement E3 conditionné par la signature.
\item \textbf{Référence ultérieure :} 1776.932465\allowbreak{}134090\allowbreak{}1 MeV/c\({}^{2}\) (valeur centrale utilisée ; présentation arrondie : 1776.93±\allowbreak{}0.09 MeV/c\({}^{2}\)).
\item \textbf{Identité et métrologie externes :} Particle Data Group (PDG); PDG2026:\allowbreak{}observable:\allowbreak{}S035:\allowbreak{}S035M:\allowbreak{}332893; schéma non conservé dans l’extrait documentaire; échelle non conservée dans l’extrait documentaire; niveau de confiance non conservé dans l’extrait documentaire.
\item \textbf{Différence relative:} -40.264615\allowbreak{}601 ppm.
\item \textbf{Portée :} évaluation exacte une fois la signature fixée ; la sélection physique de cette signature constitue une opération indépendante.
\end{itemize}
\Needspace{7\baselineskip}
\subsection{Quark haut et antiquark haut}
\begin{itemize}
\item \textbf{Identifiant du catalogue :} SM-Q-U1.
\item \textbf{Type de descripteur :} coindexation angulaire (n\_A,s\_C)\(\in\)Z\({}^{2}\).
\item \textbf{Descripteur :} (11,7).
\item \textbf{Étape :} carte angulaire conditionnelle ; coordonnée initiale 2.165924\allowbreak{}536173\allowbreak{}907663 MeV/c\({}^{2}\).
\item \textbf{Évaluation HMT--MD reproductible :} 2.165924\allowbreak{}536173\allowbreak{}907663 MeV/c\({}^{2}\).
\item \textbf{Nature mathématique :} carte angulaire ; aucune signature Z n’a été déclarée\({}^{5}\) ni aucun raffinement E3.
\item \textbf{Référence ultérieure :} 2.16±\allowbreak{}0.07 MeV/c\({}^{2}\).
\item \textbf{Identité et métrologie externes :} Particle Data Group (PDG); PDG2026:\allowbreak{}observable:\allowbreak{}Q002:\allowbreak{}Q002M:\allowbreak{}333588; schéma non conservé dans l’extrait documentaire; échelle non conservée dans l’extrait documentaire; niveau de confiance 90\%.
\item \textbf{Différence relative:} --- ppm.
\item \textbf{Portée :} évaluation conditionnelle ; la signature n’a pas été choisie à partir de la valeur comparée ; schéma et échelle externes requis pour le résidu.
\end{itemize}
\Needspace{7\baselineskip}
\subsection{Quark bas et antiquark bas}
\begin{itemize}
\item \textbf{Identifiant du catalogue :} SM-Q-D1.
\item \textbf{Type de descripteur :} coindexation angulaire (n\_A,s\_C)\(\in\)Z\({}^{2}\).
\item \textbf{Descripteur :} (17,8).
\item \textbf{Étape :} carte angulaire conditionnelle ; coordonnée initiale 4.679550\allowbreak{}162948\allowbreak{}874172 MeV/c\({}^{2}\).
\item \textbf{Évaluation HMT--MD reproductible :} 4.679550\allowbreak{}162948\allowbreak{}874172 MeV/c\({}^{2}\).
\item \textbf{Nature mathématique :} carte angulaire ; aucune signature Z n’a été déclarée\({}^{5}\) ni aucun raffinement E3.
\item \textbf{Référence ultérieure :} 4.70±\allowbreak{}0.07 MeV/c\({}^{2}\).
\item \textbf{Identité et métrologie externes :} Particle Data Group (PDG); PDG2026:\allowbreak{}observable:\allowbreak{}Q001:\allowbreak{}Q001M:\allowbreak{}333589; schéma non conservé dans l’extrait documentaire; échelle non conservée dans l’extrait documentaire; niveau de confiance 90\%.
\item \textbf{Différence relative:} --- ppm.
\item \textbf{Portée :} évaluation conditionnelle ; la signature n’a pas été choisie à partir de la valeur comparée ; schéma et échelle externes requis pour le résidu.
\end{itemize}
\Needspace{7\baselineskip}
\subsection{Quark charme et antiquark charme}
\begin{itemize}
\item \textbf{Identifiant du catalogue :} SM-Q-U2.
\item \textbf{Type de descripteur :} coindexation angulaire (n\_A,s\_C)\(\in\)Z\({}^{2}\).
\item \textbf{Descripteur :} (61,8).
\item \textbf{Étape :} carte angulaire conditionnelle ; coordonnée initiale 1270.500838\allowbreak{}641911\allowbreak{}135286 MeV/c\({}^{2}\).
\item \textbf{Évaluation HMT--MD reproductible :} 1270.500838\allowbreak{}641911\allowbreak{}135286 MeV/c\({}^{2}\).
\item \textbf{Nature mathématique :} carte angulaire ; aucune signature Z n’a été déclarée\({}^{5}\) ni aucun raffinement E3.
\item \textbf{Référence ultérieure :} 1.2729±\allowbreak{}0.0045 GeV/c\({}^{2}\).
\item \textbf{Identité et métrologie externes :} Particle Data Group (PDG); PDG2026:\allowbreak{}observable:\allowbreak{}Q004:\allowbreak{}Q004M:\allowbreak{}333604; schéma non conservé dans l’extrait documentaire; échelle non conservée dans l’extrait documentaire; niveau de confiance 90\%.
\item \textbf{Différence relative:} --- ppm.
\item \textbf{Portée :} évaluation conditionnelle ; la signature n’a pas été choisie à partir de la valeur comparée ; schéma et échelle externes requis pour le résidu.
\end{itemize}
\Needspace{7\baselineskip}
\subsection{Quark étrange et antiquark étrange}
\begin{itemize}
\item \textbf{Identifiant du catalogue :} SM-Q-D2.
\item \textbf{Type de descripteur :} coindexation angulaire (n\_A,s\_C)\(\in\)Z\({}^{2}\).
\item \textbf{Descripteur :} (41,-3).
\item \textbf{Étape :} carte angulaire conditionnelle ; coordonnée initiale 92.940093\allowbreak{}907845\allowbreak{}640641 MeV/c\({}^{2}\).
\item \textbf{Évaluation HMT--MD reproductible :} 92.940093\allowbreak{}907845\allowbreak{}640641 MeV/c\({}^{2}\).
\item \textbf{Nature mathématique :} carte angulaire ; aucune signature Z n’a été déclarée\({}^{5}\) ni aucun raffinement E3.
\item \textbf{Référence ultérieure :} 92.9 ±\allowbreak{}0.7 MeV/c\({}^{2}\).
\item \textbf{Identité et métrologie externes :} Particle Data Group (PDG); PDG2026:\allowbreak{}observable:\allowbreak{}Q003:\allowbreak{}Q003M:\allowbreak{}333590; schéma non conservé dans l’extrait documentaire; échelle non conservée dans l’extrait documentaire; niveau de confiance 90\%.
\item \textbf{Différence relative:} --- ppm.
\item \textbf{Portée :} évaluation conditionnelle ; la signature n’a pas été choisie à partir de la valeur comparée ; schéma et échelle externes requis pour le résidu.
\end{itemize}
\Needspace{7\baselineskip}
\subsection{Quark sommet et antiquark sommet}
\begin{itemize}
\item \textbf{Identifiant du catalogue :} SM-Q-U3.
\item \textbf{Type de descripteur :} coindexation angulaire (n\_A,s\_C)\(\in\)Z\({}^{2}\).
\item \textbf{Descripteur :} (100,-1).
\item \textbf{Étape :} carte angulaire conditionnelle ; coordonnée initiale 172597.479544\allowbreak{}019136\allowbreak{}261367 MeV/c\({}^{2}\).
\item \textbf{Évaluation HMT--MD reproductible :} 172597.479544\allowbreak{}019136\allowbreak{}261367 MeV/c\({}^{2}\).
\item \textbf{Nature mathématique :} carte angulaire ; aucune signature Z n’a été déclarée\({}^{5}\) ni aucun raffinement E3.
\item \textbf{Référence ultérieure :} 172.60±\allowbreak{}0.27 GeV/c\({}^{2}\).
\item \textbf{Identité et métrologie externes :} Particle Data Group (PDG); PDG2026:\allowbreak{}observable:\allowbreak{}Q007:\allowbreak{}Q007TP:\allowbreak{}333614; schéma non conservé dans l’extrait documentaire; échelle non conservée dans l’extrait documentaire; niveau de confiance non conservé dans l’extrait documentaire.
\item \textbf{Différence relative:} --- ppm.
\item \textbf{Portée :} évaluation conditionnelle ; la signature n’a pas été choisie à partir de la valeur comparée ; schéma et échelle externes requis pour le résidu.
\end{itemize}
\Needspace{7\baselineskip}
\subsection{Quark fond et antiquark fond}
\begin{itemize}
\item \textbf{Identifiant du catalogue :} SM-Q-D3.
\item \textbf{Type de descripteur :} coindexation angulaire (n\_A,s\_C)\(\in\)Z\({}^{2}\).
\item \textbf{Descripteur :} (71,-5).
\item \textbf{Étape :} carte angulaire conditionnelle ; coordonnée initiale 4190.119763\allowbreak{}299790\allowbreak{}218075 MeV/c\({}^{2}\).
\item \textbf{Évaluation HMT--MD reproductible :} 4190.119763\allowbreak{}299790\allowbreak{}218075 MeV/c\({}^{2}\).
\item \textbf{Nature mathématique :} carte angulaire ; aucune signature Z n’a été déclarée\({}^{5}\) ni aucun raffinement E3.
\item \textbf{Référence ultérieure :} 4.186±\allowbreak{}0.006 GeV/c\({}^{2}\).
\item \textbf{Identité et métrologie externes :} Particle Data Group (PDG); PDG2026:\allowbreak{}observable:\allowbreak{}Q005:\allowbreak{}Q005M:\allowbreak{}333611; schéma non conservé dans l’extrait documentaire; échelle non conservée dans l’extrait documentaire; niveau de confiance 90\%.
\item \textbf{Différence relative:} --- ppm.
\item \textbf{Portée :} évaluation conditionnelle ; la signature n’a pas été choisie à partir de la valeur comparée ; schéma et échelle externes requis pour le résidu.
\end{itemize}
\Needspace{7\baselineskip}
\subsection{Descripteur spectral reproductible du photon : secteur nul, hélicité et caractère électromagnétique}
\begin{itemize}
\item \textbf{Identifiant du catalogue :} SM-GA-EM.
\item \textbf{Type de descripteur :} invariant sectoriel nul.
\item \textbf{Descripteur :} carte nulle.
\item \textbf{Étape :} carte nulle ; coordonnée initiale 0 MeV/c\({}^{2}\).
\item \textbf{Évaluation HMT--MD reproductible :} 0 MeV/c\({}^{2}\).
\item \textbf{Nature mathématique :} sortie nulle exacte.
\item \textbf{Référence ultérieure :} <1×10\({}^{-18}\) eV/c\({}^{2}\).
\item \textbf{Identité et métrologie externes :} Particle Data Group (PDG); PDG2026:\allowbreak{}observable:\allowbreak{}S000:\allowbreak{}S000M:\allowbreak{}332461; schéma non conservé dans l’extrait documentaire; échelle non conservée dans l’extrait documentaire; niveau de confiance non conservé dans l’extrait documentaire.
\item \textbf{Différence relative:} --- ppm.
\item \textbf{Portée :} exact dans le secteur nul ; la borne externe n’est pas interprétée comme une masse positive.
\item \textbf{Borne photonique :} référence externe illustrative non comparable ; l’extrait identifie l’observable, mais ne conserve pas le niveau de confiance.
\end{itemize}
\Needspace{7\baselineskip}
\subsection{Réalisation du gluon comme section de la représentation adjointe de \(SU(3)\)}
\begin{itemize}
\item \textbf{Identifiant du catalogue :} SM-GA-GL.
\item \textbf{Type de descripteur :} invariant sectoriel nul.
\item \textbf{Descripteur :} carte nulle.
\item \textbf{Étape :} carte nulle ; coordonnée initiale 0 MeV/c\({}^{2}\).
\item \textbf{Évaluation HMT--MD reproductible :} 0 MeV/c\({}^{2}\).
\item \textbf{Nature mathématique :} sortie nulle exacte.
\item \textbf{Referencia posterior:} ---.
\item \textbf{Identité et métrologie externes :} source non consignée ; observable non individualisé ; schéma non consigné ; échelle non consignée ; niveau de confiance non consigné.
\item \textbf{Différence relative:} --- ppm.
\item \textbf{Portée :} exact dans le secteur nul ; la nullité appartient à la carte de jauge de couleur.
\end{itemize}
\Needspace{7\baselineskip}
\subsection{Bosons \texorpdfstring{\(W^{+}\)}{W+} et \texorpdfstring{\(W^{-}\)}{W−}}
\begin{itemize}
\item \textbf{Identifiant du catalogue :} SM-GA-W.
\item \textbf{Type de descripteur :} signature complète σ\(\in\)Z\({}^{5}\).
\item \textbf{Descripteur :} (94,-1,0,-2,4).
\item \textbf{Étape :} tour bêta E3 ; coordonnée initiale 80381.592550\allowbreak{}220666\allowbreak{}304144 MeV/c\({}^{2}\).
\item \textbf{Évaluation HMT--MD reproductible :} 80378.964909\allowbreak{}704547\allowbreak{}024865 MeV/c\({}^{2}\).
\item \textbf{Nature mathématique :} raffinement E3 conditionné par la signature.
\item \textbf{Référence ultérieure :} 80.3625 ±\allowbreak{}0.0077 GeV/c\({}^{2}\).
\item \textbf{Identité et métrologie externes :} Particle Data Group (PDG); PDG2026:\allowbreak{}observable:\allowbreak{}S043:\allowbreak{}S043M:\allowbreak{}332464; schéma non conservé dans l’extrait documentaire; échelle non conservée dans l’extrait documentaire; niveau de confiance non conservé dans l’extrait documentaire.
\item \textbf{Différence relative:} 204.882995\allowbreak{}235 ppm.
\item \textbf{Portée :} évaluation exacte une fois la signature fixée ; la sélection physique de cette signature constitue une opération indépendante.
\item \textbf{Convention résonante :} l’extrait externe ne fixe pas conjointement la convention de pôle, de Breit--Wigner ou sur couche de masse, la largeur et la covariance ; le résidu n’est interprété que comme information scalaire.
\end{itemize}
\Needspace{7\baselineskip}
\subsection{Descripteur spectral reproductible du boson \(Z^0\) : mélange neutre, polarisation et caractère électrofaible}
\begin{itemize}
\item \textbf{Identifiant du catalogue :} SM-GA-Z.
\item \textbf{Type de descripteur :} signature complète σ\(\in\)Z\({}^{5}\).
\item \textbf{Descripteur :} (95,-1,-11,0,-4).
\item \textbf{Étape :} tour bêta E3 ; coordonnée initiale 91299.748286\allowbreak{}598170\allowbreak{}117317 MeV/c\({}^{2}\).
\item \textbf{Évaluation HMT--MD reproductible :} 91187.533444\allowbreak{}313407\allowbreak{}844855 MeV/c\({}^{2}\).
\item \textbf{Nature mathématique :} raffinement E3 conditionné par la signature.
\item \textbf{Référence ultérieure :} 91.187873\allowbreak{}297221\allowbreak{}344 GeV/c\({}^{2}\) (valeur centrale utilisée ; présentation arrondie : 91.1879±\allowbreak{}0.0020 GeV/c\({}^{2}\)).
\item \textbf{Identité et métrologie externes :} Particle Data Group (PDG); PDG2026:\allowbreak{}observable:\allowbreak{}S044:\allowbreak{}S044M:\allowbreak{}332513; schéma non conservé dans l’extrait documentaire; échelle non conservée dans l’extrait documentaire; niveau de confiance non conservé dans l’extrait documentaire.
\item \textbf{Différence relative:} -3.726952\allowbreak{}89 ppm.
\item \textbf{Portée :} évaluation exacte une fois la signature fixée ; la sélection physique de cette signature constitue une opération indépendante.
\item \textbf{Convention résonante :} l’extrait externe ne fixe pas conjointement la convention de pôle, de Breit--Wigner ou sur couche de masse, la largeur et la covariance ; le résidu n’est interprété que comme information scalaire.
\end{itemize}
\Needspace{7\baselineskip}
\subsection{Réalisation du boson de Higgs : mode scalaire, couplage et caractère de masse}
\begin{itemize}
\item \textbf{Identifiant du catalogue :} SM-H-1.
\item \textbf{Type de descripteur :} coindexation angulaire (n\_A,s\_C)\(\in\)Z\({}^{2}\).
\item \textbf{Descripteur :} (97,9).
\item \textbf{Étape :} carte angulaire conditionnelle ; coordonnée initiale 125292.466042\allowbreak{}536133\allowbreak{}000433 MeV/c\({}^{2}\).
\item \textbf{Évaluation HMT--MD reproductible :} 125292.466042\allowbreak{}536133\allowbreak{}000433 MeV/c\({}^{2}\).
\item \textbf{Nature mathématique :} carte angulaire ; aucune signature Z n’a été déclarée\({}^{5}\) ni aucun raffinement E3.
\item \textbf{Référence ultérieure :} 125.130943\allowbreak{}828161\allowbreak{}51 GeV/c\({}^{2}\) (valeur centrale utilisée ; présentation arrondie : 125.13±\allowbreak{}0.11 GeV/c\({}^{2}\)).
\item \textbf{Identité et métrologie externes :} Particle Data Group (PDG); PDG2026:\allowbreak{}observable:\allowbreak{}S126:\allowbreak{}S126M:\allowbreak{}332733; schéma non conservé dans l’extrait documentaire; échelle non conservée dans l’extrait documentaire; niveau de confiance non conservé dans l’extrait documentaire.
\item \textbf{Différence relative:} 1290.825509\allowbreak{}927 ppm.
\item \textbf{Portée :} évaluation conditionnelle ; la signature n’a pas été choisie à partir de la valeur comparée.
\item \textbf{Convention résonante :} l’extrait externe ne fixe pas conjointement la convention de pôle, de Breit--Wigner ou sur couche de masse, la largeur et la covariance ; le résidu n’est interprété que comme information scalaire.
\end{itemize}
\Needspace{7\baselineskip}
\subsection{Pion neutre \texorpdfstring{\(\pi^{0}\)}{pi0}}
\begin{itemize}
\item \textbf{Identifiant du catalogue :} CMP-PI0.
\item \textbf{Type de descripteur :} signature complète σ\(\in\)Z\({}^{5}\).
\item \textbf{Descripteur :} (44,-4,-25,1,-2).
\item \textbf{Étape :} tour bêta E3 composée ; coordonnée initiale 135.350257\allowbreak{}251504\allowbreak{}283349 MeV/c\({}^{2}\).
\item \textbf{Évaluation HMT--MD reproductible :} 134.976724\allowbreak{}327872\allowbreak{}125739 MeV/c\({}^{2}\).
\item \textbf{Nature mathématique :} raffinement E3 conditionné par la signature composée.
\item \textbf{Référence ultérieure :} 134.976827\allowbreak{}767685 MeV/c\({}^{2}\).
\item \textbf{Identité et métrologie externes :} référence composée conservée par le supplément E3 ; la convention, la covariance et l’édition sont consignées dans la source métrologique de référence.
\item \textbf{Différence relative:} -0.766352\allowbreak{}378 ppm.
\item \textbf{Portée :} évaluation exacte une fois la signature composée fixée ; le générateur de liaison demeure un morphisme propre.
\end{itemize}
\Needspace{7\baselineskip}
\subsection{Pions chargés \(\pi^\pm\) : composition quark–antiquark, charge et masse}
\begin{itemize}
\item \textbf{Identifiant du catalogue :} CMP-PIPM.
\item \textbf{Type de descripteur :} signature complète σ\(\in\)Z\({}^{5}\).
\item \textbf{Descripteur :} (44,1,-2,2,-1).
\item \textbf{Étape :} tour bêta E3 composée ; coordonnée initiale 139.596265\allowbreak{}529971\allowbreak{}781015 MeV/c\({}^{2}\).
\item \textbf{Évaluation HMT--MD reproductible :} 139.570485\allowbreak{}171572\allowbreak{}191375 MeV/c\({}^{2}\).
\item \textbf{Nature mathématique :} raffinement E3 conditionné par la signature composée.
\item \textbf{Référence ultérieure :} 139.570390\allowbreak{}983681 MeV/c\({}^{2}\).
\item \textbf{Identité et métrologie externes :} référence composée conservée par le supplément E3 ; la convention, la covariance et l’édition sont consignées dans la source métrologique de référence.
\item \textbf{Différence relative:} 0.674841\allowbreak{}494 ppm.
\item \textbf{Portée :} évaluation exacte une fois la signature composée fixée ; le générateur de liaison demeure un morphisme propre.
\end{itemize}
\Needspace{7\baselineskip}
\subsection{Kaons chargés \(K^\pm\) : composition avec étrangeté, charge et masse}
\begin{itemize}
\item \textbf{Identifiant du catalogue :} CMP-KPM.
\item \textbf{Type de descripteur :} signature complète σ\(\in\)Z\({}^{5}\).
\item \textbf{Descripteur :} (54,-1,17,-2,2).
\item \textbf{Étape :} tour bêta E3 composée ; coordonnée initiale 492.762503\allowbreak{}210140\allowbreak{}547854 MeV/c\({}^{2}\).
\item \textbf{Évaluation HMT--MD reproductible :} 493.677085\allowbreak{}649530\allowbreak{}612476 MeV/c\({}^{2}\).
\item \textbf{Nature mathématique :} raffinement E3 conditionné par la signature composée.
\item \textbf{Référence ultérieure :} 493.676599\allowbreak{}458041 MeV/c\({}^{2}\).
\item \textbf{Identité et métrologie externes :} référence composée conservée par le supplément E3 ; la convention, la covariance et l’édition sont consignées dans la source métrologique de référence.
\item \textbf{Différence relative:} 0.984838\allowbreak{}03 ppm.
\item \textbf{Portée :} évaluation exacte une fois la signature composée fixée ; le générateur de liaison demeure un morphisme propre.
\end{itemize}
\Needspace{7\baselineskip}
\subsection{Kaon neutre \texorpdfstring{\(K^{0}\)}{K0}}
\begin{itemize}
\item \textbf{Identifiant du catalogue :} CMP-K0.
\item \textbf{Type de descripteur :} signature complète σ\(\in\)Z\({}^{5}\).
\item \textbf{Descripteur :} (54,0,32,3,-2).
\item \textbf{Étape :} tour bêta E3 composée ; coordonnée initiale 495.816066\allowbreak{}594426\allowbreak{}768595 MeV/c\({}^{2}\).
\item \textbf{Évaluation HMT--MD reproductible :} 497.611186\allowbreak{}497750\allowbreak{}457359 MeV/c\({}^{2}\).
\item \textbf{Nature mathématique :} raffinement E3 conditionné par la signature composée.
\item \textbf{Référence ultérieure :} 497.610767\allowbreak{}531258 MeV/c\({}^{2}\).
\item \textbf{Identité et métrologie externes :} référence composée conservée par le supplément E3 ; la convention, la covariance et l’édition sont consignées dans la source métrologique de référence.
\item \textbf{Différence relative:} 0.841956\allowbreak{}243 ppm.
\item \textbf{Portée :} évaluation exacte une fois la signature composée fixée ; le générateur de liaison demeure un morphisme propre.
\end{itemize}
\Needspace{7\baselineskip}
\subsection{Proton : composition baryonique, neutralisation de couleur et caractère de masse}
\begin{itemize}
\item \textbf{Identifiant du catalogue :} CMP-P.
\item \textbf{Type de descripteur :} signature complète σ\(\in\)Z\({}^{5}\).
\item \textbf{Descripteur :} (59,0,9,1,0).
\item \textbf{Étape :} tour bêta E3 composée ; coordonnée initiale 937.314779\allowbreak{}258699\allowbreak{}634563 MeV/c\({}^{2}\).
\item \textbf{Évaluation HMT--MD reproductible :} 938.271751\allowbreak{}546984\allowbreak{}898299 MeV/c\({}^{2}\).
\item \textbf{Nature mathématique :} raffinement E3 conditionné par la signature composée.
\item \textbf{Référence ultérieure :} 938.272089\allowbreak{}43 MeV/c\({}^{2}\).
\item \textbf{Identité et métrologie externes :} référence composée conservée par le supplément E3 ; la convention, la covariance et l’édition sont consignées dans la source métrologique de référence.
\item \textbf{Différence relative:} -0.360111\allowbreak{}975 ppm.
\item \textbf{Portée :} évaluation exacte une fois la signature composée fixée ; le générateur de liaison demeure un morphisme propre.
\end{itemize}
\Needspace{7\baselineskip}
\subsection{Neutron : composition baryonique, neutralisation de couleur et caractère de masse}
\begin{itemize}
\item \textbf{Identifiant du catalogue :} CMP-N.
\item \textbf{Type de descripteur :} signature complète σ\(\in\)Z\({}^{5}\).
\item \textbf{Descripteur :} (59,1,-34,0,1).
\item \textbf{Étape :} tour bêta E3 composée ; coordonnée initiale 943.123155\allowbreak{}648642\allowbreak{}014982 MeV/c\({}^{2}\).
\item \textbf{Évaluation HMT--MD reproductible :} 939.565810\allowbreak{}472507\allowbreak{}023313 MeV/c\({}^{2}\).
\item \textbf{Nature mathématique :} raffinement E3 conditionné par la signature composée.
\item \textbf{Référence ultérieure :} 939.565421\allowbreak{}94 MeV/c\({}^{2}\).
\item \textbf{Identité et métrologie externes :} référence composée conservée par le supplément E3 ; la convention, la covariance et l’édition sont consignées dans la source métrologique de référence.
\item \textbf{Différence relative:} 0.413523\allowbreak{}633 ppm.
\item \textbf{Portée :} évaluation exacte une fois la signature composée fixée ; le générateur de liaison demeure un morphisme propre.
\end{itemize}
% Las tablas históricas abreviadas permanecen conservadas a continuación en
% este propietario, pero quedan fuera del tronco científico vigente: no
% definen el dominio 81/56/13/324, la firma, el carácter ni las torres.
\endinput
\clearpage
\begin{landscape}
\section{Maillage historique étendu conservé}
\begingroup
\fontsize{6.8}{7.7}\selectfont
\setlength{\tabcolsep}{1.15pt}
\renewcommand{\arraystretch}{1.06}
\begin{longtable}{L{2.93cm}L{2.93cm}L{2.93cm}L{2.93cm}L{2.93cm}L{2.93cm}L{2.93cm}L{2.93cm}}
\toprule
\rowcolor{HMTNavy}\HMTtablehead{État} & \HMTtablehead{n} & \HMTtablehead{k} & \HMTtablehead{t} & \HMTtablehead{Évaluation HMT historique (MeV/\allowbreak{}c\({}^{2}\))} & \HMTtablehead{Référence ultérieure (MeV/\allowbreak{}c\({}^{2}\))} & \HMTtablehead{Différence consignée (ppm)} & \HMTtablehead{Force mathématique} \\
\midrule
\endfirsthead
\multicolumn{8}{l}{\sffamily\scriptsize\color{HMTGray}Suite du tableau}\\[-1pt]
\toprule
\rowcolor{HMTNavy}\HMTtablehead{État} & \HMTtablehead{n} & \HMTtablehead{k} & \HMTtablehead{t} & \HMTtablehead{Évaluation HMT historique (MeV/\allowbreak{}c\({}^{2}\))} & \HMTtablehead{Référence ultérieure (MeV/\allowbreak{}c\({}^{2}\))} & \HMTtablehead{Différence consignée (ppm)} & \HMTtablehead{Force mathématique} \\
\midrule
\endhead
\midrule
\endfoot
\bottomrule
\endlastfoot
mu & 64 & 146 & 1 & 105.656529 & 105.658374 & -17.461938\allowbreak{}2274 & preuve numérique historique conditionnelle ; ne constitue pas une règle physique de sélection antérieure à la comparaison \\
\rowcolor{HMTLight}tau & 64 & 349 & -1 & 1776.835 & 1776.86 & -14.069763\allowbreak{}5154 & preuve numérique historique conditionnelle ; ne constitue pas une règle physique de sélection antérieure à la comparaison \\
π±\allowbreak{} & 37 & 125 & 0 & 139.570582 & 139.57039 & 1.375649\allowbreak{}94982 & preuve numérique historique conditionnelle ; ne constitue pas une règle physique de sélection antérieure à la comparaison \\
\rowcolor{HMTLight}K±\allowbreak{} & 52 & 177 & 1 & 493.679 & 493.677 & 4.051231\allowbreak{}87833 & preuve numérique historique conditionnelle ; ne constitue pas une règle physique de sélection antérieure à la comparaison \\
K\({}^{0}\) & 53 & 178 & 1 & 497.616 & 497.611 & 10.048009\allowbreak{}3889 & preuve numérique historique conditionnelle ; ne constitue pas une règle physique de sélection antérieure à la comparaison \\
\rowcolor{HMTLight}p & 55 & 195 & 1 & 938.27244 & 938.272081 & 0.382618\allowbreak{}226919 & preuve numérique historique conditionnelle ; ne constitue pas une règle physique de sélection antérieure à la comparaison \\
n & 55 & 197 & 0 & 939.538 & 939.565413 & -29.176254\allowbreak{}9161 & preuve numérique historique conditionnelle ; ne constitue pas une règle physique de sélection antérieure à la comparaison \\
\rowcolor{HMTLight}eta & 46 & 158 & 0 & 547.836 & 547.862 & -47.457206\allowbreak{}3768 & preuve numérique historique conditionnelle ; ne constitue pas une règle physique de sélection antérieure à la comparaison \\
ρ\({}^{0}\) & 50 & 176 & 0 & 775.267 & 775.26 & 9.029228\allowbreak{}90385 & preuve numérique historique conditionnelle ; ne constitue pas une règle physique de sélection antérieure à la comparaison \\
\rowcolor{HMTLight}φ(1020) & 54 & 190 & 0 & 1019.475 & 1019.461 & 13.732747\allowbreak{}0104 & preuve numérique historique conditionnelle ; ne constitue pas une règle physique de sélection antérieure à la comparaison \\
J/\allowbreak{}ψ & 64 & 237 & 0 & 3096.979 & 3096.916 & 20.342818\allowbreak{}4684 & preuve numérique historique conditionnelle ; ne constitue pas une règle physique de sélection antérieure à la comparaison \\
\rowcolor{HMTLight}Υ(1S) & 89 & 334 & 0 & 9460.34 & 9460.3 & 4.228195\allowbreak{}72318 & preuve numérique historique conditionnelle ; ne constitue pas une règle physique de sélection antérieure à la comparaison \\
\end{longtable}
\endgroup
\end{landscape}
Ces douze valeurs ne sont pas fusionnées avec les vingt sorties du tableau de référence. Le maillage conserve cinq étiquettes supplémentaires ---η, ρ\({}^{0}\), φ(1020), J/ψ et Υ(1S)---, mais ne construit pas encore l’application état composé→parcours et enregistrement généalogique→(n,k,t). Le quantum interne CT108, 0.058177\allowbreak{}641733\allowbreak{}144319\allowbreak{}230789\allowbreak{}692282\allowbreak{}953757 dans la normalisation hbar\_int=\allowbreak{}c\_int=\allowbreak{}t0=\allowbreak{}1, est une section de la droite de masse et non une espèce.

Le tableau ne délimite pas le domaine. Le domaine interne demeure l’atlas de 81 cellules, 56 familles, 13 multisections et 324 parcours orientés ; les réalisations physiques typées et l’inventaire de 471 masses et de 384 largeurs sont développés ensuite.

\Needspace{7\baselineskip}

\endgroup

\section{Correspondance entre l’atlas \(81/56/13/324\) et les inventaires exhaustifs de \(471\) masses et \(384\) largeurs}
\label{chap:portal-atlas-inventarios-masas}
L’ \hyperref[chap:inventario-universal-471-masas-384-anchuras]{inventaire de 471 masses et 384 largeurs} est un catalogue métrologique externe à destination typée. Ce n’est ni un recensement de prédictions ni le cardinal du domaine matériel.
